
\documentclass[11pt]{article}
\usepackage[T1]{fontenc}
\usepackage[utf8]{inputenc}
\usepackage[english]{babel}
\usepackage{lmodern}
\usepackage{microtype}
\usepackage{geometry}
\usepackage{setspace}
\usepackage{amsmath,amssymb,amsthm,mathtools,mathrsfs}
\usepackage{enumitem}
\usepackage{booktabs,longtable,array}
\usepackage[all]{xy}
\usepackage{tikz-cd}
\usepackage{hyperref}
\hypersetup{colorlinks=true,linkcolor=blue,citecolor=blue,urlcolor=blue,hypertexnames=false}
\geometry{margin=1.05in}
\setstretch{1.12}
\setlength{\parindent}{0pt}
\setlength{\parskip}{0.42em}
\sloppy

\theoremstyle{plain}
\newtheorem{theorem}{Theorem}[section]
\newtheorem{proposition}[theorem]{Proposition}
\newtheorem{lemma}[theorem]{Lemma}
\newtheorem{corollary}[theorem]{Corollary}
\newtheorem{theoreme}[theorem]{Theorem}
\newtheorem{propositionfr}[theorem]{Proposition}
\newtheorem{lemme}[theorem]{Lemma}
\newtheorem{corollaire}[theorem]{Corollary}

\newtheorem*{theorem*}{Theorem}
\newtheorem*{proposition*}{Proposition}
\newtheorem*{lemma*}{Lemma}
\newtheorem*{corollary*}{Corollary}
\newtheorem*{definition*}{Definition}
\newtheorem*{remark*}{Remark}
\newtheorem*{exercise*}{Exercise}
\newtheorem*{problem*}{Problem}

\theoremstyle{definition}
\newtheorem{definition}[theorem]{Definition}
\newtheorem{example}[theorem]{Example}
\newtheorem{exercise}[theorem]{Exercise}
\newtheorem{problem}[theorem]{Problem}
\newtheorem{notation}[theorem]{Notation}
\theoremstyle{remark}
\newtheorem{remark}[theorem]{Remark}
\newtheorem{remarks}[theorem]{Remarks}
\newtheorem{remarque}[theorem]{Remark}
\newtheorem{remarques}[theorem]{Remarks}
\newtheorem{scholium}[theorem]{Scholium}
\newenvironment{statement}[1]{\par\medskip\noindent\textbf{#1.}\ }{\par\medskip}
\newenvironment{proofsketch}{\par\noindent\textit{Proof.}\ }{\par\medskip}
\newenvironment{prooftext}{\par\noindent\textit{Proof.}\ }{\par\medskip}
\newenvironment{remarktext}{\par\medskip\noindent\textit{Remark.}\ }{\par\medskip}


% Blackboard symbols
\newcommand{\bbA}{\mathbb A}
\newcommand{\bbC}{\mathbb C}
\newcommand{\bbF}{\mathbb F}
\newcommand{\bbG}{\mathbb G}
\newcommand{\bbN}{\mathbb N}
\newcommand{\bbP}{\mathbb P}
\newcommand{\bbQ}{\mathbb Q}
\newcommand{\bbR}{\mathbb R}
\newcommand{\bbZ}{\mathbb Z}
\newcommand{\bA}{\mathbb A}
\newcommand{\bC}{\mathbb C}
\newcommand{\bF}{\mathbb F}
\newcommand{\bG}{\mathbb G}
\newcommand{\bN}{\mathbb N}
\newcommand{\bP}{\mathbb P}
\newcommand{\bQ}{\mathbb Q}
\newcommand{\bR}{\mathbb R}
\newcommand{\bZ}{\mathbb Z}
\renewcommand{\AA}{\mathbb A}
\newcommand{\CC}{\mathbb C}
\newcommand{\FF}{\mathbb F}
\newcommand{\GG}{\mathbb G}
\newcommand{\NN}{\mathbb N}
\newcommand{\PP}{\mathbb P}
\newcommand{\QQ}{\mathbb Q}
\newcommand{\RR}{\mathbb R}
\newcommand{\ZZ}{\mathbb Z}
\newcommand{\Z}{\mathbb Z}\newcommand{\Q}{\mathbb Q}\newcommand{\N}{\mathbb N}\newcommand{\A}{\mathbb A}\newcommand{\Pp}{\mathbb P}

% Calligraphic symbols
\newcommand{\cA}{\mathcal A}\newcommand{\cB}{\mathcal B}\newcommand{\cC}{\mathcal C}
\newcommand{\cD}{\mathcal D}\newcommand{\cE}{\mathcal E}\newcommand{\cF}{\mathcal F}
\newcommand{\cG}{\mathcal G}\newcommand{\cH}{\mathcal H}\newcommand{\cI}{\mathcal I}
\newcommand{\cJ}{\mathcal J}\newcommand{\cK}{\mathcal K}\newcommand{\cL}{\mathcal L}
\newcommand{\cM}{\mathcal M}\newcommand{\cN}{\mathcal N}\newcommand{\cO}{\mathcal O}
\newcommand{\cP}{\mathcal P}\newcommand{\cQ}{\mathcal Q}\newcommand{\cR}{\mathcal R}
\newcommand{\cS}{\mathcal S}\newcommand{\cT}{\mathcal T}\newcommand{\cU}{\mathcal U}
\newcommand{\cV}{\mathcal V}\newcommand{\cW}{\mathcal W}\newcommand{\cX}{\mathcal X}
\newcommand{\cY}{\mathcal Y}\newcommand{\cZ}{\mathcal Z}
\newcommand{\U}{\mathcal U}\newcommand{\V}{\mathcal V}\newcommand{\C}{\mathcal C}
\newcommand{\D}{\mathcal D}\newcommand{\E}{\mathcal E}\newcommand{\F}{\mathcal F}
\newcommand{\G}{\mathcal G}\newcommand{\Hh}{\mathcal H}\newcommand{\M}{\mathcal M}
\newcommand{\sA}{\mathscr A}\newcommand{\sB}{\mathscr B}\newcommand{\sC}{\mathscr C}

% Common operators and categories
\DeclareMathOperator{\Spec}{Spec}
\DeclareMathOperator{\Spf}{Spf}
\DeclareMathOperator{\Proj}{Proj}
\DeclareMathOperator{\Hom}{Hom}
\DeclareMathOperator{\End}{End}
\DeclareMathOperator{\Aut}{Aut}
\DeclareMathOperator{\Isom}{Isom}
\DeclareMathOperator{\Ext}{Ext}
\DeclareMathOperator{\Exalcomm}{Exalcomm}
\DeclareMathOperator{\Tor}{Tor}
\DeclareMathOperator{\Biext}{Biext}
\DeclareMathOperator{\Pic}{Pic}
\DeclareMathOperator{\Div}{Div}
\DeclareMathOperator{\Der}{Der}
\DeclareMathOperator{\Ker}{Ker}
\DeclareMathOperator{\Coker}{Coker}
\DeclareMathOperator{\Imop}{Im}
\DeclareMathOperator{\Imn}{Im}
\DeclareMathOperator{\Ob}{Ob}
\DeclareMathOperator{\ob}{Ob}
\DeclareMathOperator{\Gal}{Gal}
\DeclareMathOperator{\GL}{GL}
\DeclareMathOperator{\PGL}{PGL}
\DeclareMathOperator{\SL}{SL}
\DeclareMathOperator{\rank}{rank}
\DeclareMathOperator{\Reg}{Reg}
\DeclareMathOperator{\Sing}{Sing}
\DeclareMathOperator{\Supp}{Supp}
\DeclareMathOperator{\Tr}{Tr}
\DeclareMathOperator{\Fr}{Fr}
\DeclareMathOperator{\Frob}{Frob}
\DeclareMathOperator{\length}{length}
\DeclareMathOperator{\card}{card}
\DeclareMathOperator{\gr}{gr}
\DeclareMathOperator{\Tot}{Tot}
\DeclareMathOperator{\Cone}{Cone}
\DeclareMathOperator{\Gr}{Gr}
\DeclareMathOperator{\id}{id}
\DeclareMathOperator{\pr}{pr}
\DeclareMathOperator{\supp}{supp}
\DeclareMathOperator{\codim}{codim}
\DeclareMathOperator{\coh}{coh}
\DeclareMathOperator{\cd}{cd}
\DeclareMathOperator{\cl}{cl}
\DeclareMathOperator{\RHom}{RHom}
\DeclareMathOperator{\Rlim}{Rlim}
\newcommand{\Set}{\mathbf{Set}}\newcommand{\Ens}{\mathbf{Set}}
\newcommand{\Cat}{\mathbf{Cat}}\newcommand{\Ab}{\mathbf{Ab}}
\newcommand{\Top}{\operatorname{Top}}\newcommand{\TOP}{\operatorname{TOP}}
\newcommand{\BIEXT}{\mathbf{BIEXT}}\newcommand{\GRAB}{\mathbf{GRAB}}
\newcommand{\Mod}{\mathbf{Mod}}

% Miscellaneous
\newcommand{\Gm}{\mathbb G_m}
\newcommand{\Ga}{\mathbb G_a}
\newcommand{\Ql}{\mathbb Q_\ell}
\newcommand{\Zl}{\mathbb Z_\ell}
\newcommand{\Tl}{T_\ell}
\newcommand{\Tp}{T_p}
\newcommand{\Vl}{V_\ell}
\newcommand{\bL}{\mathbf L}
\newcommand{\longto}{\longrightarrow}
\newcommand{\into}{\hookrightarrow}
\newcommand{\onto}{\twoheadrightarrow}
\newcommand{\ov}[1]{\overline{#1}}
\newcommand{\wt}[1]{\widetilde{#1}}
\newcommand{\ol}[1]{\overline{#1}}
\newcommand{\ul}[1]{\underline{#1}}
\newcommand{\uA}{\underline{A}}\newcommand{\uH}{\underline{H}}\newcommand{\uExt}{\underline{\operatorname{Ext}}}\newcommand{\uL}{\underline{L}}\newcommand{\uM}{\underline{M}}\newcommand{\uZ}{\underline{Z}}\newcommand{\uB}{\underline{B}}\newcommand{\uC}{\underline{C}}\newcommand{\uD}{\underline{D}}\newcommand{\uE}{\underline{E}}\newcommand{\uG}{\underline{G}}\newcommand{\uP}{\underline{P}}\newcommand{\uQ}{\underline{Q}}
\newcommand{\et}{\mathrm{\acute{e}t}}
\newcommand{\ab}{\mathrm{ab}}
\newcommand{\frakm}{\mathfrak m}
\newcommand{\fm}{\mathfrak m}
\newcommand{\ouv}{\operatorname{Open}}
\newcommand{\Open}{\operatorname{Open}}
\newcommand{\HH}{\mathrm H}
\newcommand{\AR}{\mathrm{AR}}
\newcommand{\R}{\mathrm R}
\newcommand{\longueur}{\operatorname{length}}
\newcommand{\rang}{\operatorname{rank}}
\newcommand{\PGCD}{\operatorname{gcd}}
\newcommand{\qedtext}{\hfill$\square$}
\newcommand{\Autext}{\operatorname{Aut\,ext}}
\newcommand{\prof}{\operatorname{prof}}


% Additional compatibility macros for cumulative compilation
\providecommand{\adn}{\operatorname{adn}}
\providecommand{\adf}{\operatorname{adf}}
\providecommand{\Catfl}{\operatorname{Catfl}}
\providecommand{\Catfib}{\operatorname{Catfib}}
\providecommand{\Fib}{\operatorname{Fib}}
\providecommand{\Fl}{\operatorname{Fl}}
\providecommand{\Sh}{\operatorname{Sh}}
\providecommand{\Psh}{\operatorname{PSh}}
\providecommand{\Cov}{\operatorname{Cov}}
\providecommand{\Car}{\operatorname{Car}}
\providecommand{\coker}{\operatorname{coker}}
\providecommand{\kernel}{\operatorname{ker}}
\providecommand{\image}{\operatorname{im}}
\providecommand{\coim}{\operatorname{coim}}
\providecommand{\Id}{\operatorname{Id}}
\providecommand{\tr}{\operatorname{tr}}
\providecommand{\Nm}{\operatorname{Nm}}
\providecommand{\ch}{\operatorname{ch}}
\providecommand{\td}{\operatorname{td}}
\providecommand{\chern}{\operatorname{ch}}
\providecommand{\codim}{\operatorname{codim}}
\providecommand{\Torind}{\operatorname{Tor-ind}}
\providecommand{\ind}{\operatorname{ind}}
\providecommand{\pro}{\operatorname{pro}}
\providecommand{\Ind}{\operatorname{Ind}}
\providecommand{\Pro}{\operatorname{Pro}}
\providecommand{\liminj}{\mathop{\varinjlim}}
\providecommand{\limproj}{\mathop{\varprojlim}}
\providecommand{\minus}{-}
\providecommand{\plus}{+}
\providecommand{\xto}[1]{\xrightarrow{#1}}
\providecommand{\xot}[1]{\xleftarrow{#1}}

\providecommand{\fc}{\operatorname{fc}}
\providecommand{\bd}{\operatorname{bd}}
\providecommand{\tf}{\operatorname{tf}}
\providecommand{\etale}{\ifmmode\acute{e}tale\else \'etale\fi}
\providecommand{\Sw}{\operatorname{Sw}}
\providecommand{\sw}{\operatorname{sw}}
\providecommand{\Art}{\operatorname{Art}}
\providecommand{\Cond}{\operatorname{Cond}}
\providecommand{\Swan}{\operatorname{Swan}}
\providecommand{\wh}{\widehat}
\providecommand{\Cl}{\operatorname{Cl}}
\providecommand{\K}{\mathrm K}
\providecommand{\Todd}{\operatorname{Todd}}
\providecommand{\GRR}{\operatorname{GRR}}
\providecommand{\Op}{\operatorname{Op}}
\providecommand{\Fil}{\operatorname{Fil}}
\providecommand{\grK}{\operatorname{grK}}
\providecommand{\eps}{\varepsilon}
\providecommand{\Parf}{\operatorname{Parf}}
\providecommand{\Perf}{\operatorname{Perf}}
\providecommand{\Dparf}{\operatorname{D}_{\mathrm{parf}}}
\providecommand{\Dqc}{\operatorname{D}_{\mathrm{qc}}}
\providecommand{\Dcoh}{\operatorname{D}_{\mathrm{coh}}}
\providecommand{\detn}{\operatorname{det}}
\providecommand{\Inv}{\operatorname{Inv}}
\providecommand{\Picard}{\operatorname{Picard}}
\providecommand{\rad}{\operatorname{rad}}
\providecommand{\Rep}{\operatorname{Rep}}
\providecommand{\charac}{\operatorname{char}}
\providecommand{\uF}{\underline{F}}
\providecommand{\uI}{\underline{I}}
\providecommand{\uJ}{\underline{J}}
\providecommand{\uK}{\underline{K}}
\providecommand{\uN}{\underline{N}}
\providecommand{\uO}{\underline{O}}
\providecommand{\uR}{\underline{R}}
\providecommand{\uS}{\underline{S}}
\providecommand{\uT}{\underline{T}}
\providecommand{\uU}{\underline{U}}
\providecommand{\uV}{\underline{V}}
\providecommand{\uW}{\underline{W}}
\providecommand{\uX}{\underline{X}}
\providecommand{\uY}{\underline{Y}}
\providecommand{\iso}{\xrightarrow{\sim}}
\providecommand{\isom}{\xrightarrow{\sim}}
\providecommand{\Rad}{\operatorname{Rad}}
\providecommand{\Jac}{\operatorname{Jac}}
\providecommand{\Add}{\operatorname{Add}}
\providecommand{\Endom}{\operatorname{End}}
\providecommand{\Homol}{\operatorname{Hom}}
\providecommand{\Cart}{\operatorname{Cart}}
\providecommand{\Cartes}{\operatorname{Cart}}
\providecommand{\rg}{\operatorname{rank}}
\providecommand{\rk}{\operatorname{rank}}
\providecommand{\bbB}{\mathbb{B}}
\providecommand{\bbD}{\mathbb{D}}
\providecommand{\bbE}{\mathbb{E}}
\providecommand{\bbH}{\mathbb{H}}
\providecommand{\bbI}{\mathbb{I}}
\providecommand{\bbJ}{\mathbb{J}}
\providecommand{\bbK}{\mathbb{K}}
\providecommand{\bbL}{\mathbb{L}}
\providecommand{\bbM}{\mathbb{M}}
\providecommand{\bbO}{\mathbb{O}}
\providecommand{\bbS}{\mathbb{S}}
\providecommand{\bbT}{\mathbb{T}}
\providecommand{\bbU}{\mathbb{U}}
\providecommand{\bbV}{\mathbb{V}}
\providecommand{\bbW}{\mathbb{W}}
\providecommand{\bbX}{\mathbb{X}}
\providecommand{\bbY}{\mathbb{Y}}
\providecommand{\bB}{\mathbb{B}}
\providecommand{\bD}{\mathbb{D}}
\providecommand{\bE}{\mathbb{E}}
\providecommand{\bH}{\mathbb{H}}
\providecommand{\bI}{\mathbb{I}}
\providecommand{\bJ}{\mathbb{J}}
\providecommand{\bK}{\mathbb{K}}
\providecommand{\bM}{\mathbb{M}}
\providecommand{\bO}{\mathbb{O}}
\providecommand{\bS}{\mathbb{S}}
\providecommand{\bT}{\mathbb{T}}
\providecommand{\bU}{\mathbb{U}}
\providecommand{\bV}{\mathbb{V}}
\providecommand{\bW}{\mathbb{W}}
\providecommand{\bX}{\mathbb{X}}
\providecommand{\bY}{\mathbb{Y}}
\providecommand{\Qcoh}{\operatorname{Qcoh}}
\providecommand{\Coh}{\operatorname{Coh}}
\providecommand{\QCoh}{\operatorname{QCoh}}
\providecommand{\Loclib}{\operatorname{Loclib}}
\providecommand{\LocFree}{\operatorname{LocFree}}
\providecommand{\Lib}{\operatorname{Lib}}
\providecommand{\Projmod}{\operatorname{Projmod}}
\providecommand{\Ltens}{\overset{\mathbf L}{\otimes}}
\providecommand{\Lotimes}{\overset{\mathbf L}{\otimes}}
\providecommand{\Rtens}{\overset{\mathbf R}{\otimes}}
\providecommand{\Pt}{\operatorname{Pt}}
\providecommand{\Points}{\operatorname{Points}}
\providecommand{\Go}{\operatorname{G}}
\providecommand{\Ko}{\operatorname{K}}
\providecommand{\Kzero}{K_0}
\providecommand{\Goi}{\operatorname{G}_i}
\providecommand{\Gzero}{\operatorname{G}_0}
\providecommand{\Got}{\operatorname{G}}
\providecommand{\Gou}{\operatorname{G}}
\providecommand{\us}{_{\!S}}
\providecommand{\Obj}{\operatorname{Obj}}
\providecommand{\Mor}{\operatorname{Mor}}
\providecommand{\Ch}{\operatorname{Ch}}
\providecommand{\HHH}{\operatorname{H}}
\providecommand{\Log}{\operatorname{Log}}
\providecommand{\Exp}{\operatorname{Exp}}
\providecommand{\Sym}{\operatorname{Sym}}
\providecommand{\ProjS}{\operatorname{Proj}}
\providecommand{\Comp}{\operatorname{Comp}}
\providecommand{\Totcomp}{\operatorname{Comp}}
\providecommand{\Kdot}{K_\bullet}
\providecommand{\Gdot}{G_\bullet}
\providecommand{\Kprime}{K\prime}
\providecommand{\LL}{\mathbf L}
\providecommand{\VV}{\mathbb V}
\providecommand{\AA}{\mathbb A}
\providecommand{\Drap}{\operatorname{Flag}}
\providecommand{\Flag}{\operatorname{Flag}}
\providecommand{\Gl}{\operatorname{GL}}
\providecommand{\Pgl}{\operatorname{PGL}}
\providecommand{\Sl}{\operatorname{SL}}
\providecommand{\Filt}{\operatorname{Fil}}
\providecommand{\GrF}{\operatorname{Gr}}
\providecommand{\trdeg}{\operatorname{trdeg}}
\providecommand{\dimtr}{\operatorname{trdeg}}
\providecommand{\detm}{\operatorname{det}}
\providecommand{\rankm}{\operatorname{rank}}
\providecommand{\bmu}{\boldsymbol{\mu}}
\providecommand{\muell}{\boldsymbol{\mu}_\ell}
\providecommand{\car}{\operatorname{char}}
\providecommand{\char}{\operatorname{char}}
\providecommand{\red}{\mathrm{red}}
\providecommand{\nil}{\mathrm{nil}}
\providecommand{\num}{\mathrm{num}}
\providecommand{\rat}{\mathrm{rat}}
\providecommand{\alg}{\mathrm{alg}}
\providecommand{\NS}{\operatorname{NS}}
\providecommand{\NeronSeveri}{\operatorname{NS}}
\providecommand{\Tors}{\operatorname{Tors}}
\providecommand{\Free}{\operatorname{Free}}
\providecommand{\Td}{\operatorname{Td}}
\providecommand{\topf}{\mathrm{top}}
\providecommand{\geom}{\mathrm{geom}}
\providecommand{\ie}{i.e.}
\providecommand{\eg}{e.g.}
\providecommand{\cf}{cf.}
\providecommand{\loccit}{\emph{loc. cit.}}
\providecommand{\opcit}{\emph{op. cit.}}
\providecommand{\Esp}{\operatorname{Esp}}
\providecommand{\Topos}{\operatorname{Topos}}
\providecommand{\Sch}{\operatorname{Sch}}
\providecommand{\Hhom}{\operatorname{Hom}}
\providecommand{\HHom}{\operatorname{Hom}}
\providecommand{\Hhomtop}{\operatorname{Hom}_{\mathrm{top}}}
\providecommand{\Homtop}{\operatorname{Hom}_{\mathrm{top}}}
\providecommand{\sob}{\mathrm{sob}}
\providecommand{\contracted}{\mathbin{\times}}
\providecommand{\acts}{\curvearrowright}
\providecommand{\dfn}{\mathrm{def}}
\providecommand{\Ppoints}{\operatorname{Pts}}
\providecommand{\Ner}{\operatorname{Ner}}
\providecommand{\Nerve}{\operatorname{Ner}}
\providecommand{\Mmorsite}{\operatorname{MorSite}}
\providecommand{\Morsite}{\operatorname{MorSite}}
\providecommand{\morsite}{\operatorname{MorSite}}
\providecommand{\Res}{\operatorname{Res}}
\providecommand{\Rest}{\operatorname{Res}}
\providecommand{\Prol}{\operatorname{Prol}}
\providecommand{\Compct}{\operatorname{Comp}}
\providecommand{\Ttors}{\operatorname{Tors}}
\providecommand{\tors}{\operatorname{Tors}}
\providecommand{\BTors}{\operatorname{BTors}}
\providecommand{\Ssex}{\operatorname{Sec}}
\providecommand{\Ssect}{\operatorname{Sec}}
\providecommand{\sections}{\operatorname{Sec}}
\providecommand{\Hhomind}{\operatorname{Hom}_{\mathrm{ind}}}
\providecommand{\Homind}{\operatorname{Hom}_{\mathrm{ind}}}
\providecommand{\Ppoint}{\operatorname{Point}}
\providecommand{\Point}{\operatorname{Point}}
\providecommand{\Ffib}{\operatorname{Fib}}
\providecommand{\Fibfun}{\operatorname{Fib}}
\providecommand{\Vvois}{\operatorname{Nbhd}}
\providecommand{\Vois}{\operatorname{Nbhd}}
\providecommand{\quav}{\longrightarrow}
\providecommand{\im}{\operatorname{im}}
\providecommand{\coimage}{\operatorname{coim}}
\providecommand{\kar}{\operatorname{Kar}}
\providecommand{\Kar}{\operatorname{Kar}}
\providecommand{\Ppointess}{\operatorname{Point}^{\mathrm{ess}}}
\providecommand{\Pointess}{\operatorname{Point}^{\mathrm{ess}}}
\providecommand{\Hhomtopess}{\operatorname{Hom}_{\mathrm{top}}^{\mathrm{ess}}}
\providecommand{\Homtopess}{\operatorname{Hom}_{\mathrm{top}}^{\mathrm{ess}}}
\providecommand{\Ouv}{\operatorname{Open}}
\providecommand{\Fer}{\operatorname{Closed}}
\providecommand{\Oo}{\mathcal O}
\providecommand{\Bb}{\mathcal B}
\providecommand{\cat}{\operatorname{cat}}
\providecommand{\SGA}{SGA}
\providecommand{\EGA}{EGA}
\providecommand{\cons}{\mathrm{cons}}
\providecommand{\construct}{\mathrm{cons}}
\providecommand{\emptytopos}{\varnothing}
\providecommand{\pttopos}{\ast}
\providecommand{\Uu}{\mathcal U}
\providecommand{\Zz}{\mathcal Z}
\providecommand{\Ff}{\mathcal F}
\providecommand{\Gg}{\mathcal G}
\providecommand{\Sup}{\bigcup}
\providecommand{\Inf}{\bigcap}
\providecommand{\Quot}{\operatorname{Quot}}
\providecommand{\Qu}{\operatorname{Quot}}
\providecommand{\loc}{\operatorname{loc}}
\providecommand{\Loc}{\operatorname{Loc}}
\providecommand{\glob}{\operatorname{glob}}
\providecommand{\Glob}{\operatorname{Glob}}
\providecommand{\scrHom}{\mathscr{H}\!om}
\providecommand{\sHom}{\mathscr{H}\!om}
\providecommand{\ad}{\operatorname{ad}}
\providecommand{\can}{\operatorname{can}}
\providecommand{\resp}{resp.}
\providecommand{\Hhomtoplocan}{\operatorname{Hom}_{\mathrm{top,loc.an}}}
\providecommand{\Homtoplocan}{\operatorname{Hom}_{\mathrm{top,loc.an}}}
\providecommand{\Hhomtopan}{\operatorname{Hom}_{\mathrm{top,an}}}
\providecommand{\Homtopan}{\operatorname{Hom}_{\mathrm{top,an}}}
\providecommand{\Pptgeom}{\operatorname{Pt}_{\mathrm{geom}}}
\providecommand{\Ptgeom}{\operatorname{Pt}_{\mathrm{geom}}}
\providecommand{\Cech}{\v{C}ech}
\providecommand{\checkC}{\check C}
\providecommand{\Mcat}{\mathcal M}
\providecommand{\Modcat}{\mathcal M}
\providecommand{\B}{\mathcal{B}}
\providecommand{\H}{\mathcal{H}}
\providecommand{\I}{\mathcal{I}}
\providecommand{\J}{\mathcal{J}}
\providecommand{\L}{\mathcal{L}}
\providecommand{\O}{\mathcal{O}}
\providecommand{\P}{\mathcal{P}}
\providecommand{\S}{\mathcal{S}}
\providecommand{\T}{\mathcal{T}}
\providecommand{\W}{\mathcal{W}}
\providecommand{\X}{\mathcal{X}}
\providecommand{\Y}{\mathcal{Y}}
\providecommand{\Acat}{\mathcal A}
\providecommand{\Bcat}{\mathcal B}
\providecommand{\Ccat}{\mathcal C}
\providecommand{\Dcat}{\mathcal D}
\providecommand{\Sex}{\operatorname{Sec}}
\providecommand{\Sec}{\operatorname{Sec}}
\providecommand{\Sect}{\operatorname{Sec}}
\providecommand{\op}{\mathrm{op}}
\providecommand{\Eext}{\operatorname{Ext}}
\providecommand{\Extsheaf}{\mathscr{E}\!xt}
\providecommand{\sk}{\operatorname{sk}}
\providecommand{\cosk}{\operatorname{cosk}}
\begin{document}
\begin{titlepage}
\centering
\vspace*{2.3in}
{\LARGE\textbf{SGA 5}}\\[0.5em]
{\Large $\ell$-adic Cohomology and \(L\)-Functions}\\[1.5em]
{\large English translation}
\vfill
\end{titlepage}
\tableofcontents


\clearpage
\section*{SGA 5, Volume Introduction}
\addcontentsline{toc}{section}{SGA 5, Volume Introduction}

\section*{Introduction}

This volume brings together the exposés from Grothendieck's seminar at the IHES in 1965--66, initially distributed in the form of mimeographed notes. Compared with the original version, the only important changes concern Exposé II, which is not reproduced, and Exposé III, which has been entirely rewritten and enlarged by an appendix numbered III B. Apart from a few detailed modifications and added footnotes, the other exposés have been left as they stood.

The heart of the seminar is the Lefschetz formula in étale cohomology (III, III B, XII), and its application to the cohomological interpretation of $L$-functions (XIV). All the results announced by Grothendieck in his Bourbaki seminar talk [2] are proved here completely. The trace formulae established in (III, III B, and XII), by different methods, are more general than is necessary for proving merely the rationality of the $L$-functions. A complete and much shorter proof of the latter appears in Deligne's exposé (SGA $4^{1/2}$, Report), where Grothendieck's method from Exposés XII and XIV is followed, but stripped of all superfluous generality. We nevertheless hope that the formulae of III and III B may be useful in other situations. As for the rest of the seminar, it consists of two exposés on the theory of passage to the limit that gives rise to $\ell$-adic cohomology (V, VI), and of various complements to the duality formalism (I, VII) and to the Lefschetz formula (VII, X).

Here, more precisely, is the content of this volume.

Exposé I is independent of the rest of the seminar. Its main result is that, on a regular scheme satisfying certain additional local hypotheses, and assuming the availability of resolution of singularities and the purity theorem, the constant sheaf with value $\Z/n\Z$, for $n$ prime to the residue characteristics, is dualizing (I.3.4.1). This theorem applies in particular to excellent regular schemes of characteristic zero, thanks to Hironaka and to Artin's results (SGA 4, XIX). The case of a regular scheme of dimension $1$ is treated separately, by another, very elementary method; see also SGA $4^{1/2}$, Duality. Let us also point out that, by a different argument using neither resolution nor purity, Deligne proves in SGA $4^{1/2}$, Finiteness Theorem, that if $f:X\to S$ is a scheme of finite type over a regular scheme of dimension $0$ or $1$, then the complex $f^!(\Z/n\Z)_S$, for $n$ prime to the residue characteristics, is dualizing. It is not known, however, whether this result, which is more useful than the biduality theorem of Exposé I, and which overlaps it without containing it, extends to the case where the base scheme $S$ is excellent, regular, and of dimension $>1$.

As was said above, Exposé II, entitled ``Künneth formulae for cohomology with arbitrary supports'', does not appear in this volume. Written by L. Illusie from Grothendieck's handwritten notes, it was devoted to theorems on cohomological properness and generic local acyclicity; but these were proved only under resolution hypotheses. A proof of the same results without these hypotheses, obtained since then by Deligne (SGA $4^{1/2}$, Finiteness Theorem), made publication of that exposé unnecessary. Certain complements, concerning for example the case of normal-crossing divisors, were incorporated into an appendix to \emph{loc. cit.} We add that theorems analogous to those of \emph{loc. cit.}, in the case of sheaves of sets and noncommutative groups, are proved in Mme Raynaud's exposé (SGA 1, XIII).

Exposé III contains a proof of the Lefschetz--Verdier formula for cohomological correspondences between complexes of sheaves on schemes of finite type over a field. In Exposé III B, one shows how to compute the local terms of the formula in certain cases, for instance that of the Frobenius correspondence on curves. We refer the reader to the introductions of Exposés III and III B for more details on their contents.

Exposé IV, by A. Grothendieck, devoted to the construction of the cohomology class associated with a cycle, was written up by P. Deligne and published in SGA $4^{1/2}$.

After technical preliminaries in Exposé V on adic inverse systems, Exposé VI defines the notion of a constructible $\Z_\ell$-sheaf, respectively a constructible twisted constant sheaf, that is, a ``smooth'' sheaf in Deligne's terminology ([1], SGA $4^{1/2}$, Arcata). It then extends to constructible $\Z_\ell$-sheaves, by passage to the limit from $\Z/\ell^n$-sheaves, the formalism of higher direct images with proper supports. Exposé VI does not, however, give a definition of a derived category of $\Z_\ell$-sheaves, with operations $\otimes^{\mathrm L}$ and $\RHom$ extending those available in the derived category of $\Z/\ell^n$-sheaves. A fortiori, the extension to $\Z_\ell$-sheaves of the duality formalism of SGA 4, XVIII is not considered. This question was the subject of Jouanolou's unpublished thesis. The reader may nevertheless consult the beginning of [1], where the sheaves $\underline{\operatorname{Ext}}^i$ of $\Z_\ell$-sheaves are defined.

Exposé VII, independent of the rest of the seminar, is undoubtedly one of the most interesting. It contains: (a) the computation of the cohomology of several standard varieties, namely affine spaces, projective spaces, and flag varieties; (b) an account of the theory of Chern classes in étale cohomology; (c) a proof of the self-intersection formula
\[
i^*i_*x=x\,c_d(N)
\]
in étale cohomology, respectively in the Chow ring, the proof in the latter case being due to Mumford, together with various applications of this formula, including computation of the cohomology of blow-ups and the Gauss--Bonnet formula.

Exposé VIII contains sorites on finiteness conditions in derived categories and Grothendieck groups. These questions are taken up again in a more general framework in SGA 6, I and IV.

Exposé IX, by J.-P. Serre, entitled ``Introduction to Brauer theory'', was published in [4] and is not reproduced in this volume. One of its principal results, due to Swan, concerning the existence of a projective module having as character the ``Swan character'', is used in Exposé X, where an Euler--Poincaré formula is proved for a sheaf on a curve, slightly more general than the one appearing in Raynaud's exposé [3].

Exposé XI, written by I. Bucur, does not exist: it was lost during a move, and the author had no copy of it. It was devoted to Grothendieck's theory of commutative traces, generalizing that of Stallings [5]. A more sophisticated version of this theory, in a sheaf-theoretic framework, is presented in (III B 5, 6). Applications to Lefschetz formulae are given in Exposés XII and III B. The trace formula of Exposé XII is proved independently of the general formula of Exposé III, but it is shown in (III B 6) that the local terms appearing there are indeed those of the general formula, and that the latter implies it.

The absence of Exposé XIII is due only to numbering: the contents of Exposés X through XII had initially been planned to extend over four exposés.

Exposé XIV, sometimes also numbered XV, contains generalities on the Frobenius correspondence, including in particular the comparison between ``geometric'' Frobenius and ``arithmetic'' Frobenius. Starting from the trace formula of Exposé XII, or from III B, it gives Grothendieck's proof of the rationality of $L$-functions. The two important points are reduction to the case of curves, and the passage to the limit that makes it possible to reduce to a formula for $\Z/\ell^n$-sheaves.

During the oral seminar, Grothendieck gave an exposé on open problems and stated a few conjectures. That exposé, which was intended to close the seminar, was unfortunately not written up, nor was his very beautiful introductory exposé, which surveyed the Euler--Poincaré and Lefschetz formulae in various contexts: topological, complex analytic, and algebraic.

I thank P. Deligne for having convinced me to write, in a new version of Exposé III, a proof of the Lefschetz--Verdier formula, thereby removing one of the obstacles to the publication of this seminar. I am very grateful to him for the improvements in exposition that he suggested to me, and for the help he gave me in preparing this volume for the publisher. I also thank Mme Lièvremont, who typed Exposés III and III B carefully and in a very short time.

I cannot end this introduction without paying tribute to the memory of I. Bucur, who died of cancer in September 1976. Those who, like me, had the privilege of having him as a friend will not forget his exquisite kindness, his smiling humor, and the charm of his conversation.

\begin{flushright}
Paris, 19 February 1977\\
Luc Illusie
\end{flushright}

\bigskip


\clearpage
\section*{SGA 5, Exposé I, Introduction and Sections 1--2}
\addcontentsline{toc}{section}{SGA 5, Exposé I, Introduction and Sections 1--2}

\section*{Introduction}

This exposé is devoted to the biduality theorem in étale cohomology. The theory developed here and in SGA A, Exposé XVIII (where SGA A denotes SGA 4), for the étale topology and constructible torsion sheaves on preschemes, is formally very close to the theory developed in Hartshorne~[H] for the Zariski topology and coherent sheaves; the latter served, to a very large extent, as its model. But while, in the case of ``continuous'' coefficients, the existence of dualizing complexes presents no special difficulty, in the case of ``discrete'' coefficients it is much more delicate to prove, and seems to require resolution of singularities in an essential way. See, however, Deligne's biduality theorem (SGA $4^{1/2}$, Finiteness Theorem 4.3).

Section 1 contains the definition of dualizing complexes and the formal properties which follow from it, notably the interpretation of the dualizing functors as ``exchangers of functors.''

In Section 2 one proves that, on a connected locally noetherian prescheme, a dualizing complex is determined uniquely, up to a translation of degrees and tensoring by an invertible sheaf.

In Section 3 we come to the central theorem of the exposé, Theorem 3.4.1, which asserts that on a ``good'' regular prescheme $X$---for example, an excellent noetherian regular prescheme of characteristic zero\footnote{This last restriction will become unnecessary once Hironaka-style resolution of singularities and the purity theorem are available for such preschemes.}---the complex $(\ZZ/n\ZZ)_X$ is dualizing.

Section 4 gives the applications of the theory to the local duality theorem. This theorem says that, on a ``good'' prescheme $X$ endowed with a closed point $x$ whose residue field is separably closed, and for a constructible sheaf $F$, there is a perfect pairing between $\underline{\Hh}^i_x(F')$ and $\underline{\Hh}^{-i}_x(F)$, where $F'$ is the dual of $F$, that is,
\[
   F'=\R\,\underline{\Hom}(F,K_X),
\]
with $K_X$ a dualizing complex.

Finally, in Section 5, one proves directly that on a regular prescheme of dimension $1$ the complex $(\ZZ/n\ZZ)_X$ is dualizing, for $n$ prime to the residual characteristics.

\section*{1. Definition and formal properties of dualizing complexes}

Fix a coefficient ring $A=\ZZ/n\ZZ$. If $X$ is a prescheme, we denote by $A_X$ the constant sheaf on $X_{\mathrm{\acute et}}$ with value $A$, and by $\D(X)$ the derived category of the category of $A_X$-modules. We denote by $\D_c(X)$ the full triangulated subcategory of $\D(X)$ consisting of those complexes $F$ such that each $\Hh^i(F)$ is a constructible $A_X$-module. We put
\[
   \D_c^{*}(X)=\D^{*}(X)\cap\D_c(X),
\]
where $*=\varnothing,+,-,$ or $b$.

\begin{statement}{Definition 1.1}
An $A_X$-module $I$ is called \emph{quasi-injective} if
\[
   \underline{\Ext}^{i}(F,I)=0
\]
for every constructible $A_X$-module $F$ and every $i>0$.
\end{statement}

\begin{statement}{Remarks 1.2}
\begin{enumerate}[label=\alph*)]
\item In contrast with what happens for ``continuous'' coefficients, cf. [H], II.7.17, a quasi-injective sheaf is not in general injective. For example, let $X=\Spec(\RR)$ and $A=\ZZ/2\ZZ$. Then $A_X$ is quasi-injective, but is not injective and is not even of finite injective dimension, since
\[
   \Hh^*(X;A_X)=\Hh^*(\ZZ/2\ZZ;\ZZ/2\ZZ)
\]
is a polynomial algebra on one generator of dimension $1$.

\item Let $I$ be a quasi-injective $A_X$-module. It is false in general that the relation $\underline{\Ext}^i(F,I)=0$ for $i>0$, valid for constructible $F$, remains valid for all $F$. Indeed, take $A=\FF_\ell$. Let $k$ be a field, let $\ol{k}$ be a separable closure, let
\[
   f:\ol{x}=\Spec(\ol{k})\longrightarrow X=\Spec(k)
\]
be the canonical morphism, and suppose that, for every connected étale cover $X'\to X$, there exists a connected étale cover $X''\to X'$ such that $[X'':X']$ is divisible by $\ell$; one may take, for example, $k=\mathbb Q$. Consider the exact sequence
\[
   0\longrightarrow A_X \xrightarrow{\,\varepsilon\,} f_*f^*A_X\longrightarrow F\longrightarrow0,
   \tag{*}
\]
defined by the adjunction morphism $\varepsilon$. In view of the hypothesis, the section of $\Ext^1(F,A_X)$ defined by $(*)$ does not vanish over any étale $U$ over $X$; hence
\[
   \underline{\Ext}^1(F,A_X)_{\ol{x}}\ne0.
\]
But one checks immediately that $A_X$ is quasi-injective.
\end{enumerate}
\end{statement}

\begin{statement}{Proposition 1.3, cf. [H] 1.7.6}
Let $X$ be a prescheme and $N\in\ZZ$. The following conditions on $K\in\ob\D^+(X)$ are equivalent:
\begin{enumerate}[label=(\roman*)]
\item $K$ is isomorphic, in $\D(X)$, to a bounded complex $I$ such that each $I^i$ is quasi-injective and $I^i=0$ for $i>N$;
\item for every $F\in\ob\D_c^b(X)$ such that $\Hh^i(F)=0$ for $i<0$, one has
\[
   \underline{\Ext}^i(F,K)=0\qquad(i>N);
\]
\item for every constructible $A_X$-module $F$, one has
\[
   \underline{\Ext}^i(F,K)=0\qquad(i>N).
\]
\end{enumerate}
\end{statement}

\begin{prooftext}
The implication (i)$\Rightarrow$(ii) is proved as follows. Let $F\in\ob\D_c^b(X)$ with $\Hh^i(F)=0$ for $i<0$, and let $I\in\ob\D^+(X)$ be bounded, with each $I^i$ quasi-injective and $I^i=0$ for $i>N$. We must show that $\underline{\Ext}^i(F,I)=0$ for $i>N$. By induction on the number of indices $i$ for which $\Hh^i(F)\ne0$, using dévissage, one reduces to the case $N=0$ and $I$ concentrated in degree $0$. Then there is a biregular spectral sequence
\[
   E_2^{pq}=\underline{\Ext}^p(\Hh^{-q}(F),I)\Rightarrow \underline{\Ext}^*(F,I).
\]
Since $I$ is quasi-injective, $E_2^{pq}=0$ for $p>0$ and all $q$, hence
\[
   E_2^{0i}=\underline{\Hom}(\Hh^{-i}(F),I)\xrightarrow{\sim}\underline{\Ext}^{i}(F,I),
\]
and the assertion follows.

The implication (ii)$\Rightarrow$(iii) is trivial.

For (iii)$\Rightarrow$(i), we may suppose that $K$ is bounded below and has injective components. Applying the hypothesis to $F=A_X$ gives $\Hh^i(K)=0$ for $i>N$. We may therefore replace $K$ by an isomorphic complex $K'$, bounded, with $K'^i$ injective for $i<N$ and $K'^i=0$ for $i>N$. One then observes that $K'^N$ is quasi-injective, which completes the proof. \qedtext
\end{prooftext}

\begin{statement}{Remark 1.4}
The editor does not know whether Proposition 1.3 remains true if, in condition (ii), the hypothesis $F\in\ob\D_c^b(X)$ is replaced by $F\in\ob\D_c^+(X)$.
\end{statement}

\begin{statement}{Definition 1.5}
The \emph{quasi-injective dimension} of $K$, denoted
\[
   \dim.\,q.\,inj(K),
\]
is the greatest lower bound, in $\overline{\RR}$, of the integers $N$ satisfying the equivalent conditions of Proposition 1.3.
\end{statement}

If $K$ has finite quasi-injective dimension, then the functor
\[
   \R\,\underline{\Hom}(\,\cdot\,,K):\D(X)\longrightarrow\D(X)
\]
induces a functor $\D_c^b(X)\to\D^b(X)$.

Later we shall see that if $X$ is smooth of dimension $d$ over a field $k$ and $n$ is prime to $\operatorname{char}(k)$, then
\[
   \dim.\,q.\,inj(A_X)=2d.
\]

The example of Remark 1.2(a) shows that a complex of finite quasi-injective dimension need not have finite injective dimension. Nevertheless one has the following result.

\begin{statement}{Proposition 1.5.2}
Let $X$ be a locally noetherian prescheme and let $K\in\ob\D^+(X)$. Suppose that there exists an integer $N$ such that $\cd_{\mathbf n}(X)\le N$, where $\mathbf n$ is the set of prime divisors of $n$ in the notation of SGA A, Exposé X, 1. Then, for $N'\in\ZZ$,
\[
   \dim.\,q.\,inj(K)\le N'\quad\Longrightarrow\quad \dim.\,inj(K)\le N+N'.
\]
\end{statement}

We shall use the following lemma in the proof.

\begin{statement}{Lemma 1.5.1}
Let $C$ be an abelian category satisfying (AB5) and having a family of generators $(U_i)_{i\in I}$. Let $K\in\ob\D^+(C)$ and $N\in\ZZ$. The following conditions are equivalent:
\begin{enumerate}[label=(\roman*)]
\item $\dim.\,inj(K)\le N$;
\item for every $i\in I$ and every quotient $F$ of $U_i$,
\[
   \Ext^n(F,K)=0\qquad(n>N).
\]
\end{enumerate}
\end{statement}

\begin{prooftext}
It is enough to prove (ii)$\Rightarrow$(i). By a well-known argument, cf. for example [H] 1.7.6, one is reduced to proving that if $K'\in\ob(C)$ is such that $\Ext^1(F,K')=0$ whenever $F$ is a quotient of one of the $U_i$, then $K'$ is injective. Equivalently, if every morphism from a subobject $V$ of some $U_i$ to $K'$ extends to $U_i$, then $K'$ is injective. This follows at once from the proof of Lemma 1, p. 136, of Tohoku. \qedtext
\end{prooftext}

\begin{prooftext}[Proof of Proposition 1.5.2]
Assume $\dim.\,q.\,inj(K)\le N'$ and prove $\dim.\,inj(K)\le N+N'$. Since the category of $A_X$-modules is generated by the sheaves $A_{U,X}$, where $U$ is étale of finite type over $X$, Lemma 1.5.1 reduces us to proving
\[
   \Ext^i(X;F,K)=0\qquad(i>N+N')
\]
when $F$ is a quotient of such an $A_{U,X}$, hence constructible by SGA A, Exposé IX, 2.9. Consider the spectral sequence
\[
   E_2^{pq}=\Hh^p\bigl(X;\underline{\Ext}^{q}(F,K)\bigr)\Rightarrow\Ext^*(X;F,K).
\]
The hypothesis $\cd_{\mathbf n}(X)\le N$ implies $E_2^{pq}=0$ for $p>N$, while $\dim.\,q.\,inj(K)\le N'$ implies $E_2^{pq}=0$ for $q>N'$. Thus $E_2^{pq}=0$ for $p+q>N+N'$, giving the desired assertion. \qedtext
\end{prooftext}

\begin{statement}{Remark 1.5.3}
The hypothesis of Proposition 1.5.2 is satisfied, with $N=2d$, if $X$ is a prescheme of finite type and dimension $d$ over a separably closed field $k$, with $n$ prime to $\operatorname{char}(k)$, by SGA A, Exposé X, 4.
\end{statement}

\begin{statement}{Proposition 1.6}
Let $f:X\to Y$ be a separated quasi-finite morphism of noetherian preschemes, and let $K\in\ob\D^+(Y)$. Then, for $N\in\ZZ$,
\[
   \dim.\,q.\,inj(K)\le N\quad\Longrightarrow\quad
   \dim.\,q.\,inj(\R^!f(K))\le N.
\]
\end{statement}

\begin{prooftext}
By the Main Theorem of EGA IV, 8.12.6, there is a factorization $f=uf'$, where $f'$ is an open immersion and $u$ is finite. Since
\[
   \R^!f=\R^!f'\,\R^!u,
\]
we are reduced to treating separately the case of an open immersion and the case of a finite morphism. If $f$ is an open immersion, the assertion is trivial, since every constructible sheaf on $X$ is then induced by a constructible sheaf on $Y$. Suppose then that $f$ is finite, and let $F$ be a constructible $A_X$-module. The duality isomorphisms of SGA A, Exposé XVIII, 3.1.9.6 give
\[
   f_*\underline{\Ext}^i\bigl(F,\R^!f(K)\bigr)\xrightarrow{\sim}\underline{\Ext}^{i}(f_*F,K).
\]
Since $f_*F$ is constructible by SGA A, Exposé IX, 2.14, it follows that if $\dim.\,q.\,inj(K)\le N$, then
\[
   f_*\underline{\Ext}^i\bigl(F,\R^!f(K)\bigr)=0\qquad(i>N),
\]
and hence $\underline{\Ext}^i(F,\R^!f(K))=0$ for $i>N$, by SGA A, Exposé VIII, 5.5. \qedtext
\end{prooftext}

Let now $X$ be a prescheme and $K\in\ob\D^+(X)$. We denote by $\underline D_K$ the functor
\[
   \R\,\underline{\Hom}(\,\cdot\,,K).
\]
We are going to define, for $F\in\ob\D(X)$, a functorial morphism
\[
   F\longrightarrow \underline D_K\underline D_KF,
\]
cf. [H], V.1.2. The construction works more generally on a ringed topos. We may first suppose that $K$ is represented by a complex of injectives, so that the symbols $\R$ may be suppressed. The identity morphism
\[
   \underline{\Hom}(F,K)\longrightarrow\underline{\Hom}(F,K)
\]
defines, by Cartan's adjunction formula, a morphism
\[
   F\otimes\underline{\Hom}(F,K)\longrightarrow K,
\]
and a second application of the same formula gives
\[
   F\longrightarrow\underline{\Hom}\bigl(\underline{\Hom}(F,K),K\bigr).
\]
This is the announced morphism, which we shall call the canonical morphism. We say that the pair $(F,K)$ is \emph{bidualizing}, or \emph{satisfies biduality}, or that $F$ is \emph{reflexive relative to $K$}, if the canonical morphism
\[
   F\longrightarrow \underline D_K\underline D_KF
\]
is an isomorphism.

From now on, unless expressly stated otherwise, all preschemes under consideration are assumed locally noetherian.

\begin{statement}{Definition 1.7}
Let $X$ be a prescheme. A complex $K\in\ob\D^+(X)$ is called \emph{dualizing} if the following conditions hold:
\begin{enumerate}[label=(\roman*)]
\item $K$ has finite quasi-injective dimension;
\item for every $F\in\ob\D_c^-(X)$, one has $\underline D_K(F)\in\ob\D_c^+(X)$;
\item every $F\in\ob\D_c^b(X)$ is reflexive relative to $K$.
\end{enumerate}
\end{statement}

\begin{statement}{Remarks 1.8}
\begin{enumerate}[label=\alph*)]
\item By standard dévissage, truncating $F$ and using the fact that $\D_c(X)$ is triangulated, condition (ii) is equivalent to:

(ii bis) for every constructible $A_X$-module $F$ and every $i\in\ZZ$, $\underline{\Ext}^i(F,K)$ is constructible.

\item Condition (ii bis) implies $K\in\ob\D_c^+(X)$. Unlike the example of coherent sheaves, EGA $0_{\mathrm{III}}$, 12.3.3, the converse is perhaps not automatic. It is true, however, in the ``good'' cases considered below; for example when resolution of singularities and purity are available on $X$, in particular if $X$ is an excellent prescheme of characteristic zero.

\item By dévissage on $F$, condition (iii) is equivalent to:

(iii bis) every constructible $A_X$-module $F$ is reflexive relative to $K$.

\item If $K\in\ob\D^+(X)$ satisfies (i) and (ii), respectively if $K$ is dualizing, then $\underline D_K$ induces a functor, respectively an equivalence of categories,
\[
   (\D_c^b(X))^\circ\longrightarrow\D_c^b(X).
\]

\item If $K\in\ob\D^+(X)$ has finite injective dimension and satisfies (ii), then for every $F\in\ob\D_c(X)$ one has $\underline D_K(F)\in\D_c(X)$, and $\underline D_K$ induces a functor
\[
   (\D_c^+(X))^\circ\longrightarrow\D_c^-(X).
\]
If, moreover, $K$ is dualizing, every $F\in\ob\D_c(X)$ is reflexive relative to $K$, and $\underline D_K$ induces equivalences of categories
\[
   (\D_c(X))^\circ\xrightarrow{\sim}\D_c(X),\qquad
   (\D_c^-(X))^\circ\xrightarrow{\sim}\D_c^+(X),\qquad
   (\D_c^+(X))^\circ\xrightarrow{\sim}\D_c^-(X).
\]
It is not known whether this conclusion remains valid without the restriction that $K$ have finite injective dimension.
\end{enumerate}
\end{statement}

We now turn to the ``exchange formulas.'' We shall not return to the notion of the tor-dimension of a complex, see SGA A, Exposé XVII, 4.1.9, and also SGA 6, I, 5. Retain only the following point: if $X$ is a prescheme and $G\in\ob\D^-(X)$, the condition $\underline{\Tor}_i(F,G)=0$ for $i>N$ and all $A_X$-modules $F$ is equivalent to the same condition with $F$ constructible, by SGA A, Exposé IX, 2.9 and the commutation of tensor product with filtered inductive limits. Thus there is no need to introduce a notion of ``quasi-tor-dimension.''

\begin{statement}{Lemma 1.9, cf. [H] II.4.3}
\begin{enumerate}[label=\alph*)]
\item Let $X$ be a prescheme and let $F,G\in\ob\D_c^-(X)$. One has $F\otimes^{\mathrm L}G\in\ob\D_c(X)$ if $F\in\ob\D_c^-(X)$, or if $G$ has finite tor-dimension.
\item Let $f:X\to Y$ be a morphism of preschemes and let $G\in\ob\D_c^-(Y)$. Then $f^*(G)\in\ob\D_c(X)$.
\end{enumerate}
\end{statement}

\begin{prooftext}
Part (b) follows trivially from SGA A, Exposé IX, 2.4. We prove (a). Either hypothesis on $(F,G)$ implies that the spectral sequence
\[
   E_2^{pq}=\sum_{q_1+q_2=q}\underline{\Tor}_{-p}\bigl(\Hh^{q_1}(F),\Hh^{q_2}(G)\bigr)
       \Rightarrow \Hh^*(F\otimes^{\mathrm L}G)
\]
is biregular. Hence one is reduced to the case where $F$ and $G$ are concentrated in degree $0$, that is, are constructible $A_X$-modules. The question is local, so one may suppose $X$ noetherian. Then $F$ admits, by SGA A, Exposé IX, 2.7, a left resolution $F'$ in which the $F'^i$ are of the form $A_{U,X}$ with $U$ étale and of finite type over $X$. We have
\[
   F\otimes^{\mathrm L}G\xleftarrow{\sim}F'\otimes G,
\]
and the result follows since the $A_{U,X}\otimes G$ are constructible. \qedtext
\end{prooftext}

\begin{statement}{Notation 1.10}
If $X$ is a prescheme, we denote by $\D^b(X)_{\mathrm{torf}}$, respectively $\D^b(X)_{\mathrm{qinjf}}$, the full subcategory of $\D(X)$ consisting of objects of finite tor-dimension, respectively of finite quasi-injective dimension.
\end{statement}

\begin{statement}{Proposition 1.11, cf. [H] V.2.6}
Let $X$ be a prescheme and $K\in\ob\D^+(X)$.
\begin{enumerate}[label=\alph*)]
\item For $F,G\in\ob\D^-(X)$, respectively for $F\in\ob\D(X)$ and $G\in\ob\D^b(X)_{\mathrm{torf}}$, there is a functorial isomorphism
\[
   \underline D_K(F\otimes^{\mathrm L}G)\xrightarrow{\sim}
   \R\,\underline{\Hom}(F,\underline D_K(G)).
\]

\item If $K$ satisfies conditions (i) and (ii) of Definition 1.7, then
\[
   G\in\ob\D_c^b(X)_{\mathrm{torf}}
   \Longrightarrow
   \underline D_K(G)\in\ob\D_c^b(X)_{\mathrm{qinjf}}
\]
and $\underline D_K(G)$ satisfies condition 1.7(ii).

\item If $K$ is dualizing, then for $F\in\ob\D_c^-(X)$ and $G\in\ob\D_c^b(X)$, respectively for $F\in\ob\D(X)$ and $G\in\ob\D_c^b(X)$ with $\underline D_K(G)\in\ob\D_c^b(X)_{\mathrm{torf}}$, there is a functorial isomorphism
\[
   \underline D_K\bigl(F\otimes^{\mathrm L}\underline D_K(G)\bigr)
      \xrightarrow{\sim}\R\,\underline{\Hom}(F,G).
\]
Moreover,
\[
   \underline D_K(G)\in\ob\D_c^b(X)_{\mathrm{torf}}
   \Longrightarrow
   G\in\ob\D_c^b(X)_{\mathrm{qinjf}},
\]
and $G$ satisfies condition 1.7(ii).

\item If $K$ is dualizing and has finite injective dimension, then the implications in (b) and (c) are equivalences.
\end{enumerate}
\end{statement}

\begin{prooftext}
Part (a) is precisely Cartan's adjunction formula, SGA 6, I, 7.4.

For (b), let $G\in\ob\D_c^b(X)_{\mathrm{torf}}$. If $F$ is a variable constructible $A_X$-module, then the cohomology objects of $F\otimes^{\mathrm L}G$ are constructible, by Lemma 1.9, and vanish outside two fixed bounds independent of $F$. The same therefore holds for the cohomology objects of $\underline D_K(F\otimes^{\mathrm L}G)$, and (a) gives the assertion.

For (c), the first assertion follows from (a) applied to $\underline D_K(G)$ in place of $G$, together with the reflexivity of $G$. The second follows by a similar argument to (b).

For (d), prove for instance the converse implication of (b). Suppose that $\underline D_K(G)\in\ob\D_c^b(X)_{\mathrm{qinjf}}$ and satisfies condition 1.7(ii). Then, as $F$ varies through constructible modules, the cohomology objects of $\underline D_K(F\otimes^{\mathrm L}G)$ are constructible and vanish outside two fixed bounds. The same is consequently true for the cohomology objects of $\underline D_K\underline D_K(F\otimes^{\mathrm L}G)$; but
\[
   F\otimes^{\mathrm L}G\xrightarrow{\sim}\underline D_K\underline D_K(F\otimes^{\mathrm L}G)
\]
because $K$ is dualizing and of finite injective dimension, cf. 1.8(d). Hence $G$ has finite tor-dimension. The converse implication in (c) is proved analogously. \qedtext
\end{prooftext}

\begin{statement}{Remarks 1.11.1}
\begin{enumerate}[label=\alph*)]
\item The editor does not know whether assertion (d) of Proposition 1.11 is valid under the sole hypothesis that $K$ be dualizing.
\item One expresses the preceding proposition pictorially by saying that the dualizing functor $\underline D_K$ exchanges the functors $\otimes^{\mathrm L}$ and $\R\,\underline{\Hom}$, as well as the notions of finite tor-dimension and finite quasi-injective dimension. The exchange is, however, fully satisfactory only when $K$ has finite injective dimension.
\end{enumerate}
\end{statement}

\begin{statement}{Proposition 1.12}
Let $Y$ be a noetherian prescheme, let $f:X\to Y$ be a morphism of finite type, and let $K\in\ob\D^+(Y)$. Suppose that $n$ is prime to the residual characteristics of $Y$, and put
\[
   K_X=\R^!f(K_Y),\qquad \underline D_X=\underline D_{K_X},\qquad \underline D_Y=\underline D_{K_Y},
\]
in the notation of SGA A, Exposé XVIII, 3.1.
\begin{enumerate}[label=\alph*)]
\item For $F\in\ob\D^-(X)$ there is a functorial isomorphism
\[
   \R f_*\underline D_X(F)\xrightarrow{\sim}\underline D_Y\R f_*(F).\tag{i}
\]
If, moreover, $K_X$ and $K_Y$ are dualizing,\footnote{We shall see below that under fairly general conditions $K_Y$ dualizing implies $K_X$ dualizing.} then for $F\in\ob\D_c^b(X)$ there is a functorial isomorphism
\[
   \R f_!\underline D_X(F)\xrightarrow{\sim}\underline D_Y\R f_*(F).\tag{ii}
\]

\item For $F\in\D_c^-(Y)$ there is a functorial isomorphism\footnote{If the functor $\R^!f$ has finite cohomological dimension, as for example when $Y$ is of finite type over a field (SGA A, Exposé X, 4.3 and Exposé XVIII, 3.1.7), the isomorphism (b)(i) is valid for $F\in\ob\D_c(Y)$.}
\[
   \R^!f\,\underline D_Y(F)\xrightarrow{\sim}\underline D_X f^*(F).\tag{i}
\]
If, moreover, $K_X$ and $K_Y$ are dualizing, then for $F\in\D_c^b(Y)$ there is a functorial isomorphism
\[
   f^*\underline D_YF\xrightarrow{\sim}\underline D_X\R^!f\,F.\tag{ii}
\]
\end{enumerate}
\end{statement}

\begin{prooftext}
The isomorphism (a)(i) is just the duality isomorphism of SGA A, Exposé XVIII, 3.1.9.6.

The isomorphism (a)(ii) is obtained by applying $\underline D_Y$ to both sides of (a)(i), written with argument $\underline D_XF$, and using the finiteness theorem SGA A, Exposé XVII, 5.3.6, which asserts that $\R f_*(\underline D_X(F))\in\ob\D_c(Y)$.

The isomorphism (b)(i) is the induction formula of SGA A, Exposé XVIII, 3.1.12.2.

The isomorphism (b)(ii) is obtained by applying $\underline D_X$ to both sides of (b)(i), written with argument $\underline D_YF$, and using 1.9(b). \qedtext
\end{prooftext}

\begin{statement}{Scholium 1.12.1}
In figurative language, the dualizing functors $\underline D_X$ and $\underline D_Y$ exchange the usual and unusual direct images, $\R f_*$ and $\R f_!$, and also the usual and unusual inverse images, $f^*$ and $\R^!f$.
\end{statement}

\begin{statement}{Corollary 1.13}
Under the hypotheses of Proposition 1.12, suppose in addition that $f$ is proper, and let $F\in\ob\D^-(X)$. If $(F,K_X)$ is bidualizing, then $(\R f_*F,K_Y)$ is bidualizing. Conversely, if $f$ is finite and $(\R f_*F,K_Y)$ is bidualizing, then $(F,K_X)$ is bidualizing.
\end{statement}

\begin{prooftext}
Suppose $(F,K_X)$ is bidualizing. Applying $\R f_*$ to the canonical isomorphism $F\xrightarrow{\sim}\underline D_X^2F$, one obtains
\begin{align*}
\R f_*F &\xrightarrow{\sim}\R f_*\underline D_X\underline D_XF \\
&\xrightarrow{\sim}\underline D_Y\R f_*\underline D_XF \qquad \text{by (a)(i)}\\
&\xrightarrow{\sim}\underline D_Y\underline D_Y\R f_*F \qquad \text{by (a)(i)}.
\end{align*}
The courageous reader will convince himself that the isomorphism obtained is the canonical one; this proves the first assertion.

Now suppose $f$ finite and $(f_*F,K_Y)$ bidualizing, since $\R f_*=f_*$ for $f$ finite. We prove that the canonical morphism $F\to\underline D_X^2F$ is an isomorphism. By Lemma 1.14(a) below, it suffices to show that the induced morphism
\[
   f_*F\longrightarrow f_*\underline D_X^2F
\]
is an isomorphism. But by applying (a)(i) twice, one convinces oneself that this morphism is precisely the canonical isomorphism
\[
   f_*F\longrightarrow \underline D_Y^2f_*F,
\]
which completes the proof. \qedtext
\end{prooftext}

\begin{statement}{Lemma 1.14}
Let $f:X\to Y$ be a finite morphism of preschemes.
\begin{enumerate}[label=\alph*)]
\item For a morphism $u:F\to G$ in $\D(X)$, $u$ is an isomorphism if and only if $f_*u$ is an isomorphism.
\item For $F\in\ob\D(X)$, one has $F\in\ob\D_c(X)$ if and only if $f_*F\in\ob\D_c(Y)$.
\end{enumerate}
\end{statement}

\begin{prooftext}
For (a), since $f_*$ is exact, one may restrict to the case where $F$ and $G$ are concentrated in degree zero. The assertion then follows immediately from the computation of the fibres of $f_*F$ and $f_*G$, SGA A, Exposé VIII, 5.5.

For (b), again one may restrict to the case where $F$ is concentrated in degree zero. Necessity is known, SGA A, Exposé IX, 2.14. We prove sufficiency: $f_*F$ constructible implies $F$ constructible. The question is local upstairs, hence a fortiori downstairs, so we may suppose $Y$ and hence $X$ affine. By EGA IV, 17.16.6, there exists a finite family $(Y_i)$ of affine subpreschemes of $Y$, pairwise disjoint and with union $Y$, such that, if $X_i=X\times_YY_i$, the induced morphism $f_i:X_i\to Y_i$ factors as
\[
   X_i\xrightarrow{u_i}X'_i\xrightarrow{v_i}Y_i,
\]
where $u_i$ is finite radicial and surjective and $v_i$ is finite étale. Put $F|_{X_i}=F_i$. By the base-change theorem for finite morphisms, SGA A, Exposé VIII, 5.5, one has
\[
   f_{i*}(F_i)\xrightarrow{\sim}(f_*F)|_{Y_i},
\]
and by SGA A, Exposé IX, 2.8, one is reduced to proving the assertion for $(f_i,F_i)$. Put $F'_i=u_{i*}(F_i)$, so that $f_{i*}(F_i)=v_{i*}(F'_i)$. It is clear from the definition of constructibility, respectively follows from SGA A, Exposé VIII, 1.1, that $F_i$, respectively $F'_i$, is constructible if and only if $F'_i$, respectively $v_{i*}(F'_i)$, is constructible. This proves the claim. \qedtext
\end{prooftext}

\begin{statement}{Corollary 1.15}
Let $f:X\to Y$ be a separated quasi-finite morphism of noetherian preschemes. Let $K_Y\in\ob\D^+(Y)$, and put, as in Proposition 1.12,
\[
   K_X=\R^!f(K_Y),\qquad \underline D_X=\underline D_{K_X},\qquad \underline D_Y=\underline D_{K_Y}.
\]
If $K_Y$ is dualizing, then $K_X$ is dualizing.
\end{statement}

\begin{prooftext}
We already know, by Proposition 1.6, that if $K_Y$ has finite quasi-injective dimension, then the same is true of $K_X$. We show that if $K_Y$ satisfies condition (ii), respectively condition (iii), of Definition 1.7, then $K_X$ satisfies the same condition.

By the Main Theorem, EGA IV, 8.12.6, $f$ factors as $f=f'i$, where $i$ is an open immersion and $f'$ is finite. Since
\[
   \R^!f\simeq\R^!i\,\R^!f',
\]
one must treat separately the case of an open immersion and the case of a finite morphism. If $f$ is an open immersion, the assertion is trivial, since a constructible sheaf on $X$ is induced by a constructible sheaf on $Y$, and formation of $\R\,\underline{\Hom}$ commutes with restriction to an open.

Suppose then that $f$ is finite. It follows from Corollary 1.13 that if $K_Y$ satisfies (iii), then $K_X$ satisfies (iii), taking account of the finiteness theorem SGA A, Exposé IX, 2.14. It remains to prove that if $K_Y$ satisfies (ii), then $K_X$ satisfies (ii). Let $F\in\ob\D_c^-(X)$. We must show that $\underline D_X(F)\in\ob\D_c^+(X)$. By the finiteness theorem, $f_*F\in\ob\D_c^-(Y)$. Hence, by the hypothesis on $K_Y$,
\[
   \underline D_Y f_*F\in\ob\D_c^+(Y).
\]
But
\[
   \underline D_Y f_*F\xrightarrow{\sim}f_*\underline D_X(F)
\]
by Proposition 1.12(a)(i), and Lemma 1.14(b) completes the proof. \qedtext
\end{prooftext}

\section*{2. Uniqueness of the dualizing complex}

In this section we suppose that $A$ is local, that is,
\[
   A=\ZZ/\ell^\nu\ZZ,
\]
where $\ell$ is prime.

\begin{statement}{Theorem 2.1, cf. [H] V.3.1}
Let $X$ be a connected locally noetherian prescheme, and let $K$ be a dualizing complex on $X$ in the sense of Definition 1.7. Let also $K'\in\ob\D_c^+(X)$. Then $K'$ is dualizing if and only if there exist an invertible $A_X$-module $L$ and an integer $r$, necessarily determined up to unique isomorphism, such that
\[
   K'\xrightarrow{\sim}K\otimes^{\mathrm L}L[r].
\]
\end{statement}

\begin{prooftext}
Let $L$ be an invertible $A_X$-module and $r$ an integer, and put $L^\vee=\underline{\Hom}(L,A_X)$. For $F\in\ob\D(X)$ and $G\in\ob\D^+(X)$ there are canonical isomorphisms
\[
\R\,\underline{\Hom}\bigl(F\otimes^{\mathrm L}L^\vee[-r],G\bigr)
   \xrightarrow{\sim}
\R\,\underline{\Hom}\bigl(F,G\otimes^{\mathrm L}L[r]\bigr)
   \xrightarrow{\sim}
\R\,\underline{\Hom}(F,G)\otimes^{\mathrm L}L[r]. \tag{*}
\]
It follows immediately that, if $K'\simeq K\otimes^{\mathrm L}L[r]$, then $K'$ is dualizing.

Moreover, if we write $\underline D_X=\underline D$ and $\underline D_{K'}=\underline D'$, then for $F\in\ob\D(X)$ the isomorphism $(*)$ gives
\[
   \underline D'F\xrightarrow{\sim}\underline DF\otimes^{\mathrm L}L^\vee[-r].
\]
Since $\underline DK\simeq A_X$, one obtains
\[
   \underline D'K\xrightarrow{\sim}L^\vee[-r],
\]
which shows that $L^\vee[r]$, hence $L$ and $r$, are determined essentially uniquely by $K'$.

Conversely, suppose that $K'$ is dualizing, and put
\[
   P=\underline D'K=\underline D'\underline DA_X=\R\,\underline{\Hom}(K,K').
\]
We shall prove:
\begin{align*}
\text{(a)}\qquad K'&\xrightarrow{\sim}K\otimes^{\mathrm L}P,\\
\text{(b)}\qquad P&\xrightarrow{\sim}L[r]\quad\text{for an invertible }A_X\text{-module }L.
\end{align*}
This will complete the proof of the theorem.

First observe that, by Proposition 1.9(b)(ii), for $F\in\ob\D_c(X)$ there is an isomorphism
\[
   F\otimes^{\mathrm L}P\xrightarrow{\sim}\underline D'\underline DF. \tag{**}
\]
Applying $(**)$ to $F=K$ gives
\[
   K\otimes^{\mathrm L}P\xrightarrow{\sim}\underline D'\underline DK\xrightarrow{\sim}\underline D'A_X=K',
\]
which proves (a).

Applying $(**)$ to
\[
   F=P'=\underline D\underline D'A_X
\]
gives
\[
   P'\otimes^{\mathrm L}P
      =\underline D'\underline D\underline D\underline D'A_X
      \xrightarrow{\sim}\underline D'\underline D'A_X
      \xrightarrow{\sim}A_X.
\]
Thus (b) follows from the following lemma. \qedtext
\end{prooftext}

\begin{statement}{Lemma 2.2, cf. [H] V.3.3}
Let $X$ be a connected locally noetherian prescheme, and let $P,P'\in\ob\D_c^-(X)$ be such that
\[
   P\otimes^{\mathrm L}P'\xrightarrow{\sim}A_X.
\]
Then there exist an invertible $A_X$-module $L$ and an integer $r$ such that
\[
   P\xrightarrow{\sim}L[r],\qquad P'\xrightarrow{\sim}L^\vee[-r].
\]
\end{statement}

\begin{prooftext}
The proof has two steps.

First suppose that $X=\Spec(k)$, with $k$ separably closed. One merely repeats the easy proof given in [H]. The essential point is this: by EGA $0_{\mathrm I}$, 5.4.3, if $M$ and $M'$ are finite-type $A$-modules such that
\[
   M\otimes M'\xrightarrow{\sim}A,
\]
then $M$ and $M'$ are invertible and $M'\simeq M^\vee$. This is where the hypothesis that $A$ is local is used.

We now pass to the general case. By the special case just treated, for every $x\in X$ there exists an integer $r(x)$ such that
\[
   P_{\ol{x}}\simeq A_{\ol{x}}[r(x)],\qquad
   P'_{\ol{x}}\simeq A_{\ol{x}}[-r(x)].
\]
Equivalently, $\Hh^i(P_{\ol{x}})=0$, respectively $\Hh^i(P'_{\ol{x}})=0$, except for $i=-r(x)$, respectively $i=r(x)$, and the corresponding nonzero cohomology modules are isomorphic to $A_{\ol{x}}$. It remains to show that the function $x\mapsto r(x)$ is locally constant.

Assume this for the moment. Since $X$ is connected, $r$ is constant, say equal to $r_0$. Then
\[
   \Hh^i(P)=0\ (i\ne-r_0),\qquad \Hh^i(P')=0\ (i\ne r_0),
\]
and
\[
   \Hh^{-r_0}(P)\otimes\Hh^{r_0}(P')\xrightarrow{\sim}A.
\]
Put $L=\Hh^{-r_0}(P)$ and $L'=\Hh^{r_0}(P')$. Let $u:\ol{x}\to\ol{y}$ be a specialization morphism, and let
\[
   u^*:L_{\ol{y}}\to L_{\ol{x}},\qquad u'^*:L'_{\ol{y}}\to L'_{\ol{x}}
\]
be the corresponding specialization morphisms. There is a commutative square
\[
\begin{array}{ccc}
L_{\ol{y}}\otimes L'_{\ol{y}} & \xrightarrow{\sim} & A_{\ol{y}}\\
\downarrow {u^*\otimes u'^*} && \downarrow \\
L_{\ol{x}}\otimes L'_{\ol{x}} & \xrightarrow{\sim} & A_{\ol{x}}.
\end{array}\tag{*}
\]
Choose a basis $e$, respectively $f$, of $L_{\ol{x}}$, respectively $L_{\ol{y}}$, and let $e'$, respectively $f'$, be the dual basis of $L'_{\ol{x}}$, respectively $L'_{\ol{y}}$, meaning that $e\otimes e'$ and $f\otimes f'$ map to $1$. If $u^*(f)=\lambda e$ and $u'^*(f')=\lambda'e'$, the commutativity of $(*)$ gives $\lambda\lambda'=1$. Hence $u^*$ is an isomorphism, and $u'^*$ is the contragredient isomorphism. By SGA A, Exposé IX, 2.13, it follows that $L$ and $L'$ are locally constant, hence of finite presentation as $A_X$-modules. Consequently, for every $x\in X$ the canonical map
\[
   \underline{\Hom}(L,F)_{\ol{x}}\longrightarrow\Hom(L_{\ol{x}},F_{\ol{x}})
\]
is an isomorphism for every $A_X$-module $F$. Since the geometric fibres of $L$ are invertible, $L$ itself is invertible. Similarly $L'$ is invertible, and $L'\simeq L^\vee$. Hence
\[
   P\simeq L[-r_0],\qquad P'\simeq L^\vee[r_0].
\]

It remains only to prove that $r$ is locally constant. For $m\in\ZZ$, the condition $r(x)=m$ is equivalent to
\[
   \Hh^m(P')_{\ol{x}}\ne0.
\]
Thus, by SGA A, Exposé IX, 2.4, the function $r$ is locally constructible. Therefore, to verify that it is locally constant, EGA IV, 1.10.1, allows one to suppose that $X$ is local. Let $x$ be its closed point and $U$ the open complement.

We prove that $r$ is constant. By induction on the dimension of $X$, we may already suppose that $r(y)$ is constant and equal to $r_0$ for $y\in U$; the point is to prove $r(x)=r_0$.

Assume the contrary, and derive a contradiction. By hypothesis
\[
   P|_U\xrightarrow{\sim}L[r_0]
\]
with $L$ an invertible $A_U$-module. After translating degrees we may suppose $r_0=0$. On the other hand $r(x)\ne0$; write $r(x)=a>0$. Let
\[
   i:U\to X,
   \qquad
   j:x\to X
\]
be the canonical inclusions. Consider the exact sequence
\[
   0\longrightarrow i_!(P|_U)\longrightarrow P\longrightarrow j_*j^*P\longrightarrow0.\tag{**}
\]
By hypothesis the only nonzero cohomology objects of $P$ are $\Hh^0(P)$ and $\Hh^{-a}(P)$, and
\[
   \Hh^0(P)\xrightarrow{\sim}i_!L.
\]
It follows that $(**)$ splits naturally, via the map $P\to\Hh^0(P)$, hence
\[
   P\xrightarrow{\sim} i_!(P|_U)\oplus j_*j^*P.
\]
Therefore
\begin{align*}
P\otimes^{\mathrm L}P'
&\xrightarrow{\sim} i_!(P|_U)\otimes^{\mathrm L}P'
      \oplus j_*j^*P\otimes^{\mathrm L}P' \\
&\xrightarrow{\sim} i_!\bigl(P|_U\otimes^{\mathrm L}P'|_U\bigr)
      \oplus j_*\bigl(j^*P\otimes^{\mathrm L}j^*P'\bigr)\\
&\xrightarrow{\sim} i_!A_U\oplus j_*A_x,
\end{align*}
which is impossible, since $i_!A_U\oplus j_*A_x$ is plainly not isomorphic to $A_X$; already the degree-zero cohomology sheaves are not isomorphic. This proves Lemma 2.2, and hence Theorem 2.1. \qedtext
\end{prooftext}

\bigskip


\clearpage
\section*{SGA 5, Exposé I continued; Exposé III and III B to 5.10}
\addcontentsline{toc}{section}{SGA 5, Exposé I continued; Exposé III and III B to 5.10}

\section*{Exposé I. Dualizing Complexes, continued}

\section*{3. Existence of dualizing complexes}

\subsection*{3.1. Preliminaries}

As was said in the introduction, the object of this section is to prove that, on a ``good'' regular prescheme $X$, the sheaf $A_X$ is dualizing, where $n$ is assumed prime to the residual characteristics.

In fact, one would like this to hold on excellent regular preschemes of equal characteristic; but there is no doubt that the proof given below will extend to them once resolution of singularities and the purity theorem are available in that case.

We first present the various ingredients which will be used.

\begin{statement}{3.1.1. Singular locus}
Let $X$ be a prescheme, locally noetherian. We shall say that $X$ satisfies condition \emph{(reg)} if the following equivalent properties hold (EGA IV, 6.12.3):
\begin{enumerate}[label=(\roman*)]
\item for every sub-prescheme $Y$ of $X$, the set $\Reg(Y)$ of regular points of $Y$ is open in $Y$;
\item for every closed integral sub-prescheme $Y$ of $X$, the set $\Reg(Y)$ contains a non-empty open subset of $Y$.
\end{enumerate}
The condition (reg) is satisfied if $X$ is locally of finite type over a complete noetherian local ring, for example over a field (EGA IV, 6.12.8), or over a Dedekind ring whose field of fractions has characteristic zero, for example over $\ZZ$ (EGA IV, 6.12.6); more generally, it is satisfied if $X$ is excellent (EGA IV, 7.8.6).
\end{statement}

\begin{statement}{3.1.2. Dévissage of constructible sheaves}
Let $X$ be a noetherian prescheme, and let $C$ be a strictly full subcategory of the category $\operatorname{Cons}(X)$ of constructible $A_X$-modules. Suppose that $C$ is stable under direct factors and extensions. The following conditions are equivalent:
\begin{enumerate}[label=(\roman*)]
\item $C=\operatorname{Cons}(X)$;
\item $C$ contains sheaves of the form $i_!(F)$, where $i:Y\to X$ is the immersion of an integral sub-prescheme and $F$ is locally constant on $Y$;
\item $C$ contains sheaves of the form $f_*i_!(A'_U)$, where $i:U\to Y$ is an open subset of an integral prescheme, $f:Y\to X$ is a finite morphism such that the generic point of $Y$ is separable over its image, and $A'=\ZZ/\ell\ZZ$, for $\ell$ a prime divisor of $n$.
\end{enumerate}
If $X$ satisfies (reg) (3.1.1), these conditions are again equivalent to the following:
\begin{enumerate}[label=(\roman* bis)]
\item the analogue of (ii), with $Y$ moreover regular;
\item the analogue of (iii), with $U$ regular.
\end{enumerate}
If, in addition, $X$ is affine, then in (ii), respectively in (ii bis), one may suppose $Y$ closed in an open subset of the form $D(f)$, with $f\in \Gamma(X,\cO_X)$.
\end{statement}

\begin{prooftext}
The equivalence of (i)--(iii) was proved in SGA A, Exposé IX, 2.5 and 5.8. The remaining assertions follow by noetherian induction. \qedtext
\end{prooftext}

\begin{statement}{3.1.3. Cohomological dimension}
Let $X$ be a locally noetherian prescheme. We denote by $(\mathrm{cdloc})_n$ the following condition. For every sub-prescheme $Y$ of $X$, every geometric point $\ol y$ of $Y$, and every
\[
   f\in\Gamma\bigl(\wt Y(\ol y),\cO_{\wt Y,\ol y}\bigr),
\]
where $\wt Y(\ol y)$ denotes the strict localization of $Y$ at $\ol y$ (SGA A, Exposé VIII, 7.1), one has
\[
   \cd_{\mathbf n}\bigl(\wt Y(\ol y)-V(f)\bigr)
       \le \dim \cO_{Y,\ol y}
       \quad (=\dim \cO_{Y,y}\text{ by EGA IV, 6.1.3}),
\]
where $\mathbf n$ is the set of prime divisors of $n$.

The condition $(\mathrm{cdloc})_n$ is satisfied if $X$ is locally of finite type over a field (SGA A, Exposé XIV, 3.5), or is an excellent prescheme of characteristic zero (SGA A, Exposé XIX, 6.2).\footnote{Added in 1977: or, by a recent result of G. Ofer, if $X$ is excellent, regular, and of dimension $2$.}
\end{statement}

\begin{statement}{3.1.4. Absolute cohomological purity}
Let $X$ be a regular prescheme, and let $Y$ be a regular closed sub-prescheme of $X$, of codimension $d$ at every point. One says that the pair $(X,Y)$ satisfies the absolute cohomological purity theorem, relative to $n$, if
\[
   \uH^i_Y(A_X)=0\qquad (i\ne 2d),
\]
and if the ``local fundamental class'' homomorphism (SGA $4^{1/2}$, Cycle 2.2)
\[
   A_Y\longrightarrow \uH^{2d}_Y\bigl((\mu_n)^{\otimes d}_X\bigr)
\]
is an isomorphism.

One says that the absolute cohomological purity theorem is true on $X$ if every pair $(U,Y)$ as above satisfies it, where $U$ is an open subset of $X$.

Absolute purity is true on $X$ if $X$ is smooth over a perfect field of characteristic prime to $n$ (SGA A, Exposé XVI, 3.9), or if $X$ is an excellent prescheme of characteristic zero (SGA A, Exposé XIX, 3.2), or if $X$ is regular of dimension $1$ (cf. 5.1).\footnote{Added in 1977: or, by a recent result of G. Ofer, if $X$ is excellent, regular, and of dimension $2$.}
\end{statement}

\begin{statement}{3.1.5. Hironaka-style resolution of singularities}
\begin{enumerate}[label=\alph*)]
\item Let $X$ be a regular prescheme, and let $D\ge0$ be a divisor on $X$. One says that $D$ has \emph{strict normal crossings} if there exists a finite family $(f_i)_{i\in I}$ of elements of $\Gamma^*(X,\cO_X)$ such that
\[
   D=\sum_{i\in I}\Div(f_i)
\]
(EGA IV, 21), and such that, for every $x\in\Supp(D)$, the germs $(f_i)_x$ which lie in $\fm_x$ form part of a regular system of parameters. One says that $D$ has \emph{normal crossings} if, locally for the étale topology, $D$ has strict normal crossings.

The standard example of a normal crossings divisor is furnished by a locally finite union of coordinate hyperplanes in affine space.

\item Let $X$ be a locally noetherian prescheme. We shall say that $X$ is \emph{strongly desingularizable} if the following condition is satisfied: for every finite morphism $f:Y\to X$, where $Y$ is integral and its generic point is separable over its image, and every regular open subset $U$ of $Y$, there exists a proper birational morphism $h:Y'\to Y$, with $Y'$ regular, such that $h$ induces an isomorphism from $U'=h^{-1}(U)$ onto $U$, and $Y'-U'$ is a normal crossings divisor.
\end{enumerate}
Here are cases where $X$ is strongly desingularizable:
\begin{enumerate}[label=\alph*)]
\item $X$ is excellent of characteristic zero;
\item $X$ is locally noetherian, regular, and of dimension $\le1$;
\item $X$ is of finite type over $\ZZ$ or over a perfect field, and of dimension $\le2$.\footnote{Added in 1977: more generally, if $X$ is excellent noetherian of dimension $\le2$, by H. Hironaka, \emph{Desingularization of excellent surfaces}, notes by B. Bennett, Bowdoin College, 1967, to be supplemented if necessary by H. Hironaka, \emph{Certain numerical characters of singularities}, J. Math. Kyoto Univ. 10 (1970), 151--187.}
\end{enumerate}
\end{statement}

\begin{statement}{3.1.6. Cohomological finiteness}
Let $f:X\to Y$ be a morphism of preschemes, and let $F\in\ob\D_c^+(X)$. One has $\R f_*(F)\in\ob\D_c^+(Y)$ in the following cases:
\begin{enumerate}[label=\alph*)]
\item $f$ is proper and of finite presentation (SGA A, Exposé XIV, 1.1);
\item ``when resolution of singularities and the purity theorem are available'', in particular:
\begin{enumerate}[label=(\roman*)]
\item if $X$ and $Y$ are locally of finite type over a field $k$ of characteristic zero and $f$ is a $k$-morphism of finite type (SGA A, Exposé XVI, 5.1);
\item if $Y$ is excellent of characteristic zero and $f$ is of finite type (SGA A, Exposé XIX, 5.1);
\item if $X$ is locally of finite type over a field and $\dim(X)\le2$ (SGA A, Exposé XIX, 5.1).
\end{enumerate}
\end{enumerate}
\end{statement}

\subsection*{3.2. The quasi-injective dimension of $A_X$}

\begin{statement}{Proposition 3.2.1}
Let $X$ be a regular prescheme whose residual characteristics are prime to $n$. Suppose that $X$ satisfies (reg) (3.1.1) and $(\mathrm{cdloc})_n$ (3.1.3), and that absolute purity is true on $X$ (3.1.4). Then, for every constructible $A_X$-module $F$ and every $x\in X$, one has
\[
   \uExt^i(F,A_X)_x=0\qquad\text{for } i>2\dim\cO_{X,x}.\tag{*}
\]
\end{statement}

\begin{prooftext}
The question being local, we may suppose $X$ affine. By dévissage (3.1.2), it is enough to verify $(*)$ for $F=i_!(G)$, where $i:Y\to X$ is the immersion of a regular integral sub-prescheme, closed in an open subset $D(f)$, $f\in A(X)$, and $G$ is a locally constant constructible sheaf on $Y$. We must compute
\[
   \uExt^q(F,A_X)=\Hh^q\bigl(\R\,\underline{\Hom}(F,A_X)\bigr).
\]
The projection formula for $i$ (SGA A, Exposé XVIII) gives
\[
 \R\,\underline{\Hom}(i_!G,A_X)
   \xrightarrow{\sim}
 \R i_*\R\,\underline{\Hom}(G,\R^!i(A_X)).
\]
By purity,
\[
   \R^!i(A_X)\xrightarrow{\sim}(\mu_n)^{\otimes d}_Y[-2d],
\]
where $d=\codim(Y,X)$. Put $T=(\mu_n)^{\otimes d}_Y$. Since $G$ is locally constant constructible, for every geometric point $\ol y$ of $Y$ and every $p$ one has
\[
   \uExt^p(G,T)_{\ol y}\xrightarrow{\sim}\Ext^p(G_{\ol y},T_{\ol y}).
\]
But $T_{\ol y}$ is injective, being isomorphic to $A=\ZZ/n\ZZ$. Hence $\uExt^p(G,T)=0$ for $p\ne0$, and
\[
   \R\,\underline{\Hom}(G,T)\xrightarrow{\sim}\underline{\Hom}(G,T).
\]
Thus
\[
   \R\,\underline{\Hom}(F,A_X)
      \xrightarrow{\sim} \R i_*(\mathcal H)[-2d],
   \qquad \mathcal H=\underline{\Hom}(G,T).
\]
Let $Z$ be the schematic closure of $Y$ in $X$ (EGA I, 9.5.10). Then $Y=Z\cap D(f)$ (EGA $0_{\mathrm I}$, 2.1.6). Denote by $j:Y\to Z$ and $k:Z\to X$ the canonical immersions. We have $\R i_*=k_*\R j_*$, so we are reduced to computing $\R j_*(\mathcal H)$. If $\ol x$ is a geometric point of $Z$, let $\wt Z=\Spec(\cO_{Z,\ol x})$ be the strict localization, let $\wt f$ be the image of $f$ in $\cO_{Z,\ol x}$, let $\wt Y=\wt Z\times_ZY=\wt Z-V(\wt f)$, and let $\wt{\mathcal H}$ be the inverse image of $\mathcal H$ on $\wt Y$. By SGA A, Exposé VIII, 5.2,
\[
   \R^qj_*(\mathcal H)_{\ol x}
      \xrightarrow{\sim}\Hh^q(\wt Y,\wt{\mathcal H}),
\]
hence
\[
   \uExt^q(F,A_X)_{\ol x}
      \xrightarrow{\sim}\Hh^{q-2d}(\wt Y,\wt{\mathcal H}).
\]
The condition $(\mathrm{cdloc})_n$ implies that this group is zero for
\[
   q-2d>\dim\cO_{Z,\ol x}.
\]
Using
\[
   d=\codim_{\ol x}(Z,X)\le\dim\cO_{X,\ol x}
\]
(EGA IV, 5.1.3) and
\[
   \dim\cO_{Z,\ol x}
     =\dim\cO_{X,\ol x}-\codim_{\ol x}(Z,X)
\]
(EGA $0_{\mathrm{IV}}$, 16.5.12 and EGA IV, 5.1.9), one obtains
\[
   \uExt^q(F,A_X)_x=0\qquad(q>2\dim\cO_{X,x}),
\]
as required. \qedtext
\end{prooftext}

\begin{statement}{Remark 3.2.2}
The bound obtained in 3.2.1 is the best possible. Indeed, let $x$ be a point of $X$, and let $Y$ be the closed integral sub-prescheme of $X$ whose generic point is $x$. Since $X$ satisfies (reg), there exists a regular open subset $U$ of $Y$ containing $x$. By purity,
\[
   \uExt^{2d}(A_U,A_X)|U=(\mu_n)^{\otimes d}_U,
\]
where $d=\dim\cO_{X,x}=\codim(Y,X)$.
\end{statement}

\begin{statement}{Corollary 3.2.3}
Let $X$ be a regular prescheme of dimension $d$. If $X$ is smooth over a field $k$ of characteristic prime to $n$, or if $X$ is excellent of characteristic zero, then $X$ satisfies the conclusion of 3.2.1; consequently
\[
   \dim.\,q.\,inj(A_X)\le 2d.
\]
\end{statement}

\begin{prooftext}
If $X$ is excellent of characteristic zero, or if $k$ is perfect of characteristic prime to $n$ and $X$ is smooth over $k$, then the hypotheses of 3.2.1 are satisfied by 3.1.1, 3.1.3, and 3.1.4. The conclusion is still valid when $X$ is smooth over an arbitrary field $k$, as follows immediately after the base change $\Spec(k')\to\Spec(k)$, where $k'$ is a perfect closure of $k$ (SGA A, Exposé VIII, 1.1). \qedtext
\end{prooftext}

\begin{statement}{Lemma 3.2.4}
Let $X$ be a prescheme satisfying the conclusion of 3.2.1 and of dimension $\le d$. If there exists an integer $N$ such that $X$ has cohomological dimension $\le N$, then
\[
   \dim.\,inj(A_X)\le 2d+N.
\]
If, moreover, for every closed subset $Y$ of $X$, the implication
\[
   \codim(Y,X)\ge p \quad\Rightarrow\quad \cd_{\mathbf n}(Y)\le N-2p
\]
holds, then
\[
   \dim.\,inj(A_X)\le N.
\]
\end{statement}

\begin{prooftext}
The first assertion follows immediately from 1.5. For the second, one argues as in 1.5. We are reduced to proving, for $F$ constructible on $X$, the relation
\[
   \Ext^i(X;F,A_X)=0\qquad(i>N).
\]
For this, inspect the spectral sequence
\[
   E_2^{pq}=\Hh^p\bigl(X;\uExt^q(F,A_X)\bigr)
      \Rightarrow \Ext^*(X;F,A_X).
\]
By 3.2.1, the sheaves $\uExt^{2i}(F,A_X)$ and $\uExt^{2i-1}(F,A_X)$ are concentrated in codimension $\ge i$. Therefore $E_2^{p,2i}=E_2^{p,2i-1}=0$ for $p>N-2i$, by the hypothesis. Hence $E_2^{pq}=0$ for $p+q>N$, and the assertion follows. \qedtext
\end{prooftext}

\begin{statement}{Corollary 3.2.5}
If $X$ is a smooth prescheme of dimension $d$ over a separably closed field $k$ of characteristic prime to $n$, then
\[
   \dim.\,inj(A_X)=2d.
\]
\end{statement}

\begin{prooftext}
By 3.2.3 and SGA A, Exposé X, 4.2, we are in the situation of 3.2.4. The hypotheses of SGA A, Exposé X, 4.2 are satisfied by SGA A, Exposé X, 2.1 and 3.2; hence $\dim.\,inj(A_X)\le2d$. The equality follows from purity: if $x$ is a closed point of codimension $d$ in $X$, then
\[
   \Ext^{2d}(X,A_x,A_X)=\Hh^{2d}_x(A_X)
      \xrightarrow{\sim}(\mu_n)^{\otimes d}(k)\ne0.
\]
\qedtext
\end{prooftext}

\subsection*{3.3. Constructibility of the sheaves $\uExt^i(F,A_X)$}

\begin{statement}{Proposition 3.3.1}
Let $X$ be a prescheme.
\begin{enumerate}[label=\alph*)]
\item The following conditions are equivalent:
\begin{enumerate}[label=(\roman*)]
\item for $F\in\ob\D_c^-(X)$ and $G\in\ob\D_c^+(X)$, one has
\[
   \R\,\underline{\Hom}(F,G)\in\ob\D_c^+(X);
\]
\item[(i bis)] for every pair $(F,G)$ of constructible $A_X$-modules, one has $\R\,\underline{\Hom}(F,G)\in\ob\D_c^+(X)$, i.e. the sheaves $\uExt^i(F,G)$ are constructible for all $i$;
\item for every open immersion $i:U\to X$ and every constructible $A_U$-module $F$, the sheaves $\R^pi_*(F)$ are constructible for all $p$;
\item[(ii bis)] for every immersion $i:Y\to X$ and every constructible $A_Y$-module $F$, the sheaves $\R^pi_*(F)$ are constructible for all $p$;
\item[(ii ter)] for every immersion $i:Y\to X$ and every locally constant constructible $A_Y$-module $F$, the sheaves $\R^pi_*(F)$ are constructible for all $p$;
\item for every finite morphism, respectively every closed immersion, $f:Y\to X$, and every $F\in\ob\D_c^+(X)$, one has $\R^!f(F)\in\ob\D_c^+(Y)$.
\end{enumerate}
\item These conditions are satisfied in each of the following two cases:
\begin{enumerate}[label=(\roman*)]
\item $\dim(X)\le1$;
\item resolution of singularities and purity are available on $X$ in the sense of 3.1.6(b).
\end{enumerate}
\end{enumerate}
\end{statement}

\begin{prooftext}
The equivalence of (i) and (i bis) follows from the standard spectral sequence passing from sheaf-Ext to hyper-Ext.

(i bis) $\Rightarrow$ (ii). Let $\ol F$ be a constructible $A_X$-module extending $F$, for example $\ol F=i_!F$. By duality (SGA A, Exposé XVIII, 3.1.10),
\[
   \R i_*F
    =\R i_*\R\,\underline{\Hom}(A_U,i^!\ol F)
    \xrightarrow{\sim}
      \R\,\underline{\Hom}(i_!A_U,\ol F),
\]
and this proves the claim.

(ii) $\Rightarrow$ (ii bis). Factor $i$ as a closed immersion followed by an open immersion.

(ii bis) $\Rightarrow$ (i bis). The question is local, so we may suppose $X$ affine. By the usual dévissage (SGA A, Exposé IX, 2.5), it is enough to verify (ii) for $F=i_!M$, where $i:Y\to X$ is the immersion of a sub-prescheme and $M$ is a locally constant constructible $A_Y$-module. By the projection formula (SGA A, Exposé XVIII),
\[
   \R\,\underline{\Hom}(i_!M,G)
      \xrightarrow{\sim}
   \R i_*\R\,\underline{\Hom}(M,\R i^!G),
\]
and by a spectral sequence we are reduced to proving that
\[
   \R\,\underline{\Hom}(M,\R i^!G)\in\ob\D_c^+(Y).
\]
First prove $\R i^!G\in\ob\D_c^+(Y)$. After factoring $i$, we may suppose that $Y$ is closed. Let $j:U=X-Y$ be the complementary open immersion. Applying $\R\,\underline{\Hom}(\ ,G)$ to
\[
   0\longrightarrow A_{U,X}\longrightarrow A_X\longrightarrow A_Y\longrightarrow0
\]
gives a distinguished triangle. Since $\R\,\underline{\Hom}(A_X,G)=G$, and
\[
   \R\,\underline{\Hom}(A_{U,X},G)\xrightarrow{\sim}\R j_*j^!G
\]
belongs to $\D_c^+(X)$ by hypothesis, it follows that
\[
   \R\,\underline{\Hom}(A_Y,G)\xrightarrow{\sim}i_*\R i^!G
\]
belongs to $\D_c^+(X)$. Resolving $M$ locally on $Y$, for the étale topology, by free $A_Y$-modules, another dévissage gives
\[
   \R\,\underline{\Hom}(M,\R i^!G)\in\ob\D_c^+(Y).
\]
This proves the implication.

(ii bis) $\Rightarrow$ (ii ter) is trivial.

(ii ter) $\Rightarrow$ (ii). Let $i:U\to X$ be an open immersion and let $F$ be a constructible $A_U$-module. We must prove $\R i_*F\in\ob\D_c^+(X)$. The question being local, suppose $X$ noetherian and argue by noetherian induction on $X$. We may suppose $U\ne\varnothing$. Then there exists a non-empty open subset $j:V\to U$ such that $j^*F$ is locally constant (SGA A, Exposé IX, 2.4). Let $M$ be the mapping cylinder of
\[
   F\longrightarrow \R j_*j^*F.
\]
Applying $\R i_*$ gives a distinguished triangle
\[
   \R i_*F\longrightarrow \R i_*\R j_*j^*F
       \longrightarrow \R i_*M\xrightarrow{[1]}.
\]
Condition (ii ter) implies
\[
   \R i_*\R j_*j^*F\xrightarrow{\sim}\R(ij)_*j^*F\in\ob\D_c^+(X).
\]
Consequently,
\[
   \R j_*j^*F\xrightarrow{\sim}i^*\R(ij)_*j^*F\in\ob\D_c^+(U).
\]
Thus $M\in\ob\D_c^+(U)$, since $F\in\ob\D_c^+(U)$. Moreover, the cohomology of $M$ is supported on a closed subset $Z$ of $U$ distinct from $U$. Applying the induction hypothesis to $U\cap\ol Z\to\ol Z$, we deduce $\R i_*M\in\ob\D_c^+(X)$, and hence $\R i_*F\in\ob\D_c^+(X)$.

The equivalence of (i)--(ii ter) with condition (iii) is left to the reader; use 1.14.

For b)(i), normalization reduces one to the case where $X$ is regular, and one uses purity, which will be proved in 5.1. For b)(ii), see 3.1.6(b). \qedtext
\end{prooftext}

\subsection*{3.4. The biduality theorem on regular preschemes}

\begin{statement}{Theorem 3.4.1, cf. [H], V, 2.2}
Let $X$ be a regular locally noetherian prescheme of finite dimension. Suppose that $X$ is strongly desingularizable (3.1.5), satisfies (reg) (3.1.1) and $(\mathrm{cdloc})_n$ (3.1.3), and that absolute purity (3.1.4) is true on every prescheme smooth over $X$. Then $A_X$ is dualizing.
\end{statement}

\begin{prooftext}
We must prove that $A_X$ satisfies conditions (i)--(iii) of 1.6. Condition (i) follows from 3.2.1. Condition (ii) follows from 3.3.1(b). It remains to prove condition (iii). Thus, by 1.7, for every constructible $A_X$-module $F$ we must prove that the canonical homomorphism
\[
   F\longrightarrow \underline D_X\underline D_XF,
   \qquad \underline D_X=\R\,\underline{\Hom}(\ ,A_X),
\]
is an isomorphism. The question being local, suppose $X$ noetherian. By dévissage (3.1.2, iii bis), it is enough to treat sheaves
\[
   F=f_*i_!(A'_U),
\]
where $i:U\to Y$ is the inclusion of a regular open subset of an integral prescheme, $f:Y\to X$ is finite and the generic point of $Y$ is separable over its image, and $A'=\ZZ/\ell\ZZ$ for a prime divisor $\ell$ of $n$.

Because $X$ is strongly desingularizable, there is a proper birational morphism $h:Y'\to Y$, with $Y'$ regular, inducing an isomorphism from $U'=h^{-1}U$ to $U$, and such that $Y'-U'$ is a normal crossings divisor. By base change (SGA A, Exposé XVII),
\[
   i_!(A'_U)\xrightarrow{\sim}\R h_*i'_!(A'_{U'}),
\]
where $i':U'\to Y'$ is the inclusion. Hence, with $g=fh$,
\[
   f_*i_!(A'_U)
     \xrightarrow{\sim} f_*\R h_*i'_!(A'_{U'})
     \xrightarrow{\sim} \R g_*i'_!(A'_{U'}).
\]
By 1.11, it suffices to prove that the pair
\[
   \bigl(i'_!(A'_{U'}),\R^!g(A_X)\bigr)
\]
is bidualizing. Locally immersing $Y'$ in a prescheme $X'$ smooth over $X$, and applying purity on $X'$, one sees that $\R^!g(A_X)$ is locally isomorphic to $A_{Y'}$, up to a translation of degrees. We are therefore reduced to the following special case.
\end{prooftext}

\begin{statement}{Lemma 3.4.2}
Let $X$ be a locally noetherian regular prescheme on which absolute purity is true. Let $Y=Y_1\cup\cdots\cup Y_p$ be a strict normal crossings divisor on $X$, and put $U=X-Y$. Then the pair $(A'_{U,X},A_X)$ is bidualizing.
\end{statement}

\begin{prooftext}
By the exact sequence
\[
   0\longrightarrow A'_{U,X}\longrightarrow A'_X\longrightarrow A'_Y\longrightarrow0,
\]
it is equivalent to verify that $(A'_Y,A_X)$ is bidualizing. We argue by induction on $p$.

If $p=1$, let $i:Y\to X$ be the immersion. By 1.11, it is enough to verify that
\[
   (A'_Y,\R^!i(A_X))
\]
is bidualizing. By purity,
\[
   \R^!i(A_X)\xrightarrow{\sim}(\mu_n)^{\otimes1}_Y[-2].
\]
Thus the canonical homomorphism $A'_Y\to\underline D_Y\underline D_Y(A'_Y)$ is an isomorphism, as is checked trivially fibre by fibre, using Pontryagin duality for $A$-modules.

Suppose $p\ge2$, and put $Y'=Y_1\cup\cdots\cup Y_{p-1}$. We have an exact sequence
\[
   0\longrightarrow A'_{(Y-Y'),X}\longrightarrow A'_Y\longrightarrow A'_{Y'}\longrightarrow0.
\]
By the induction hypothesis, $(A'_{Y'},A_X)$ is bidualizing. By the five lemma, we are reduced to biduality for $(A'_{(Y-Y'),X},A_X)$. But again there is an exact sequence
\[
   0\longrightarrow A'_{(Y-Y'),X}\longrightarrow A'_{Y_p}\longrightarrow A'_{Y_p\cap Y'}\longrightarrow0.
\]
Since $(A'_{Y_p},A_X)$ is bidualizing, it remains to prove the same for $(A'_{Y_p\cap Y'},A_X)$. Applying 1.11 again and using purity, we are finally reduced to biduality for this pair on $Y_p$, where the induction hypothesis applies. \qedtext
\end{prooftext}

This completes the proof of Theorem 3.4.1.

\begin{statement}{Corollary 3.4.3}
Let $S$ be an excellent noetherian regular prescheme of characteristic zero. Then $A_S$ is dualizing. Moreover, if $f:X\to S$ is a morphism of finite type, then $\R^!f(A_S)$ is dualizing.
\end{statement}

\begin{prooftext}
The first assertion follows because $S$ satisfies the hypotheses of 3.4.1, as recalled in 3.1. For the second, one must first verify that $\R^!f(A_S)$ has finite quasi-injective dimension. Cover $S$ by finitely many affine open subsets $S_i$ such that $X_i=f^{-1}(S_i)$ is a finite union of affine open subsets $X_{ij}$ immersed as closed subschemes in preschemes $X'_{ij}$ smooth of finite dimension over $S_i$. From 3.2.3 and 1.6 it follows that
\[
   \dim.\,q.\,inj\bigl(\R^!f(A_S)\bigr)
       \le 2\sup \dim X'_{ij}.
\]
It remains to verify conditions (ii) and (iii) of 1.7. Since these are local on $X$ for the étale topology, we may suppose that $f$ factors as $f=f'i$, where $f':X'\to S$ is smooth of relative dimension $d$, and $i:X\to X'$ is a closed immersion. Then
\[
   \R^!f(A_S)\simeq \R^!i\,\R^!f'(A_S).
\]
But $\R^!f'(A_S)$ is locally isomorphic to $A_{X'}$, up to the translation by $2d$, hence is dualizing by the first assertion. By 1.15, $\R^!f(A_S)$ is dualizing, and in particular satisfies conditions (ii) and (iii) of 1.7. \qedtext
\end{prooftext}

\begin{statement}{Corollary 3.4.4}
Let $k$ be a field, and let $f:X\to S=\Spec(k)$ be a morphism locally of finite type, with $\dim(X)\le2$. Then $\R^!f(A_S)$ is dualizing.
\end{statement}

\begin{prooftext}
Replacing $k$ by its perfect closure, which is harmless for the étale topology (SGA A, Exposé VIII, 1.1), we may suppose $k$ perfect. Cover $X$ by affine open subsets $X_i$ of finite type over $k$. Since $\dim(X_i)\le2$, the normalization lemma makes $X_i$ finite over a scheme $X'_i$ smooth over $S$ and of dimension $\le2$. By 1.15 this reduces us to the case where $f$ is smooth, hence locally
\[
   f^!(A_S)\simeq A_X[2d].
\]
Thus it remains to show that $A_X$ is dualizing in this case; this is exactly 3.4.1, since the hypotheses are satisfied (see 3.1). \qedtext
\end{prooftext}

\begin{statement}{Remark 3.4.5}
We shall see below that if $X$ is regular of dimension $\le1$, then $A_X$ is dualizing.
\end{statement}

\section*{4. Local duality}

The purpose of this section is to present the notion of biduality developed in the preceding sections in a form closer to geometric intuition.

\subsection*{4.1}
Let $X$ be a prescheme, let $i:x\to X$ be a closed point with separably closed residue field, and let $K_X$ be a dualizing complex on $X$ (1.7). By 1.15, $K_{\{x\}}=\R^!i(K_X)$ is a dualizing complex on $x$, hence by 2.1 and 3.4.4 is isomorphic to $A[r]$ for some $r\in\ZZ$, if $A$ is local. We shall say that $K_X$ is \emph{normalized at $x$} if $r=0$ and if an isomorphism $K_{\{x\}}\xrightarrow{\sim}A$ has been chosen.

For example, if $X$ is smooth of relative dimension $d$ over the spectrum of a separably closed field, the complex $(\mu_n)_X^{\otimes d}[2d]$ is, by purity (SGA A, Exposé XVI, 3; SGA $4^{1/2}$, Cycle 2.3), canonically normalized at every closed point of $X$.

\subsection*{4.2}
Let $X$ be a prescheme, let $i:Y\to X$ be a closed sub-prescheme, and let $K_X$ be a dualizing complex on $X$. Then $K_Y=\R^!i(K_X)$ is a dualizing complex on $Y$ (1.15), and one has the complementary induction formula (1.12 b)(ii)): for $F\in\ob\D_c^b(X)$,
\[
   i^*\underline D_X(F)\xrightarrow{\sim}\underline D_Y\R^!i(F),\tag{4.2.1}
\]
where
\[
   \underline D_X=\R\,\underline{\Hom}(\ ,K_X),\qquad
   \underline D_Y=\R\,\underline{\Hom}(\ ,K_Y).
\]
In particular, take $Y=\{x\}$, where $x$ is a closed point of $X$ with separably closed residue field, and suppose $K_X$ normalized at $x$. Since $A_x$ is injective, (4.2.1) becomes
\[
   \underline D_X(F)_x\xrightarrow{\sim}\underline{\Hom}^*(\R^!i(F),A).\tag{4.2.2}
\]
Since $\R^!i(F)$ is also $\R\underline\Gamma_x(F)$ by definition, one obtains a perfect pairing in $\D(A)$,
\[
   \underline D_X(F)_x\times \R\underline\Gamma_x(F)\longrightarrow A,\tag{4.2.2'}
\]
which gives perfect pairings between finite groups
\[
   \uH^i(\underline D_XF)_x\times \uH^{-i}_x(F)\longrightarrow A.\tag{4.2.2''}
\]
For example, if $X$ is smooth of dimension $d$ over a separably closed field and $x\in X$ is closed, then, for $F\in\ob\D_c^b(X)$, there is a canonical perfect duality between finite groups
\[
   \uExt^{2d-i}\bigl(F,(\mu_n)_X^{\otimes d}\bigr)_x
      \times \uH^i_x(F)\longrightarrow A.\tag{4.2.2'''}
\]

\subsection*{4.3}
\begin{statement}{Lemma 4.3}
Let $X$ be a prescheme, let $X'$ be the strict localization of $X$ at a geometric point $\ol x$, and let $f:X'\to X$ be the canonical map. For $F\in\ob\D_c^-(X)$ and $G\in\ob\D^+(X)$, the canonical functorial homomorphism
\[
   f^*\R\,\underline{\Hom}(F,G)
      \longrightarrow
   \R\,\underline{\Hom}(f^*F,f^*G)
\]
is an isomorphism.
\end{statement}

\begin{prooftext}
As usual, one reduces to the case where $F$ and $G$ are concentrated in degree $0$. The question being local near $x$, suppose $X$ affine. The sheaf $F$ admits a left resolution (SGA A, Exposé IX, 2.7) by sheaves of the form $i_!(A_U)$, where $i:U\to X$ is étale of finite presentation. By dévissage, we are reduced to $F=i_!(A_U)$. Form the cartesian square
\[
\begin{array}{ccc}
 U'&\xrightarrow{g}&U\\
 \downarrow i'&&\downarrow i\\
 X'&\xrightarrow{f}&X.
\end{array}
\]
There are canonical isomorphisms, trivial special cases of duality (SGA A, Exposé XVIII, 3.1.10),
\[
   \R\,\underline{\Hom}(F,G)\xleftarrow{\sim}\R i_*i^*G,
   \qquad
   \R\,\underline{\Hom}(f^*F,f^*G)\xleftarrow{\sim}\R i'_*g^*i^*G.
\]
Under these isomorphisms the canonical map is identified with the base-change map
\[
   f^*\R i_*i^*G\longrightarrow \R i'_*g^*i^*G,
\]
and this is an isomorphism, as is seen immediately by passage to the limit (SGA A, Exposé VII, 5.11). \qedtext
\end{prooftext}

\begin{statement}{Lemma 4.4}
Let $X$ be a prescheme, and let $K\in\ob\D_c^+(X)$ be a complex satisfying conditions (i) and (ii) of 1.7. Then $K$ is dualizing if and only if, for every strict localization
\[
   f:\ol X(\ol x)\longrightarrow X,
\]
the complex $f^*K$ is dualizing.
\end{statement}

\begin{prooftext}
Use the fact that every constructible sheaf on a strict localization $\ol X(\ol x)$ is induced by a constructible sheaf on an étale $\ol x$-pointed scheme over $X$ (SGA A, Exposé IX, 2.7.4), and apply 4.3. \qedtext
\end{prooftext}

\subsection*{4.5}
Let $X$ be a prescheme, let $\ol x$ be a geometric point of $X$, and let
\[
   f:X'=\ol X(\ol x)\longrightarrow X
\]
be the corresponding strict localization. Let $Y$ be the integral closed sub-prescheme which is the closure of $x$ in $X$; write
\[
   g:\ol x\longrightarrow Y,
   \qquad i:Y\longrightarrow X
\]
for the canonical maps, and form the cartesian square
\[
\begin{array}{ccc}
\ol x&\xrightarrow{i'}&X'\\
\scriptstyle g\downarrow&&\downarrow\scriptstyle f\\
Y&\xrightarrow{i}&X.
\end{array}
\]
The point $\ol x$ is both the closed point of $X'$ and the strict localization of $Y$ at $\ol x$.

We denote by
\[
   \R\Gamma_{\ol x}:\D^+(X)\longrightarrow\D^+(\ol x)
\]
the functor defined by
\[
   \R\Gamma_{\ol x}=g^*\R^!i.\tag{4.5.1}
\]
Using the well-known relation between the functors $\R^!i$ and $\R j_*j^*$, where $j:U\to X$ is the open complement of $Y$, and applying passage to the limit (SGA A, Exposé VII, 5.11), one obtains a canonical isomorphism
\[
   \R\Gamma_{\ol x}\xrightarrow{\sim}\R^!i'\,f^*.\tag{4.5.2}
\]
When $x$ is a closed point with separably closed residue field, the functor $\R\Gamma_{\ol x}$ coincides with the usual functor $\R\Gamma_x$ (SGA A, Exposé V, 4.3), which justifies the notation.

Let $K_X\in\ob\D^+(X)$, and put
\[
   K_Y=\R^!i(K_X),\qquad K_{X'}=f^*K_X,
   \qquad K_{\ol x}=\R\Gamma_{\ol x}(K_X).
\]
Here $K_{\ol x}$ is not the fibre of $K_X$ at $\ol x$. Suppose $K_X$ dualizing. Then by 1.15 and 4.4, the complexes $K_{X'}$, $K_Y$, and $K_{\ol x}$ are dualizing. Hence by 2.1,
\[
   K_{\ol x}\simeq A[d]
\]
for some $d\in\ZZ$. We say that $K_X$ is \emph{normalized at} $\ol x$ if $d=0$ and an isomorphism $K_{\ol x}\simeq A$ has been chosen. When $x$ is closed with separably closed residue field, this recovers the notion introduced in 4.1.

Let $K_X$ be dualizing and put
\[
   \underline D_X=\R\,\underline{\Hom}(\ ,K_X),\qquad
   \underline D_Y=\R\,\underline{\Hom}(\ ,K_Y),
\]
etc. For $F\in\ob\D_c^b(X)$, let us compute $(\underline D_XF)_{\ol x}$. We have
\[
\begin{aligned}
(\underline D_XF)_{\ol x}
&\xrightarrow{\sim}g^*i^*\underline D_XF\\
&\xrightarrow{\sim}g^*\underline D_Y\R^!iF
   \qquad\text{by the complementary induction formula (4.2.1)}\\
&\xrightarrow{\sim}\underline D_{\ol x}g^*\R^!iF
   \qquad\text{by 4.3},
\end{aligned}
\]
that is,
\[
   (\underline D_XF)_{\ol x}
      \xrightarrow{\sim}\underline D_{\ol x}\R\Gamma_{\ol x}(F).\tag{4.5.3}
\]
Of course, one could make the same calculation through $X'$ rather than through $Y$, and would obtain the canonical isomorphism
\[
   (\underline D_XF)_{\ol x}
      \xrightarrow{\sim}\underline D_{\ol x}\R^!i' f^*F,
\]
compatible with (4.5.3) under (4.5.2).

If, moreover, $K_X$ is normalized at $\ol x$, one obtains perfect pairings generalizing (4.2.2') and (4.2.2''):
\[
   \uH^i_Y(F)_{\ol x}\times R^{-i}\Gamma_{\ol x}(F)\longrightarrow A,\tag{4.5.3'}
\]
\[
   \uH^i\bigl((\underline D_XF)_{\ol x}\bigr)\times
      \uH^{-i}_{\ol x}(F)\longrightarrow A,\tag{4.5.3''}
\]
where $\uH^i_{\ol x}=\uH^i\R\Gamma_{\ol x}$.

Conversely, let $K_X\in\ob\D^+(X)$ satisfy conditions (i) and (ii) of 1.7, and suppose that
\[
   K_{\ol x}=\R\Gamma_{\ol x}(K_X)
\]
is dualizing for every geometric point $\ol x$ of $X$. Then, for each $\ol x$, one can define a canonical map, cf. 1.12 b)(ii),
\[
   \underline D_X(F)_{\ol x}
      \longrightarrow \underline D_{\ol x}\R\Gamma_{\ol x}(F),
      \qquad F\in\ob\D_c^b(X).\tag{4.5.3 bis}
\]
If this map is an isomorphism for all $\ol x$ and all $F$, then $K_X$ is dualizing. Indeed, one must verify that
\[
   F\longrightarrow\underline D_X\underline D_X(F)
\]
is an isomorphism for all $F\in\ob\D_c^b(X)$. At an arbitrary geometric point $\ol x$, one has
\[
\begin{aligned}
F_{\ol x}
&\xrightarrow{\sim}\underline D_{\ol x}\underline D_{\ol x}(F_{\ol x})
     \qquad\text{since $K_{\ol x}$ is dualizing}\\
&\xrightarrow{\sim}\underline D_{\ol x}\R\Gamma_{\ol x}\underline D_X(F)
     \qquad\text{by 4.3 and ordinary induction}\\
&\xrightarrow{\sim}\underline D_{\ol x}\underline D_{\ol x}\bigl((F)_{\ol x}\bigr)
     \qquad\text{by hypothesis}.
\end{aligned}
\]
Modulo the verification of compatibility, this proves that $(F,K_X)$ is bidualizing.

\begin{statement}{Proposition 4.5.4}
Let $X$ be a prescheme, and let $K_X\in\ob\D^+(X)$ satisfy conditions (i) and (ii) of 1.7. Then $K_X$ is dualizing if and only if, for every geometric point $\ol x$ of $X$, the complex $K_{\ol x}$ is dualizing and the canonical map (4.5.3 bis) is an isomorphism for every $F\in\ob\D_c^b(X)$.
\end{statement}

\begin{statement}{Remark 4.5.5}
The existence of dualizing complexes must be regarded, in a sense, as an important complement to the global duality theorem. The latter says that, for a separated morphism of finite type $f:X\to Y$ with $Y$ quasi-compact, there is a canonical isomorphism
\[
   \underline D_Y\R_!f(F)\xrightarrow{\sim}\R f_*\underline D_X(F)
\]
with the notation of 1.12. To use this in the calculation of $\R_!f(F)$, when $F\in\ob\D_c^b(X)$, one needs information about $\underline D_X(F)$. If $K_X$ is dualizing, this information is supplied by (4.5.3), which allows, at least theoretically, the fibres of $\underline D_X(F)$ to be computed as cohomological invariants of a purely spatial nature.
\end{statement}

\begin{statement}{Exercise 4.5.6}
Let $X$ be a prescheme endowed with a dualizing complex (1.7). The family of functors
\[
   \uH^i_{\ol x}:\D_c^b(X)\longrightarrow A\text{-modules},
\]
where $i$ runs through $\ZZ$ and $\ol x$ runs through the geometric points of $X$, is conservative: a morphism $F\to G$ in $\D_c^b(X)$ such that all maps
\[
   \uH^i_{\ol x}(F)\longrightarrow \uH^i_{\ol x}(G)
\]
are isomorphisms is itself an isomorphism.
\end{statement}

\subsection*{4.6}
Let $X$ be a prescheme, let $j:Y\to X$ be a closed sub-prescheme, and let $i:\ol x\to Y$ be a geometric point of $Y$. Let $K_X$ be a dualizing complex on $X$, normalized at $\ol x$ in the sense of 4.5. We have seen in 4.5 that, when $Y$ is the closure of $x$, for $F\in\ob\D_c^b(X)$ there is a perfect pairing
\[
   \R^!j(F)_{\ol x}\times \underline D_X(F)_{\ol x}\longrightarrow A.
\]
It is not difficult to define an analogous pairing in the general case. Factor $i$ as
\[
   \ol x\xrightarrow{i'}\{\ol x\}\xrightarrow{i''}Y.
\]
For $F\in\ob\D_c^b(X)$, one has canonical isomorphisms
\[
\begin{aligned}
\R\Gamma_{\ol x}j^*\underline D_X(F)
&\xrightarrow{\sim}\R\Gamma_{\ol x}\underline D_Y\R^!j(F)\\
&\xrightarrow{\sim}i'^*\R^!i''\underline D_Y\R^!j(F)\\
&\xrightarrow{\sim}\underline D_{\ol x}i'^*\R^!j(F),
\end{aligned}
\]
using the complementary induction formula, the definition of $\R\Gamma_{\ol x}$, ordinary induction, and 4.3. Thus, finally, one obtains a canonical isomorphism
\[
   \R^!j(F)_{\ol x}
      \xrightarrow{\sim}
   \underline D_{\ol x}\R\Gamma_{\ol x}\,j^*\underline D_X(F),\tag{4.6.1}
\]
and hence perfect pairings
\[
   \R^!j(F)_{\ol x}\times
   \R\Gamma_{\ol x}\,j^*\underline D_X(F)\longrightarrow A,
\]
\[
   \uH^i_Y(F)_{\ol x}\times
   \uH^{-i}_{\ol x}\bigl((\underline D_XF)|Y\bigr)
      \longrightarrow A.\tag{4.6.1''}
\]

\begin{statement}{Example 4.6.2}
Let $X$ be a regular prescheme, let $\ol x$ be a geometric point of $X$, and put $d=\dim\cO_{X,\ol x}$. Let
\[
   K_X=(\mu_n^{\otimes d})_X[2d].
\]
If $X$ satisfies (reg) (3.1.1) and the purity theorem (3.1.4), then $K_X$ is canonically normalized at $\ol x$. Hence, if $K_X$ is dualizing, one has perfect pairings
\[
   \uH^i_Y(F)_{\ol x}\times
      \uH^{2d-i}_{\ol x}(F'|Y)\longrightarrow A,
\]
for every closed subset $Y$ containing $\ol x$ and every $F\in\ob\D_c^b(X)$, where
\[
   F'=\R\,\underline{\Hom}\bigl(F,(\mu_n^{\otimes d})_X\bigr).
\]
\end{statement}

We end with a strictly local variant of (4.6.1).

\subsection*{4.7}
Let $X$ be a noetherian strictly local scheme (SGA A, Exposé VIII, 4.2), with closed point $x$. Let $Y$ be a non-empty closed subset of $X$, and consider the canonical immersions
\[
   U=X-Y\xrightarrow{i}V=X-\{x\}\xrightarrow{j}X.
\]
Let $K_X$ be a dualizing complex normalized at $x$ (4.1), so that there are corresponding dualizing complexes on $U,V,Y,x$; by hypothesis
\[
   \underline D_x=\R\Hom(\ ,A),
\]
and
\[
   K_V=j^!K_X,\\ K_U=i^*K_V,
   \qquad K_Y=\R\Gamma_Y(K_X).
\]
Let $G\in\ob\D_c^b(U)$, and put $G'=\underline D_U(G)$. We shall define a natural pairing
\[
   \R\Gamma(U,G)\times\R\Gamma(V,i_*G')[-1]
       \longrightarrow A,\tag{4.7.1}
\]
or equivalently a functorial homomorphism
\[
   \R\Gamma(U,G)\overset{L}{\otimes}_A
   \R\Gamma(V,i_*G')[-1]\longrightarrow A.\tag{4.7.1'}
\]
The ring $A=\ZZ/n\ZZ$ has infinite cohomological dimension for $\ell\ge2$, so at this point one only knows how to define derived tensor products
\[
   \overset L\otimes:D^-(A)\times D^-(A)\to D^-(A)
\]
and
\[
   \overset L\otimes:D(A)\times D^b(A)_{\mathrm{torf}}\to D(A).
\]
We leave to the reader the task of imposing convenient hypotheses ensuring that the two factors in (4.7.1') belong to $D^b(A)$ or that one of them has finite Tor-dimension.\footnote{Added in 1977: this assertion is inaccurate; in fact it is easy to extend $\overset L\otimes:D^-(A)\times D^-(A)\to D^-(A)$ to $\overset L\otimes:D^-(A)\times D(A)\to D(A)$.}

First note that, by the complementary duality formula (1.12 a)(ii), there is a canonical isomorphism
\[
   i_!(G')\xrightarrow{\sim}\underline D_V(\R i_*G).\tag{4.7.2}
\]
On the other hand, there is the ordinary duality isomorphism
\[
   \R\Hom(A_V,G)\xrightarrow{\sim}\R\Hom(A_V,\R i_*G).\tag{4.7.3}
\]
By the definition of $\underline D_V$, and independently of the fact that $K_V$ is dualizing, there is a canonical map
\[
   \R i_*G\overset L\otimes \underline D_V(\R i_*G)
      \longrightarrow K_V.\tag{4.7.4}
\]
From (4.7.2), (4.7.3), and (4.7.4) one obtains a canonical map
\[
   \R\Gamma(U,G)\overset L\otimes_A \R\Gamma(V,i_*G')
      \longrightarrow \R\Hom(A_V,K_V).\tag{4.7.5}
\]
But $K_V=j^!K_X$, and there is the duality isomorphism (SGA A, 3.1.10)
\[
   \R\Hom(A_V,K_V)\xrightarrow{\sim}\R\Hom(A_{V,X},K_X).\tag{4.7.6}
\]
Finally the exact sequence
\[
   0\longrightarrow A_{V,X}\longrightarrow A_X\longrightarrow A_x\longrightarrow0\tag{4.7.7}
\]
gives a homomorphism of degree $+1$,
\[
   \R\Hom(A_{V,X},K_X)\longrightarrow
   \R\Hom(A_x,K_X)\xrightarrow{\sim}A.\tag{4.7.8}
\]
Composing (4.7.5), (4.7.6), and (4.7.8) gives the desired map (4.7.1').

We now see that the pairing (4.7.1) is perfect, i.e. that (4.7.1') identifies $\R\Gamma(V,i_*G')[-1]$ with $\underline D_x\R\Gamma(U,G)$. Suppose, as we may, that
\[
   G=(ji)^*F
\]
for some $F\in\ob\D_c^b(X)$. Put $F'=\underline D_X(F)$ and $F'_{U,X}=(ji)_!G'$. There is an exact sequence
\[
   0\longrightarrow F'_{U,X}\longrightarrow F'
      \longrightarrow F'|Y\longrightarrow0,\tag{4.7.9}
\]
giving a distinguished triangle
\[
   \R\Gamma_x(F'_{U,X})\to \R\Gamma_x(F')\to
      \R\Gamma_x(F'|Y)\xrightarrow{[1]}.
\]
By the definition of $\R\Gamma_x$,
\[
   \R\Gamma_x(F'_{U,X})=\R\Hom(A_x,F'_{U,X}).
\]
Using (4.7.7) and the fact that $(F'_{U,X})_x=0$, one gets a canonical isomorphism
\[
   \R\Hom(A_{V,X},F'_{U,X})[1]
      \xrightarrow{\sim}\R\Gamma_x(F'_{U,X}).
\]
But
\[
\begin{aligned}
\R\Hom(A_{V,X},F'_{U,X})
&=\R\Hom(j_!A_V,j_!i_!G')\\
&\xrightarrow{\sim}\R\Hom(A_V,i_!G')
     \qquad\text{by duality (SGA A, 3.1.10)},
\end{aligned}
\]
and hence there is a canonical isomorphism
\[
   \R\Gamma(V,i_!G')[-1]
      \xrightarrow{\sim}\R\Gamma_x(F'_{U,X}).\tag{4.7.10}
\]
On the other hand, the exact sequence
\[
   0\longrightarrow A_{U,X}\longrightarrow A_X\longrightarrow A_Y\longrightarrow0\tag{4.7.11}
\]
gives a distinguished triangle
\[
   \R\Gamma_Y(F)\to \R\Gamma(X,F)\to
   \R\Gamma(U,F|U)\xrightarrow{[1]}.
\]
By (4.7.1) and (4.7.10), one has a natural pairing
\[
   \R\Gamma(U,F|U)\times \R\Gamma_x(F'_{U,X})\longrightarrow A.\tag{4.7.12}
\]
By (4.6.1), one has a perfect pairing
\[
   \R\Gamma_Y(F)\times \R\Gamma_x(F'|Y)\longrightarrow A.\tag{4.7.13}
\]
Finally, by ordinary induction (1.12 b)(i)), one has a perfect pairing
\[
   \R\Gamma(X,F)\times \R\Gamma_x(F')\longrightarrow A.\tag{4.7.14}
\]
The reader patient enough to orient himself through this maze of ``canonical arrows'' will verify that the pairings (4.7.12), (4.7.13), and (4.7.14) are compatible with the arrows of the two triangles, and hence that (4.7.12), and therefore (4.7.1), are perfect.

The pairing of the two triangles gives two exact sequences transposed to one another:
\[
   \cdots\to \uH^i_Y(X,F)\to \uH^i(X,F)\to \uH^i(U,F)\to\cdots
\]
and
\[
   \cdots\leftarrow \uH^{-i}_x(Y,F'|Y)
      \leftarrow \uH^{-i}_x(X,F')
      \leftarrow \uH^{-i-1}_x(V,F'_{U,V})\leftarrow\cdots.
\]

\begin{statement}{Proposition 4.7.16}
Let $X$ be a noetherian strictly local scheme, with closed point $x$. Let $Y$ be a non-empty closed subset of $X$, put $U=X-Y$ and $V=X-x$, and let $i:U\to V$ and $j:V\to X$ be the canonical inclusions. Let $K_X$ be a dualizing complex on $X$ normalized at $x$, hence giving associated dualizing complexes on $V$, $U$, and $Y$.

For $G\in\ob\D_c^b(U)$, there is a perfect pairing
\[
   \R\Gamma(U,G)\times\R\Gamma(V,i_*G')[-1]	o A,
\]
where $G'=\underline D_U(G)$. If $G=F|U$ for $F\in\ob\D_c^b(X)$, then there is a canonical isomorphism
\[
   \R\Gamma(V,i_*G')[-1]\xrightarrow{\sim}\R\Gamma_x(F'_{U,X}),
\]
where $F'=\underline D_X(F)$ and $F'_{U,X}=(ji)_!(F'|U)=(ji)_!(G')$, and the exact localization triangles are paired perfectly as described above.
\end{statement}

\begin{statement}{Remark 4.7.17}
If, in the preceding proposition, $X$ is a strictly local regular prescheme of dimension $d$, if $Y=\{x\}$, and if $K_X=(\mu_n^{\otimes d})_X[2d]$, then $K_U=(\mu_n^{\otimes d})_U[2d]$. Assuming $K_X$ dualizing, local duality on $X$ implies, for every locally constant constructible $A$-module $G$, perfect pairings
\[
   H^i(U,G)\times H^{2d-1-i}(U,G^\circ)\longrightarrow A,
\]
where $G^\circ=\underline{\Hom}(G,(\mu_n^{\otimes d})_U)$. Formally, this corresponds to Poincaré duality on a scheme of dimension $2d-1$.
\end{statement}

\section*{5. Local duality on curves}

\begin{statement}{Theorem 5.1}
Let $X$ be a noetherian regular prescheme of dimension $1$. Suppose $n$ is prime to the residual characteristics of $X$. Then the purity theorem is true on $X$ (3.1.4), and the complex $A_X$ is dualizing (1.7).
\end{statement}

\begin{prooftext}
Assume first that purity has been proved. Since $X$ obviously satisfies (reg), and since $(\mathrm{cdloc})$ holds by SGA A, Exposé X, 2.2, 3.2.1 gives
\[
   \dim.\,q.\,inj(A_X)=2.
\]
Moreover condition (ii) of 1.7 holds by 3.3.1 b)(i), whose proof uses purity. Hence, by 4.4, we are reduced to proving that $A_X$ is dualizing when $X$ is strictly local. For purity, also a local question near a closed point, we may likewise suppose $X$ strictly local. Thus it remains to prove the theorem for strictly local $X$.

Assume then that $X$ is strictly local. Let $x$ be its closed point, let $p$ be the residual characteristic, let $\eta$ be its generic point, and put
\[
   U=X-x=\Spec(k(\eta)).
\]
Let $G=\Gal(\overline{k(\eta)}/k(\eta))$ be the fundamental group of $U$. One knows (cf. Chapter IV, §2, Exercise 2) that there is a canonical exact sequence
\[
   1\longrightarrow P\longrightarrow G\longrightarrow \ZZ_p(1)\longrightarrow1,\tag{5.1.1}
\]
where
\[
   \ZZ_p(1)=\varprojlim \mu_{p^r}(\overline{k(\eta)})
\]
and $P$ is a pro-$p$-group. Of course, the quotient is isomorphic, but not canonically, to the product of the $\ZZ_\ell(1)$ for $\ell\ne p$.

On the other hand, the étale cohomology of $U$ is interpreted as the Galois cohomology of $G$ (SGA A, Exposé VIII, 2); sheaves, respectively constructible sheaves, of $A$-modules correspond to $A$-modules, respectively finite-type $A$-modules, on which $G$ acts continuously.

The groups $H^i_X(A_X)$ are determined by the exact sequence
\[
0\to H^0_X(A_X)\to H^0(A_X)\to H^0(U,A_U)
\to H^1_X(A_X)\to0,\tag{5.1.2}
\]
and by the isomorphisms
\[
   H^i_X(A_X)\xrightarrow{\sim}H^{i-1}(U,A)
      \qquad(i\ge2).\tag{5.1.3}
\]
Trivially, $H^0(X,A_X)=H^0(U,A_U)=A$, so $H^0_X(A_X)=H^1_X(A_X)=0$.

For a Galois $G$-$A$-module $M$, the Hochschild--Serre spectral sequence gives
\[
   E_2^{pq}=H^p(H,H^q(P,M))\Rightarrow H^{p+q}(G,M).
\]
Since $n$ is prime to $p$ and $P$ is a pro-$p$-group, $H^q(P,M)=0$ for $q\ne0$, hence
\[
   H^i(G,M)\xrightarrow{\sim}H^i(H,M^P).\tag{5.1.4}
\]
An analogous argument identifies $H^i(H,N)$ with $H^i(\ZZ_p(1),N)$ when the complementary factors act trivially on $N$. Thus, with the preceding convention, there is a canonical isomorphism
\[
   H^i(G,M)\xrightarrow{\sim}H^i(\ZZ_p(1),M^P).\tag{5.1.5}
\]
The Galois cohomology of $\ZZ_p(1)$ is well known (CG, Chapter I, p. 31): for a $\ZZ_p(1)$-$A$-module $N$,
\[
H^i(\ZZ_p(1),N)=
\begin{cases}
N^{\ZZ_p(1)},& i=0,\\
N_{\ZZ_p(1)},& i=1,\\
0,& i\ge2.
\end{cases}\tag{5.1.6}
\]
It follows from (5.1.3), (5.1.5), and (5.1.6) that
\[
   H^i(U,A_X)=0\qquad(i>2),
\]
and $H^2(U,A_X)\simeq A$. This last isomorphism is not canonical, since it depends on a choice of an identification of $\ZZ(1)$ with $\ZZ/n\ZZ$. But using the Kummer exact sequence, one obtains (SGA A, Exposé XIX, 1.3) a canonical isomorphism, given by the local fundamental class,
\[
   H^2(U,\mu_n)\xrightarrow{\sim}A.\tag{5.1.7}
\]
Purity is therefore proved.

It remains to show that $A_X$ is dualizing. By 4.7.17, it suffices to show that, for every locally constant constructible $A_X$-module $M$, the pairing
\[
   \R\Gamma(U,M)\times \R\Gamma(U,\underline{\Hom}(M,A_X))
      \longrightarrow \R\Gamma(U,A_X)\xrightarrow{\sim}A[-2]\tag{5.1.8}
\]
is a perfect duality. We interpret this pairing in terms of Galois cohomology, using the following elementary lemma.
\end{prooftext}

\begin{statement}{Lemma 5.1.9}
Let $D=\Hom(\ ,A)$ be the Pontryagin-duality functor on the category of finite-type $A$-modules. Let $M$ be a finite-type $A$-module on which a group $G$ acts $A$-linearly.
\begin{enumerate}[label=(\roman*)]
\item There is a canonical isomorphism $M^G\xrightarrow{\sim}D(M_G)$.
\item If $G$ is a profinite group acting continuously on $M$ and if $G$ has pro-order prime to $n$, then the canonical map $M^G\to M_G$ is an isomorphism.
\end{enumerate}
\end{statement}

\begin{prooftext}
By definition,
\[
   M^G=\varprojlim_{g\in G}\Ker(M\xrightarrow{g-1}M),
\]
hence
\[
   D(M^G)\xrightarrow{\sim}\varinjlim_{g\in G}
        \Coker(DM\xrightarrow{g-1}DM)
      \xrightarrow{\sim}(DM)_G,
\]
which proves (i).

For (ii), let $M'$ be defined by
\[
   0\longrightarrow M'\longrightarrow M\longrightarrow M_G\longrightarrow0.
\]
Since $G$ has pro-order prime to $n$, one has $H^0(G,M')=0$. The cohomology exact sequence implies that $M^G\to M_G$ is an epimorphism. Applying $D$ and using (i), one sees that it is also a monomorphism. \qedtext
\end{prooftext}

Applying 5.1.9 to $M$ and to $D(M)$, one obtains a canonical identification
\[
   \Hom(M,A)^G\simeq \Hom(M^G,A).
\]
After choosing an isomorphism $\ZZ(1)\simeq\ZZ/n\ZZ$, the pairing (5.1.8) can be written, using (5.1.5), as
\[
   H^i(G,M)\times H^{1-i}\bigl(G,D(M^P)\bigr)
      \longrightarrow H^1(\ZZ_p(1),A)\simeq A.\tag{5.1.8'}
\]
It follows immediately from (5.1.6) and 5.1.9(i) that this is a perfect pairing. This completes the proof of 5.1.

\section*{Bibliography for Exposé I}
\begin{enumerate}[label={[\arabic*]}]
\item H. Hironaka, \emph{Resolution of singularities of an algebraic variety over a field of characteristic zero}, Ann. of Math. 79 (1964), 109.
\item S. Abhyankar, \emph{Local uniformization on algebraic surfaces over ground fields of characteristic $p\ne0$}, Ann. of Math. 63 (1956), 491--526; and \emph{Resolution of singularities of arithmetical surfaces}, in \emph{Arithmetical Algebraic Geometry}, Harper \& Row, New York, 1965.
\item R. Hartshorne, \emph{Residues and Duality}, Lecture Notes in Mathematics 20, Springer, 1966.
\item[SGA A] Séminaire de Géométrie Algébrique de l'IHES, \emph{Cohomologie étale des schémas}, by M. Artin and A. Grothendieck, 1963--64.
\item[CL] J.-P. Serre, \emph{Local Fields}, Hermann, 1962.
\item[CG] J.-P. Serre, \emph{Galois Cohomology}, Lecture Notes in Mathematics 5, Springer, 1964.
\end{enumerate}

\section*{Appendix to Exposé I}
\emph{by L. Illusie}

The preceding theory extends without difficulty to base rings more general than the constant ring $\ZZ/n\ZZ$. We briefly indicate how the theory is written for locally constant base rings, annihilated by an integer $n>0$, whose geometric fibres are noetherian rings. This is only a provisional framework: sooner or later one will have to globalize, taking as ``coefficients'' on a prescheme $X$ relative preschemes on the site $(X_{\mathrm{\acute et}},\ZZ/n\ZZ)$.

\subsection*{1. Constructible sheaves}

\begin{statement}{1.1}
Let $X$ be a prescheme, and let $A_X$ be a sheaf of rings on $X$ for the étale topology, locally constant, with geometric fibres noetherian on the left. Write $\underline{\mathrm{Mod}}_{A_X}$ for the category of left $A_X$-modules, and $D(X)$, or $D(A_X,X)$ if ambiguity is possible, for the corresponding derived category. The notion of a constructible left $A_X$-module is defined as in SGA A, Exposé IX, 2.3, replacing the expression ``locally constant with value of finite presentation'' by ``of finite presentation as an $A_X$-module.''

Let $D_c(X)$ be the full subcategory of $D(X)$ consisting of complexes $E$ such that $H^i(E)$ is a constructible $A_X$-module for every $i$. It is a triangulated subcategory of $D(X)$.
\end{statement}

\begin{statement}{1.2}
Let $E\in\ob D(X)$. One says, cf. Exposé II, 3.11, that $E$ is pseudo-coherent if, for every $i$, the sheaf $H^i(E)$ is an $A_X$-module of finite presentation, or equivalently locally constant constructible. We denote by $D_{\mathrm{coh}}(X)$, or $D_{\mathrm{coh}}(A_X,X)$, the full subcategory of $D(X)$ consisting of pseudo-coherent objects. It follows from SGA A, Exposé IX, 2.1 that this is a triangulated subcategory of $D_c(X)$.
\end{statement}

\begin{statement}{Proposition 1.3, dévissage lemma}
Let $X$ be a quasi-compact and quasi-separated, respectively noetherian, prescheme, and let $\cC$ be a strictly full subcategory of the category $\operatorname{Cons}(X)$ of constructible left $A_X$-modules. Suppose $\cC$ is stable under extensions. The following conditions are equivalent:
\begin{enumerate}[label=(\roman*)]
\item $\cC=\operatorname{Cons}(X)$;
\item $\cC$ contains sheaves of the form $i_*F$, where $i:Y\to X$ is a sub-prescheme of finite presentation, respectively an integral sub-prescheme, and $F$ is locally constant constructible on $Y$.
\end{enumerate}
If $A_X$ is constant with value $A$, a left noetherian ring annihilated by $n>0$, these conditions are equivalent to:
\begin{enumerate}[label=(\roman*)]
\item[(iii)] $\cC$ is stable under direct factors and contains sheaves of the form $f_*j_!(M_Y)$, where $j:U\to Y$ is an open immersion, respectively an open subset of an integral prescheme, $f:Y\to X$ is finite, respectively finite with the generic point of $Y$ separable over its image, and $M_Y$ is the constant left $A_Y$-module associated with a cyclic $A$-module annihilated by a prime divisor of $n$.
\end{enumerate}
Finally, if $X$ is noetherian and satisfies (reg) (I, 3.1.1), condition (ii), respectively (iii), is equivalent to (i) if $Y$, respectively $U$, is moreover assumed regular.
\end{statement}

\begin{prooftext}
The equivalence of (i) and (ii) is proved as in the case where $A_X$ is constant; see SGA A, Exposé IX, 2.4 and 2.5. Suppose $A_X$ constant with value $A$, annihilated by $n$, and prove that (iii) implies (ii). Let $E\in\ob\operatorname{Cons}(X)$; we show that $E\in\ob\cC$. By the equivalence (i) $\Leftrightarrow$ (ii), one may suppose $E=j_!F$, where $j:Z\to X$ is the immersion of a sub-prescheme of finite presentation, respectively an integral sub-prescheme, and $F$ is a locally constant constructible $A_Z$-module. After a further dévissage, we may suppose $F$ trivialized by an étale covering $g:Z'\to Z$.

Then one has an exact sequence
\[
   0\to F\to g_*g^*F\to\operatorname{coker}\to0,
\]
and one is reduced to proving that $g_*g^*F$ and the cokernel belong to $\ob\cC$. This follows at once from (iii) and the induction hypothesis on the dimension of $Z$, taking into account that $g_*g^*F$ is locally constant and that the cokernel is supported on the non-smooth locus of $g$. \qedtext
\end{prooftext}

\begin{statement}{Corollary 1.4, dévissage for pseudo-coherent objects}
Let $X$ be quasi-compact and quasi-separated, respectively noetherian, and let $\cC$ be a strictly full triangulated subcategory of $D_{\mathrm{coh}}(X)$. The following conditions are equivalent:
\begin{enumerate}[label=(\roman*)]
\item $\cC=D_{\mathrm{coh}}(X)$;
\item $\cC$ contains objects of the form $i_*F$, where $i:Y\to X$ is a sub-prescheme of finite presentation, respectively an integral sub-prescheme, and $F\in\ob D_{\mathrm{coh}}(Y)$.
\end{enumerate}
If $A_X$ is constant with value $A$ annihilated by $n$, these are equivalent to the condition that $\cC$ is stable under direct factors and contains the objects $f_*j_!(M_Y)$ with $j$, $f$, and $M_Y$ as in 1.3(iii).
\end{statement}

\begin{statement}{1.5}
With the notation of 1.1, a complex $K\in\ob D(X)$ is called \emph{dualizing} if it satisfies conditions (i), (ii), (iii) of I, 1.7. Explicitly:
\begin{enumerate}[label=(\roman*)]
\item $K\in\ob D_c^b(X)$;
\item for every $F\in\ob D_c^b(X)$, $\R\,\underline{\Hom}(F,K)\in\ob D_c^b(X)$;
\item the canonical map
\[
   F\longrightarrow \R\,\underline{\Hom}(\R\,\underline{\Hom}(F,K),K)
\]
is an isomorphism for $F\in\ob D_c^b(X)$.
\end{enumerate}
\end{statement}

\begin{statement}{Corollary 1.6, dévissage for dualizing complexes}
Let $X$ be a noetherian prescheme. Let $K\in\ob D_c^b(X)$ satisfy conditions (i) and (ii) of 1.5. Then $K$ is dualizing if and only if, for every geometric point $\ol x$ of $X$, the complex $\R\Gamma_{\ol x}(K)$ is dualizing.
\end{statement}

\begin{prooftext}
This follows immediately from 1.4, dévissage for biduality, and from I, 4.3. \qedtext
\end{prooftext}

\begin{statement}{1.7}
Keep the preceding notation, but suppose $A_X$ commutative. The notions of bidualizing pair and dualizing complex are defined as in I, 1.7. The statements I, 1.9, 1.11, and 1.14 remain valid without modification. The statements I, 1.12, 1.13, and 1.15 are also valid, provided $A$ is annihilated by an integer $n>0$ prime to the residual characteristics of $X$. Finally, the uniqueness theorem I, 2.1 is valid provided one assumes that $\Spec(A_{\ol x})$ is connected for every geometric point $\ol x$ of $X$.
\end{statement}

\subsection*{2. Dualizing complexes: definition and formal properties}

Retaining the preceding notation but assuming $A_X$ commutative, the definitions and formal statements of the preceding paragraph remain in force. In particular, the formalism of dualizing complexes is unchanged, subject to the same restriction that $A$ be killed by an integer $n>0$ invertible on the base and to the connectedness hypothesis on $\Spec(A_{\ol x})$ for uniqueness.

\subsection*{3. Absolute purity and quasi-injective dimension}

We return to the situation of absolute cohomological purity (I, 3.1.4). If $X$ is regular locally noetherian, if $n>0$ is prime to the residual characteristics of $X$, and if $i:Y\to X$ is a regular sub-prescheme of codimension $d>0$ at every point, the assertion that $(Y,X)$ satisfies absolute purity is the assertion that the local fundamental class homomorphism of SGA $4^{1/2}$, Cycle 2.2, is an isomorphism.

One has the obvious inequality
\[
   \dim\operatorname{top}_x(K)\le \dim.\,q.\,inj_x(K),
\]
and says that examples where equality holds will be seen below. It also restates, for this more general coefficient sheaf $A_X$, the formal definitions of dualizing complexes and the hypotheses under which the statements of Exposé I carry over.

\begin{statement}{Proposition 5.1, Appendix: bound for quasi-injective dimension}
Let $X$ be a regular locally noetherian prescheme, equipped with a locally constant sheaf of rings $A_X$, annihilated by an integer $n>0$ prime to the residual characteristics of $X$, with geometric fibres left noetherian. Suppose that $X$ satisfies (reg) (I, 3.1.1) and $(\mathrm{cdloc})_n$ (I, 3.1.3), and that the strong form of the purity theorem is true on $X$ (4.4). These hypotheses are all satisfied if $X$ is smooth over a perfect field, or if $X$ is excellent of characteristic zero.

Let $F\in\ob D_c^b(A_X)$, with the sheaves $\uH^i(F)$ locally constant. Let $x$ be a point of $X$ such that $F_{\ol x}\ne0$ and
\[
   \dim.\,inj_{A_{\ol x}}(F_{\ol x})<+\infty.
\]
Then
\[
   \dim.\,q.\,inj_x(F)
      \le 2\dim\cO_{X,x}+\dim.\,inj_{A_{\ol x}}(F_{\ol x}).
\]
\end{statement}

\begin{prooftext}
One repeats, step by step, the proof of I, 3.2.1. Let the right-hand side be $N$; the first aim is to prove
\[
   \uExt^i(E,F)_{\ol x}=0
      \qquad(i>N)
\]
for every constructible $A_X$-module $E$. Since the question is local, dévissage (1.3) reduces the verification to $E=i_!M$, where $i:Y\to X$ is an integral regular sub-prescheme, closed in an open $D(f)$, of codimension $d$ at every point of $D(f)$, and where $M$ is a locally constant constructible $A_Y$-module.
\end{prooftext}

\section*{Exposé III. Lefschetz Formula}

\begin{center}
\textbf{Lefschetz Formula}\\
\emph{by A. Grothendieck, written up by L. Illusie}
\end{center}

In this exposé we prove the Lefschetz formula in étale cohomology for correspondences between complexes of sheaves, in Verdier's formulation. In an earlier version, distributed by the IHES about ten years before, the formula was stated under resolution hypotheses, but the required compatibilities were not proved. Deligne's finiteness theorems (SGA $4^{1/2}$, Finiteness) made it possible to dispense with these hypotheses, provided one works over schemes separated and of finite type over a field. As for the compatibilities, the writer has checked them, painfully.

Let us say at once that the Lefschetz formula considered here, like the trace formula for constant coefficients (SGA $4^{1/2}$, Cycle) which it generalizes, is essentially tautological: it expresses certain compatibilities in the formalism of duality. Roughly speaking, it says that, under suitable conditions, the trace of an endomorphism of a hypercohomology complex $\R\Gamma(X,L)$, where $L$ is a complex of sheaves on $X$, may be computed as the ``integral'' of certain cohomology classes on $X\times X$. In practice, the endomorphisms considered arise from correspondences $C\subset X\times X$, and the classes recovered on the product are supported on $C\cap\Delta$, where $\Delta$ is the diagonal. These classes are local in nature near the fixed points, which in principle makes their computation easier. Precise statements are given below, notably in 4.2 and 4.7. In fact, the computation of these local classes in terms of more understandable invariants, without which the formula would be of little use, has so far been carried out only in special situations; see 4.10, 4.11, and the following exposé, where one treats in particular the Frobenius correspondence on a curve.

\subsection*{1. Reminders and notation}

\begin{statement}{1.1}
Let $S$ be a scheme and let $f:X\to S$ be a separated morphism of finite type. Suppose given on $S$ a sheaf of rings $A_S$ for the étale topology; put $A_X=f^*A_S$, and similarly $A_Y=g^*A_S$ for an $S$-scheme $g:Y\to S$. We suppose $A_S$ annihilated by an integer $n>0$ invertible on $S$.
\end{statement}

\begin{statement}{1.2}
We recall briefly the principal statements of the duality formalism for coefficients $A_X$ (SGA A, or Exposé I in the case of constant coefficients). For a separated morphism of finite type $f:X\to Y$, there is a functor $\R^!f$, compatible with composition, and an adjunction with $\R f_!$. It is characterized, in the situations considered below, by the duality isomorphisms
\[
   \R\Hom_Y(\R f_!E,F)\simeq \R\Hom_X(E,\R^!fF),
\]
together with their sheafified analogues.
\end{statement}

\begin{statement}{1.3. Transitivity}
For a composite $f=f_1\circ f_2$, one has canonical transitivity isomorphisms
\[
   \R^!f\simeq \R^!f_2\,\R^!f_1,
\]
compatible in the usual way for triple composites.
\end{statement}

\begin{statement}{1.4. Duality for a separated finite-type morphism}
For $f:X\to Y$ separated and of finite type, and for complexes of finite tor-dimension in the constructible derived category, one uses the usual duality isomorphisms relating $\R f_!$, $\R f_*$, and $\R^!f$. These isomorphisms are the formal tools from which the Lefschetz construction below will be assembled.
\end{statement}

\begin{statement}{1.5. Smooth case}
If $f$ is smooth, purely of relative dimension $d$, then there is a canonical isomorphism, the Poincaré duality formula,
\[
   \R^!f(A_Y)\simeq A_X(d)[2d],
\]
where $A_X(d)$ denotes the usual Tate twist. More generally this identifies the relative dualizing object in the smooth case.
\end{statement}

\begin{statement}{1.6. External tensor products}
For two separated finite-type morphisms $f_i:X_i\to S$ ($i=1,2$), there is an external tensor product on complexes and on cohomological correspondences. The compatibilities of this product with $\R f_!$, $\R f_*$, and duality are those needed for taking products of trace classes.
\end{statement}

\subsection*{2. Lefschetz formula for a proper morphism}

\begin{statement}{2.1}
Let $S$ be a scheme, and let $X$ and $Y$ be two separated finite-type $S$-schemes. With coefficients as in 1.1, let
\[
   E\in\ob D_{\mathrm{ctf}}(X,A_X),
   \qquad F\in\ob D_{\mathrm{ctf}}(Y,A_Y).
\]
One defines a canonical morphism, functorial in $E$ and $F$, which is the basic cup-product pairing underlying the trace formalism.
\end{statement}

\begin{statement}{2.2}
For every $S$-morphism $u:E\to F$, lying over an $S$-morphism $g:X\to Y$, one defines the trace of $u$ as an element of a group of the form
\[
   H^0\bigl(S,\R\Hom(\R f_{1!}E,\R f_{2*}F)\bigr),
\]
where the $f_i$ are the relevant structure morphisms. This description is intentionally functorial, since the formula will later be applied to morphisms carried by correspondences.
\end{statement}

\begin{statement}{2.3}
Suppose now $X=Y$, and let $g:X\to X$ be an $S$-endomorphism. For complexes $E,F\in\ob D_{\mathrm{ctf}}(X,A_X)$ and a morphism
\[
   u:g^*E\longrightarrow F,
\]
one obtains a trace element attached to $u$. Its local form is the class that appears in the Lefschetz formula.
\end{statement}

\begin{statement}{2.4}
Let, for $i=1,2$, $f_i:X_i\to Y_i$ be an $S$-morphism, with complexes $E_i,F_i\in\ob D^b_{\mathrm{ctf}}(X_i)$ and $P_i,Q_i\in\ob D^b_{\mathrm{ctf}}(Y_i)$. The formalism of external products gives the compatibility of traces with products. This is the categorical expression of the multiplicativity of Lefschetz classes.
\end{statement}

\begin{statement}{2.5}
The principal proposition of this section states the compatibility of the above trace construction with composition and exterior tensor product. In the special case $Y_1=Y_2=S$ and $X_1=X_2=X$, for an $S$-endomorphism $g$ of $X$, it produces the Lefschetz class associated with a pair of mutually transposed cohomological correspondences.
\end{statement}

\begin{statement}{2.6}
The proof of 2.5 occupies the rest of section 2. It follows from the compatibility of external products with the duality isomorphisms that if $w_i:L_i\to M_i$ are morphisms of type 1 over $f_i$, then the exterior tensor product
\[
   w=w_1\overset L\otimes_S w_2
\]
makes the required diagram commute. The same verification may be written with the notational variants $u_i$ and $w_i$.
\end{statement}

\begin{statement}{2.7--2.9}
When $S=\Spec(k)$, with $k$ separably closed, one has $H^0(S,A_S)=A$, and the Lefschetz formula takes its classical numerical form. In the case where $E$ is a locally constant sheaf of free modules of finite rank and $g$ is the identity, it reduces to the Euler--Poincaré formula: the trace is the alternating sum of traces on cohomology.
\end{statement}

\section*{Exposé III continued. Cohomological correspondences and trace formulae}

\subsection*{3. Cohomological correspondences}

\subsubsection*{3.1. Generalities}
We resume the notation of 2.2. Let $(S,\cO_S)$ be a scheme, let $X$ and $Y$ be $S$-schemes, and let
\[
   L\in\ob D^{-}(X), \qquad M\in\ob D^{-}(Y).
\]
For $i\in \ZZ$ put
\begin{equation}
   \Hom^i_S(L,M)=\Ext^i_{\cO_S}(L,M),
   \tag{3.1.1}
\end{equation}
where the $\Ext^i_{\cO_S}$ are the derived functors of $\Hom_{\cO_S}$. If $S=\Spec(k)$, with $k$ a field, and if $L$ and $M$ are coherent sheaves, this is the ordinary group of extensions over $S$.

The Yoneda product gives pairings
\begin{equation}
   \Hom^i_S(L,M)\otimes \Hom^j_S(M,N)
      \longrightarrow \Hom^{i+j}_S(L,N),
   \tag{3.1.2}
\end{equation}
and hence, in particular, canonical homomorphisms
\begin{equation}
   \Hom^i_S(L,M)	o
   \Hom_{\cO_S}\bigl(\Ext^j_S(M,N),\Ext^{i+j}_S(L,N)\bigr).
   \tag{3.1.3}
\end{equation}
This is the convention by which the action of a correspondence on cohomology is read below.

\subsubsection*{3.2. Correspondences with support}
Suppose given a commutative diagram
\[
\begin{array}{ccc}
C & \xrightarrow{p} & X \\
\downarrow q & & \downarrow f \\
Y & \xrightarrow{g} & S .
\end{array}
\]
For $L\in\ob D^{-}(X)$ and $M\in\ob D^{-}(Y)$ one has the formal morphism, as in 2.2.2,
\begin{equation}
   p^*L\otimes^L_{\cO_C}q^*M
      \longrightarrow p^*L\otimes^L_{\cO_C}q^*M,
   \tag{3.2.1}
\end{equation}
which, by adjunction, yields
\begin{equation}
   \Hom_C(p^*L,q^*M)	o \Hom_X(L,\R p_*q^*M).
   \tag{3.2.2}
\end{equation}

\begin{statement}{Definition 3.3}
A \emph{cohomological correspondence} from $L$ to $M$, with support in $C$, is an element
\[
   u\in \Hom_C(p^*L,q^*M)
       =H^0\bigl(C,\underline{\R\Hom}(p^*L,q^*M)\bigr).
\]
\end{statement}

\begin{statement}{3.3.1. Composition}
Let $u\in \Hom_C(p^*L,q^*M)$ and $v\in\Hom_D(p^*M,q^*N)$ be correspondences for compatible supports. Their composite $vu$ is obtained by pulling $u$ to the intermediate correspondence, composing with $v$, and applying the relevant adjunction. This convention agrees with the Yoneda product (3.1.2) on the cohomology groups on which the correspondences act.
\end{statement}

\begin{statement}{3.4. Base change}
For a cartesian square of supports
\[
\begin{array}{ccc}
C' & \longrightarrow & C \\
\downarrow & & \downarrow \\
X'\times_S Y' & \longrightarrow & X\times_S Y,
\end{array}
\]
there is a canonical base-change morphism for cohomological correspondences. The construction is the evident one: pull back the two factors, use the functorial base-change morphism for the derived inverse images, and transport the result through adjunction.
\end{statement}

\begin{statement}{Proposition 3.5}
The base-change diagrams just described are commutative. Equivalently, the formation of the morphism induced by a cohomological correspondence is compatible with passage to a cartesian base change.
\end{statement}

\begin{statement}{3.6. Proper smooth schemes over an algebraically closed field}
Assume $S=\Spec(k)$ with $k$ algebraically closed. Let $X$ and $Y$ be proper smooth $k$-schemes, and let
$L\in\ob D^b_{\mathrm{coh}}(X)$, $M\in\ob D^b_{\mathrm{coh}}(Y)$. The duality isomorphisms give, in this notation,
\begin{equation}
   \Hom(L,\cO_X)\xrightarrow{\sim} H^*(X,L)^{\vee}.
   \tag{3.6.1}
\end{equation}
Consequently a cohomological correspondence determines, by functoriality, a homomorphism on cohomology.
\end{statement}

\begin{statement}{3.6.1. Cycle class case}
If $L=M=\cO_X$, $X=Y$, and $C\subset X\times X$ is a cycle of dimension $d=\dim X$, the associated cohomological correspondence is represented by the cohomology class
\[
   [C]\in H^{2d}(X\times X)(d),
\]
and the operation induced on cohomology is cup-product with $[C]$, followed by the appropriate projection.
\end{statement}

\subsubsection*{Direct image of a correspondence}

\begin{statement}{3.7}
Let
\[
\begin{array}{ccccc}
C & & & & \\
\downarrow c & \searrow & & & \\
X\times_S Y & \xrightarrow{p_1} & X & \xrightarrow{f} & S \\
& \xrightarrow{p_2} & Y & \xrightarrow{g} & S
\end{array}
\]
be the corresponding diagram. If
\[
   u\in \Hom_C(p_1^*L,p_2^!M)
\]
is a cohomological correspondence, its \emph{direct image} is the composite
\begin{equation}
   L\xrightarrow{\mathrm{adj}} \R f_*f^*L
    \xrightarrow{\R f_*(u)} \R f_*p_{2*}p_2^!M
    \xrightarrow{\mathrm{trace}} M .
   \tag{3.7.1}
\end{equation}
\end{statement}

\begin{statement}{3.7.1}
There is a canonical morphism
\begin{equation}
   g^!f_*K\longrightarrow f'_*g'^!K,
   \tag{3.7.2}
\end{equation}
which gives rise to the commutativity diagrams needed for direct images. In the situations used here it is the transpose of the usual base-change morphism.
\end{statement}

\begin{statement}{3.8--3.9}
One defines in the same way the direct-image correspondence under a proper morphism. If $X\xrightarrow{u}Y\xrightarrow{g}S$ is a composite, the resulting operations satisfy
\[
   (g f)_*=g_*f_*.
\]
This is a direct verification from adjunction and the transitivity of the trace morphism.
\end{statement}

\subsubsection*{Lefschetz formula}

\begin{statement}{3.10}
Let $X$, $Y$, and $Z$ be $S$-schemes, and let
\[
   C\subset X\times_S Y,
   \qquad
   D\subset Y\times_S Z
\]
be correspondences. Their composite $D\circ C$ is defined by the usual fiber product over $Y$ and projection to $X\times_S Z$.
\end{statement}

\begin{statement}{Theorem 3.10.1}
With the preceding notation, if $u$ and $v$ are composable cohomological correspondences, then
\[
   \Tr_{Z/S}\bigl((v\circ u)_*\bigr)=\Tr_{Y/S}(v_*\circ u_*).
\]
This is the Lefschetz trace formula at the level of cohomological correspondences.
\end{statement}

\begin{statement}{3.11--3.12}
Interpreting the same construction in étale cohomology, cohomological correspondences act on $\ell$-adic cohomology and recover the familiar fixed-point formulae. If $X$ is a scheme over $\FF_q$ and $F:X\to X$ is Frobenius, the associated correspondence gives the Grothendieck--Lefschetz trace formula.
\end{statement}

\section*{4. Cup-products and trace formulae}

\begin{statement}{4.0}
We keep the preceding notation. Let $X$ and $Y$ be $S$-schemes, and let
\[
   L_1,L_2\in\ob D^-(X),\qquad M_1,M_2\in\ob D^-(Y).
\]
\end{statement}

\subsection*{4.1. Cup-product pairings}
Let $C\subset X\times_S Y$, with projections $p:C\to X$ and $q:C\to Y$. There are natural pairings
\begin{equation}
\Hom_C(p^*L_1,q^*M_1)\otimes \Hom_C(p^*L_2,q^*M_2)
\longrightarrow
\Hom_C\bigl(p^*(L_1\otimes^L L_2),q^*(M_1\otimes^L M_2)\bigr).
\tag{4.1.1}
\end{equation}
In particular, after specializing $L_1$ and $M_2$ to the structure object, one obtains
\begin{equation}
\Hom_C(p^*L,q^*M)\otimes \Hom_C(q^*M,p^*L)
\longrightarrow \Hom_C(\cO_C,\cO_C)=H^0(C,\cO_C).
\tag{4.1.2}
\end{equation}
There are analogous pairings for the internal derived Hom complexes:
\begin{equation}
\underline{\R\Hom}_C(p^*L_1,q^*M_1)\otimes^L
\underline{\R\Hom}_C(p^*L_2,q^*M_2)
\longrightarrow
\underline{\R\Hom}_C\bigl(p^*(L_1\otimes^L L_2),q^*(M_1\otimes^L M_2)\bigr).
\tag{4.1.4}
\end{equation}

\begin{statement}{4.2. Local cohomology}
Let $c:C\to X\times_S Y$ and $d:D\to X\times_S Y$ be morphisms, and let $L,M,L',M'$ be complexes on the two factors. The cup-product in cohomology with supports has the form
\begin{equation}
H^i_C(X\times Y,p_1^*L\otimes^L p_2^*M)\otimes
H^j_D(X\times Y,p_1^*L'\otimes^L p_2^*M')
\longrightarrow
H^{i+j}_{C\cap D}(\cdots).
\tag{4.2.1}
\end{equation}
For proper supports and correspondences
\[
   u\in\Hom(c^*p_1^*L,c^!p_2^*M),
   \qquad
   v\in\Hom(d^*p_1^*L',d^!p_2^*M'),
\]
one defines
\begin{equation}
   \langle u,v\rangle\in \Hom(\R\Gamma_{C\cap D}(\cO),K_X)
   \tag{4.2.3}
\end{equation}
by the formula
\begin{equation}
   \langle u,v\rangle=\Tr(u\cup v).
   \tag{4.2.4}
\end{equation}
\end{statement}

\subsection*{4.3. Projection formula}
The projection formulae for these cup-products are obtained from the usual projection formula in the derived category, together with the compatibility of traces with adjunction. They express that pushing forward one factor before taking the pairing gives the same result as first taking the pairing and then applying the trace.

\subsection*{4.4. The principal formula}
\begin{statement}{Theorem 4.4.1}
Let $X_i$, $Y_i$ ($i=1,2$) be $S$-schemes, and let $f_i:X_i\to Y_i$ be proper morphisms. Consider the cube 4.4.0 with cartesian horizontal faces and with $c'$, $c''$, $d'$, $d''$ proper. Under the stated tor-independence assumptions, and for $L_i\in\ob D(X_i)_{\mathrm{parf}/S}$, the square
\[
\begin{array}{ccc}
 f_*c'_*\underline{\R\Hom}(c'^*L_1,c'^!L_2)\otimes
 f_*c''_*\underline{\R\Hom}(c''^*L_2,c''^!L_1)
 & \xrightarrow{(1)} &
 f_*\underline{\R\Hom}(\R c'_*\cO_{C'},\cO_X) \\
 \downarrow & & \downarrow \\
 d'_*\underline{\R\Hom}(d'^*M_1,d'^!M_2)\otimes
 d''_*\underline{\R\Hom}(d''^*M_2,d''^!M_1)
 & \xrightarrow{(3)} &
 \underline{\R\Hom}(\R d'_*\cO_{D'},\cO_Y)
\end{array}
\]
is commutative. Here $M_i=f_{i*}L_i$; the arrow (3) is given by the cup-product 4.2.4; and (1) is obtained from the same construction after applying $f_*$ and composing with the canonical map from the tensor product of direct images to the direct image of a tensor product.
\end{statement}

\begin{statement}{Corollary 4.5}
The Lefschetz formula for cohomological correspondences follows from 4.4.1.
\end{statement}

\begin{statement}{4.6. Trace calculation}
For
\[
   u\in\Hom(c_1^*L_1,c_2^!L_2),
   \qquad
   v\in\Hom(c_2^*L_2,c_1^!L_1),
\]
one has
\begin{equation}
   \Tr(f_*u,f_*v)=f_*\Tr(u,v).
   \tag{4.6.1}
\end{equation}
\end{statement}

\begin{statement}{4.7. Tensor products of correspondences}
For $u_i\in\Hom_{X_i}(L_i,M_i)$ one defines
\[
   u_1\boxtimes u_2\in
   \Hom_{X_1\times X_2}(L_1\boxtimes L_2,M_1\boxtimes M_2).
\]
The trace is multiplicative:
\begin{equation}
   \Tr(u_1\boxtimes u_2)=\Tr(u_1)\Tr(u_2).
   \tag{4.7.1}
\end{equation}
\end{statement}

\section*{5. Complements}

\subsection*{5.1. Duality and Hom-type functors}
For a morphism $f:X\to Y$ of $S$-schemes and for $L,M\in\ob D(X)$ put
\begin{equation}
   \underline{\R\Hom}_f(L,M)=\underline{\R\Hom}(L,f^!M).
   \tag{5.1.1}
\end{equation}
For $M\in\ob D^{-}(Y)$ and $L\in\ob D^{+}(X)$ there are adjunction isomorphisms
\begin{equation}
   \Hom_X(L,f^!M)\xrightarrow{\sim}\Hom_Y(\R f_!L,M).
   \tag{5.1.2}
\end{equation}
The associated Ext-type functors are
\begin{equation}
   \Ext^i_f(L,M)=H^i\bigl(\underline{\R\Hom}_f(L,M)\bigr).
   \tag{5.2.1}
\end{equation}
For a composite $X\xrightarrow{f}Y\xrightarrow{g}Z$ there are transitivity isomorphisms
\begin{equation}
   \underline{\R\Hom}_{gf}(L,N)\xrightarrow{\sim}
   \underline{\R\Hom}_f(L,g^!N).
   \tag{5.3.1}
\end{equation}

\subsection*{5.4. Base change}
For a cartesian square
\[
\begin{array}{ccc}
X' & \xrightarrow{g'} & X \\
\downarrow f' & & \downarrow f \\
Y' & \xrightarrow{g} & Y,
\end{array}
\]
one has a canonical morphism
\begin{equation}
   g^*\underline{\R\Hom}_f(L,M)	o
   \underline{\R\Hom}_{f'}(g'^*L,g^*M),
   \tag{5.4.1}
\end{equation}
which is an isomorphism when $g$ is flat, or when $L$ is perfect.

\subsection*{5.5--5.7. Künneth formulae}
For $X_1$ and $X_2$ over $S$, and $L_i\in\ob D^{-}(X_i)$, the Künneth morphism is
\begin{equation}
   \R p_{1*}L_1\otimes^L_{\cO_S}\R p_{2*}L_2
      \longrightarrow \R p_*(L_1\boxtimes L_2),
   \tag{5.5.1}
\end{equation}
where $p_i:X_1\times_S X_2\to X_i$ and $p:X_1\times_S X_2\to S$. It is an isomorphism under tor-independence hypotheses. Similarly, for $X_i$, $Y_i$ and $f_i:X_i\to Y_i$, one has
\begin{equation}
   f_*\underline{\R\Hom}(L_1,M_1)\otimes
   f_*\underline{\R\Hom}(L_2,M_2)	o
   f_*\underline{\R\Hom}(L_1\boxtimes L_2,M_1\boxtimes M_2).
   \tag{5.6.1}
\end{equation}
The Künneth morphisms are compatible with base change, composition, and the cup-product pairings.

\subsection*{5.8--5.9. Complements on correspondences}
If
\[
   C\subset X\times_S Y,
   \qquad
   D\subset Y\times_S Z,
\]
then the composite correspondence is
\[
   D\circ C=\pr_{13}(C\times_YD).
\]
For composable cohomological correspondences $u$ and $v$, the trace pairing satisfies
\begin{equation}
   \langle v\circ u\rangle=\langle v,u\rangle,
   \tag{5.9.1}
\end{equation}
in the sense of the pairings of 4.2.

\section*{6. Cohomological operations and the Lefschetz--Verdier formula}

\begin{statement}{6.0}
This section gives generalizations of the preceding constructions, leading to the Lefschetz--Verdier formula and to the Woods Hole formula.
\end{statement}

\subsection*{6.1--6.2. Derived categories and relative duality}
Let $f:X\to S$ be a morphism of schemes. Write $D(X)$ for the derived category of sheaves of $\cO_X$-modules, and $D(X)_{\mathrm{parf}}$ for the full subcategory of perfect complexes. If $f$ is proper and of finite tor-dimension, there is a functor
\begin{equation}
   f^!:D^{+}(S)\to D^{+}(X)
   \tag{6.1.1}
\end{equation}
right adjoint to $\R f_*$, together with a duality isomorphism
\begin{equation}
   \underline{\R\Hom}_X(L,f^!M)
      \xrightarrow{\sim}
   \underline{\R\Hom}_S(\R f_*L,M).
   \tag{6.1.2}
\end{equation}
The relative dualizing object is
\[
   K_{X/S}=f^!\cO_S.
\]
If $f$ is smooth, purely of relative dimension $n$, then
\[
   K_{X/S}=\Omega^n_{X/S}[n].
\]
For such $f$ one has trace morphisms, for instance
\begin{equation}
   \Tr_f:R^{2n}f_*\Omega^n_{X/S}\to \cO_S,
   \tag{6.2.1}
\end{equation}
and more generally, for $K\in\ob D^b_{\mathrm{parf}}(X)$,
\begin{equation}
   \Tr_f:
   \R f_*\underline{\R\Hom}(K,K_{X/S})\to \cO_S.
   \tag{6.2.2}
\end{equation}

\subsection*{6.3. General Künneth formulae}
In the situation of 1.6, suppose that $X_1$ and $X_2$ are tor-independent over $S$, and likewise $Y_1$ and $Y_2$. If $L_i\in\ob D(X_i)_{\mathrm{parf}/S}$, the Künneth arrow 1.6.4 is defined, using the fact that finite relative tor-dimension is preserved under direct image. It is an isomorphism: the verification is reduced by affine localization and a Čech calculation to the affine case, where the assertion follows from tor-independence.

For $M_i\in\ob D(Y_i)_{\mathrm{parf}/S}$, with $f_1^!M_1$ or $f_2^!M_2$ bounded, the Künneth arrow 1.7.3 is also defined, and is an isomorphism when $f_1$ and $f_2$ are of finite tor-dimension. In particular, if $X_1$ and $X_2$ are of finite tor-dimension over $S$ and tor-independent, 1.7.6 is an isomorphism.

One defines, as in 2.1, the complexes $\underline{\R\Hom}(u,v)$ for pairs of arrows of type $(2,1)$ and of type $(3,4)$, assuming $f$ proper in order to avoid the pro-objects that appear with $f_!$. Subject to the appropriate boundedness conditions, the evaluation and tensor-product arrows exist and behave as in the preceding sections.

\subsection*{6.4. Cohomological correspondences with support}
Let $X_1$ and $X_2$ be $S$-schemes, and let $L_i\in\ob D(X_i)$. One considers four kinds of arrows over a morphism $f$:
\begin{align}
   \Hom_f^{(1)}(L,M)&=\Hom(M,f_*L),       \tag{6.4.1}\\
   \Hom_f^{(2)}(L,M)&=\Hom(f^*M,L),       \tag{6.4.2}\\
   \Hom_f^{(3)}(L,M)&=\Hom(L,f^!M),       \tag{6.4.3}\\
   \Hom_f^{(4)}(L,M)&=\Hom(f^!M,L).       \tag{6.4.4}
\end{align}
Their elements are called arrows of type $i$ from $L$ to $M$ over $f$, or $f$-arrows of type $i$. The partial adjunction formulae recover the usual descriptions; for instance, with suitable boundedness,
\[
   \Hom_f^{(1)}(L,M)=\Hom(M,f_*L),
\]
and, for $f$ proper,
\[
   \Hom_f^{(3)}(L,M)=\Hom(f_*L,M).
\]

\subsection*{6.5. Direct images}
Let $f_i:X_i\to Y_i$ ($i=1,2$) be proper $S$-morphisms, and put $f=f_1\times_S f_2:X\to Y$. Assume the finite tor-dimension and tor-independence hypotheses appearing in 2.4.0. If $L_i\in\ob D(X_i)_{\mathrm{parf}/S}$, then, as in 3.3.1, there is an isomorphism
\begin{equation}
   f_*\underline{\R\Hom}_S(L_1,L_2)
      \xrightarrow{\sim}
   \underline{\R\Hom}_S(f_{1*}L_1,f_{2*}L_2),
   \tag{6.5.1}
\end{equation}
and hence, after applying $\R\Gamma(Y,-)$,
\begin{equation}
   f_*:\Hom_S(L_1,L_2)	o \Hom_S(f_{1*}L_1,f_{2*}L_2).
   \tag{6.5.2}
\end{equation}
In the case $Y_1=Y_2=S$, this is again written $u\mapsto u_*$.

For a square of the form 3.7.1 with $c$ proper and $K\in\ob D^+(C)$, there is a canonical arrow
\begin{equation}
   f_*c_*c^!K\to d_*d^!f_*K,
   \tag{6.5.3}
\end{equation}
which is defined as in 3.7.3. Taking $K=\underline{\R\Hom}_S(L_1,L_2)$ and using 6.5.1 gives
\begin{equation}
   f_*c_*\underline{\R\Hom}(c_1^*L_1,c_2^!L_2)	o
   d_*\underline{\R\Hom}\bigl(d_1^*(f_{1*}L_1),d_2^!(f_{2*}L_2)\bigr),
   \tag{6.5.4}
\end{equation}
and on global Hom groups,
\begin{equation}
   f_*:\Hom(c_1^*L_1,c_2^!L_2)	o
   \Hom\bigl(d_1^*(f_{1*}L_1),d_2^!(f_{2*}L_2)\bigr).
   \tag{6.5.5}
\end{equation}
These are the analogues of 3.7.5 and 3.7.6.

\subsection*{6.6. Cup-products and pairings}
Let $X_i$ be $S$-schemes of finite tor-dimension and tor-independent, and let $L_i\in\ob D(X_i)_{\mathrm{parf}/S}$. The exterior tensor product of the evaluation morphisms
\[
   \underline{D}L_i\otimes L_i\to K_{X_i}
\]
defines cup-product pairings analogous to 4.1.2--4.1.4.

Let $c:C\to X$ and $d:D\to X$ be proper morphisms, and let $e=c\times_X d:E\to X$. For $P,Q\in\ob D(X)_{\mathrm{parf}/S}$ one has
\[
   c_*c^!P=\underline{\R\Hom}(\R c_*\cO_C,P),
\]
and similarly for $d$. Assuming $c$ and $d$ of finite tor-dimension at the relevant points, the tensor product gives
\begin{equation}
   c_*c^!P\otimes d_*d^!Q
   \to
   \underline{\R\Hom}(\R c_*\cO_C\otimes \R d_*\cO_D,P\otimes Q).
   \tag{6.6.1}
\end{equation}
When $c$ and $d$ are tor-independent, the right-hand side is rewritten, by Künneth and 6.1.1, as $e_*e^!(P\otimes Q)$. Consequently one obtains pairings
\begin{equation}
\begin{split}
 &c_*\underline{\R\Hom}(c_1^*L_1,c_2^!L_2)
 \otimes d_*\underline{\R\Hom}(d_2^*L_2,d_1^!L_1)\\
 &\hspace{3em}\longrightarrow
 \underline{\R\Hom}(\R c_*\cO_C\otimes \R d_*\cO_D,K_X),
\end{split}
\tag{6.6.2}
\end{equation}
and, on global Hom groups,
\begin{equation}
   \langle\ ,\ \rangle:
   \Hom(c_1^*L_1,c_2^!L_2)\otimes
   \Hom(d_2^*L_2,d_1^!L_1)	o
   \Hom(\R c_*\cO_C\otimes \R d_*\cO_D,K_X).
   \tag{6.6.3}
\end{equation}
These pairings are compatible with étale localization. If $c$ and $d$ are tor-independent, the target is canonically $e_*K_E$, respectively $H^0(E,K_E)$.

\subsection*{6.7. Lefschetz--Verdier formula}
\begin{statement}{Theorem 6.7}
Let $f_i:X_i\to Y_i$ ($i=1,2$) be proper $S$-morphisms, put $f=f_1\times_S f_2:X\to Y$, and let $p_i:X\to X_i$, $q_i:Y\to Y_i$ be the projections. Suppose given a cube as in 4.4.0, with cartesian horizontal faces, with $c'$, $c''$, $d'$, $d''$ proper, and with the required finite tor-dimension assumptions along the supports. Let $L_i\in\ob D(X_i)_{\mathrm{parf}/S}$ and suppose the pairs $(X_1,X_2)$, $(X_1,Y_2)$, $(Y_1,X_2)$, and $(Y_1,Y_2)$ are tor-independent over $S$, while $X_i$ and $Y_i$ have finite tor-dimension over $S$. Then the square
\[
\begin{array}{ccc}
 f_*c'_*\underline{\R\Hom}(c_1'^*L_1,c_2'^!L_2)
 \otimes f_*c''_*\underline{\R\Hom}(c_2''^*L_2,c_1''^!L_1)
 &\to&
 f_*\underline{\R\Hom}(\R c'_*\cO_{C'}\otimes \R c''_*\cO_{C''},K_X)\\
 \downarrow && \downarrow\\
 d'_*\underline{\R\Hom}(d_1'^*M_1,d_2'^!M_2)
 \otimes d''_*\underline{\R\Hom}(d_2''^*M_2,d_1''^!M_1)
 &\to&
 \underline{\R\Hom}(\R d'_*\cO_{D'}\otimes \R d''_*\cO_{D''},K_Y)
\end{array}
\]
is commutative, where $M_i=f_{i*}L_i$. The horizontal arrows are the cup-product pairings of 6.6.2; the left vertical arrow is the tensor product of the direct-image maps 6.5.4; the right vertical arrow is obtained from relative duality and the canonical base-change morphisms.
\end{statement}

The proof is the same as that of 4.4, with the necessary substitutions in the duality formalism. If $c'$ and $c''$ are tor-independent, and likewise $d'$ and $d''$, the right vertical arrow can be read simply as
\[
   f_*c_*K_C\to d_*K_D,
\]
given by the adjunction map $g_*K_C=g_*g^!K_D\to K_D$. The general statement is kept because it includes excess-intersection cases.

As a consequence, under the hypotheses of 6.7, for
\[
   u\in\Hom(c_1'^*L_1,c_2'^!L_2),
   \qquad
   v\in\Hom(c_2''^*L_2,c_1''^!L_1),
\]
one has
\begin{equation}
   \langle f_*u,f_*v\rangle=f_*\langle u,v\rangle,
   \tag{6.7.2}
\end{equation}
where $f_*u$ and $f_*v$ are direct-image correspondences in the sense of 6.5.5, and where the $f_*$ on the right is the morphism
\begin{equation}
\Hom(\R c'_*\cO_{C'}\otimes \R c''_*\cO_{C''},K_X)	o
\Hom(\R d'_*\cO_{D'}\otimes \R d''_*\cO_{D''},K_Y).
\tag{6.7.3}
\end{equation}
When $Y_1=Y_2=S=D'=D''$, formula (6.7.2) is the Lefschetz--Verdier formula.

\subsection*{6.8. The Hodge cohomology class}
Let $f:X\to S$ be smooth, purely of relative dimension $N$, and let $i:Y\hookrightarrow X$ be a regular immersion of codimension $d$; equivalently, $Y$ is locally cut out by a regular sequence of $d$ equations, or is a local complete intersection of codimension $d$. Put $g=fi$. One attaches to $Y$ an element
\begin{equation}
   \cl_{X/S}(Y)\in\Ext^d_{\cO_X}(\cO_Y,\Omega^d_{X/S}),
   \tag{6.8.1}
\end{equation}
called the \emph{cohomology class associated with} $Y$.

By the purity theorem, one has
\begin{equation}
\underline{\Ext}^{i}_{\cO_X}(\cO_Y,\Omega^d_{X/S})=
\begin{cases}
0, & i\ne d,\\
\underline{\Hom}(\Lambda^dN_{Y/X},\Omega^d_{X/S}\otimes\cO_Y), & i=d,
\end{cases}
\tag{6.8.2}
\end{equation}
where $N_{Y/X}=I/I^2$ is the conormal sheaf, $I=\Ker(\cO_X\to\cO_Y)$. Therefore
\begin{equation}
\Ext^d_{\cO_X}(\cO_Y,\Omega^d_{X/S})=
H^0\bigl(Y,\underline{\Hom}(\Lambda^dN_{Y/X},\Omega^d_{X/S}\otimes\cO_Y)\bigr).
\tag{6.8.3}
\end{equation}
If $Y$ is defined by the ideal generated by a regular sequence $t_1,\ldots,t_d$, then $\Lambda^dN_{Y/X}$ is generated by $\bar t_1\wedge\cdots\wedge\bar t_d$, and $\cl_{X/S}(Y)$ corresponds to
\begin{equation}
   \bar t_1\wedge\cdots\wedge\bar t_d
   \longmapsto dt_1\wedge\cdots\wedge dt_d.
   \tag{6.8.5}
\end{equation}
It is multiplicative for transversal intersections:
\begin{equation}
   \cl_{X/S}(Y\cap Z)=\cl_{X/S}(Y)\cl_{X/S}(Z).
   \tag{6.8.6}
\end{equation}
If $Y$ has codimension equal to the relative dimension and is finite flat over $S$, the trace is
\begin{equation}
   g_*g^*F=F\otimes g_*\cO_Y\to F,
   \tag{6.8.7}
\end{equation}
induced by $\Tr_{Y/S}:g_*\cO_Y\to\cO_S$.

\subsection*{6.9--6.10. Woods Hole formula}
Let $X$ and $Y$ be proper smooth $S$-schemes, of relative dimensions $m$ and $n$, and set $Z=X\times_SY$. Let
\[
   a:A\hookrightarrow Z,
   \qquad
   b:B\hookrightarrow Z
\]
be regular immersions such that $a_2:A\to Y$ and $b_1:B\to X$ are finite and flat. Write $p_1:Z\to X$, $p_2:Z\to Y$, $a_i=p_ia$, $b_i=p_ib$, and put $C=A\times_ZB$. Assume $C$ is finite flat over $S$. Let $L\in\ob D(X)_{\mathrm{parf}}$, $M\in\ob D(Y)_{\mathrm{parf}}$.

For every $r$ one has the decomposition
\begin{equation}
   \Omega^r_{Z/S}=\bigoplus_{i+j=r}\Omega^i_{Z/Y}\otimes\Omega^j_{Z/X}.
   \tag{6.9.0}
\end{equation}
The identifications of $p_2^!M$ and $p_1^!L$ give
\begin{equation}
\Hom(a_1^*L,a_2^!M)=
H^0\bigl(A,\underline{\R\Hom}(a_1^*L,a_2^*M)\otimes
\Lambda^mN_{A/Z}^{\vee}\otimes\Omega^m_{Z/Y}\bigr),
\tag{6.9.1}
\end{equation}
and similarly
\begin{equation}
\Hom(b_2^*M,b_1^!L)=
H^0\bigl(B,\underline{\R\Hom}(b_2^*M,b_1^*L)\otimes
\Lambda^nN_{B/Z}^{\vee}\otimes\Omega^n_{Z/X}\bigr).
\tag{6.9.2}
\end{equation}

Let $u\in\Hom(a_1^*L,a_2^*M)$ and $v\in\Hom(b_2^*M,b_1^*L)$. On $C$ define
\begin{equation}
   \langle u,v\rangle=\Tr(u_Cv_C)=\Tr(v_Cu_C)\in H^0(C,\cO_C).
   \tag{6.9.3}
\end{equation}
The cohomological correspondences $u^!$ and $v^!$ are obtained from $u$ and $v$ by composition with the trace morphisms, equivalently by multiplication with $\cl_{Z/Y}(A)$ and $\cl_{Z/X}(B)$. Their cup-product is
\begin{equation}
   \langle u^!,v^!\rangle=
   \langle u,v\rangle\cl_{Z/Y}(A)\cl_{Z/X}(B).
   \tag{6.9.4}
\end{equation}
The Lefschetz--Verdier formula gives
\begin{equation}
   \langle (u^!)_*,(v^!)_*\rangle=
   \Tr_{Z/S}\bigl(\langle u,v\rangle\cl_{Z/Y}(A)\cl_{Z/X}(B)\bigr).
   \tag{6.9.5}
\end{equation}
Thus:

\begin{statement}{Theorem 6.10}
Under the hypotheses of 6.9,
\[
   \Tr\bigl((u^!)_*\circ (v^!)_*\bigr)=
   \Tr_{Z/S}\bigl(\langle u,v\rangle\cl_{Z/Y}(A)\cl_{Z/X}(B)\bigr).
\]
\end{statement}

For $L=\cO_X$, $M=\cO_Y$, and $u,v$ the identity maps, this recovers the classical intersection-number formula. The same formalism admits variants for other types of correspondences, singular schemes, continuous coefficients, and the usual applications to étale trace formulae and congruence formulae.

\section*{Exposé III B. Calculation of local terms}
\begin{center}
{\large\bfseries by Luc Illusie}
\end{center}

\section*{0. Introduction}
The purpose of this exposé is to calculate the local terms occurring in the Lefschetz--Verdier trace formula for several kinds of correspondences. The results generalize those of Exposé III A and complete the formulae of SGA 5, Exposé III.

Let $X$ be a scheme of finite type over an algebraically closed field $k$ of characteristic $p$. One considers cohomological correspondences
\[
   u:c_1^*L\to c_2^!L,
   \qquad
   c=(c_1,c_2):C\to X\times X.
\]
The local term of $u$ at an isolated fixed point $x\in X$ is defined as the trace of the endomorphism induced by $u$ on the appropriate local complex. The plan is as follows. Section 1 treats correspondences in general position between curves. Section 2 reduces the proof to an assertion of vanishing. Section 3 reduces to the tame case. Section 4 finishes the proof of the main theorem. The second part treats equivariant correspondences and non-commutative traces.

The author thanks P. Deligne for useful suggestions made in many conversations, and in particular for pointing out and helping to correct an error in an early proof of Theorem 1.2.

\section*{I. Correspondences in general position between curves}

Throughout this part, $S=\Spec(k)$ with $k$ algebraically closed of characteristic $p$, and $\Lambda$ is a commutative noetherian ring killed by an integer prime to $p$. For every $S$-scheme $X$, let $D(X)$ be the derived category of sheaves of $\Lambda$-modules on $X$, and let $D_{\mathrm{ctf}}(X)$ be the full subcategory formed by complexes of finite tor-dimension with constructible cohomology.

\subsection*{1. Statement of the theorem and corollaries}
Let $X$ and $Y$ be the strict localizations at closed points of smooth separated $S$-curves. Let $A$ and $B$ be strictly local $S$-schemes, and let
\[
   a:A\to X\times_S Y,
   \qquad
   b:B\to X\times_S Y
\]
be $S$-morphisms. Write $p_1,p_2$ for the projections, set $a_i=p_ia$, $b_i=p_ib$, and denote by $x$ and $y$ the closed points of $X$ and $Y$. Assume that $a_2$ and $b_1$ are finite and surjective, and that
\[
   A\times_{X\times_SY}B
\]
is reduced to the closed point $z$ mapping to $(x,y)$.

Let
\[
   L\in\ob D_{\mathrm{ctf}}(X),\qquad
   M\in\ob D_{\mathrm{ctf}}(Y),
\]
and let
\[
   u\in\Hom(a_1^*L,a_2^!M),
   \qquad
   v\in\Hom(b_2^*M,b_1^!L).
\]
The induced map
\[
   u_z:L_x=\R\Gamma(X,L)\to \R\Gamma(Y,M)=M_y
\]
is the composite
\[
L_x\xrightarrow{a_1^*}\R\Gamma(A,a_1^*L)
\xrightarrow{\R\Gamma(A,u)}\R\Gamma(A,a_2^!M)
\xrightarrow{a_{2*}}\R\Gamma(Y,M),
\]
where $a_{2*}$ is defined by the adjunction map $a_{2*}a_2^!M\to M$. The morphism $v_z:M_y\to L_x$ is defined similarly. Since $L_x$ and $M_y$ are perfect complexes of $\Lambda$-modules, one can form
\begin{equation}
   \langle u_z,v_z\rangle=
   \Tr(u_zv_z)=\Tr(v_zu_z)\in\Lambda.
   \tag{1.1.1}
\end{equation}

Put $K_X=\Lambda_X(1)[2]$ and $K_Y=\Lambda_Y(1)[2]$. By duality, the functors
\[
   D_X=\R\Hom(-,K_X),
   \qquad
   D_Y=\R\Hom(-,K_Y)
\]
are dualizing on $D_{\mathrm{ctf}}(X)$ and $D_{\mathrm{ctf}}(Y)$. As in III 4.1.2 and III 4.2.3, the cohomological correspondences $u$ and $v$ have a cup-product local term
\begin{equation}
   \langle u,v\rangle\in\Lambda.
   \tag{1.1.2}
\end{equation}
If the data are obtained by strict localization from global curves $X'$, $Y'$, $A'$, $B'$, this is the local term of III 4.2.7 at the isolated point $z'$.

In general the two quantities (1.1.1) and (1.1.2) are not equal. The main theorem of the first part gives a sufficient general-position condition.

\begin{statement}{Theorem 1.2}
Under the hypotheses of 1.1, assume in addition that $A$ and $B$ are in general position: the intersection of the tangent cones at $(x,y)$ to their images in $X\times_S Y$ is reduced to $(x,y)$, and similarly $A$ and $X\times\{y\}$ are in general position, as are $\{x\}\times Y$ and $B$. Then
\begin{equation}
   \langle u_z,v_z\rangle=\langle u,v\rangle.
   \tag{1.2.1}
\end{equation}
\end{statement}

\begin{statement}{Corollary 1.3}
Let $F$ be an endomorphism of a proper smooth curve $X/S$ such that the fixed-point scheme $X^F$ is finite and consists of points at which the graph of $F$ is transverse to the diagonal. If $L\in\ob D_{\mathrm{ctf}}(X)$ and $u\in\Hom(F^*L,L)$, then
\[
   \Tr((F,u),\R\Gamma(X,L))=
   \sum_{x\in X^F}\Tr(u_x,L_x).
\]
For $L$ concentrated in degree $0$ this is the Artin--Verdier result. Passing to the limit gives the usual $\ell$-adic form.
\end{statement}

\begin{statement}{Corollary 1.3.1}
With $X$ and $F$ as above, let $L$ be a constructible $\bar{\QQ}_\ell$-sheaf on $X$, with $\ell\ne p$, and let $u\in\Hom(F^*L,L)$. Then
\[
   \sum_i(-1)^i\Tr((F,u),H^i_c(\bar X,L))
      =\sum_{x\in X^F}\Tr(u_x,L_x).
\]
\end{statement}

The hypotheses apply in particular when $(F,u)$ is a Frobenius correspondence; one obtains Grothendieck's trace formula.

There is also a dual form. If $b_2$ and $a_1$ are finite and surjective, consider the dual correspondences
\[
   Du\in\Hom(a_2^*DM,a_1^!DL),
   \qquad
   Dv\in\Hom(b_1^*DL,b_2^!DM).
\]
Under the corresponding general-position assumptions, applying Theorem 1.2 to $Du$ and $Dv$ yields
\begin{equation}
   \langle u,v\rangle=\langle (Du)_z,(Dv)_z\rangle.
   \tag{1.4.1}
\end{equation}

\begin{statement}{Corollary 1.5}
Let $X$ and $Y$ be proper smooth curves over $S$, and let $a:A\to X\times_SY$, $b:B\to X\times_SY$ be correspondences with $A$ and $B$ proper curves over $S$. Let $L\in\ob D_{\mathrm{ctf}}(X)$, $M\in\ob D_{\mathrm{ctf}}(Y)$, and let $u$ and $v$ be cohomological correspondences as above. Assume $A\times_{X\times_SY}B$ is finite over $S$ and decomposes as $C=C_d\cup C_t$, where the points of $C_d$ satisfy the general-position hypotheses for $(a_2,b_1)$ and the points of $C_t$ satisfy the dual general-position hypotheses. If $\underline u$ and $\underline v$ are the induced morphisms on cohomology, then
\begin{equation}
   \langle \underline u,\underline v\rangle=
   \sum_{z\in C_d}\langle u_z,v_z\rangle+
   \sum_{z\in C_t}\langle (Du)_z,(Dv)_z\rangle.
   \tag{1.5.1}
\end{equation}
\end{statement}

The $\ell$-adic limit form of this corollary contains Langlands' Proposition 7.12 as a special case. For the identity coefficients, taking $u$ and $v$ to be the fundamental classes of $A$ and $B$, Theorem 1.2 reads
\begin{equation}
   \deg(a_2)\deg(b_1)=\langle \cl(A),\cl(B)\rangle,
   \tag{1.7.1}
\end{equation}
which is the intersection multiplicity of the direct-image cycles under the stated general-position hypotheses.

\subsection*{2. Reduction to a vanishing theorem}
Let $U=X-\{x\}$ and $V=Y-\{y\}$, with inclusions $i:U\hookrightarrow X$ and $j:V\hookrightarrow Y$. Since $a_2$ and $b_1$ are finite, $a_2^{-1}(y)=s$ and $b_1^{-1}(x)=t$, so the inverse images of $U$ and $V$ fit into cartesian diagrams.

\begin{statement}{Lemma 2.2}
To prove Theorem 1.2 it suffices to treat the corresponding graded cases obtained from the filtrations
\[
F^nL=\begin{cases}L,& n<0,\\ i_!i^*L,& n=1,\\0,& n>1,
\end{cases}
\qquad
F^nM=\begin{cases}M,& n<0,\\ j_!j^*M,& n=1,\\0,& n>1.
\end{cases}
\]
The morphisms $u$ and $v$ are morphisms of filtered complexes and, by additivity of cup-products,
\[
   \langle u,v\rangle=
   \langle \operatorname{gr}^0u,\operatorname{gr}^0v\rangle+
   \langle \operatorname{gr}^1u,\operatorname{gr}^1v\rangle.
\]
The graded pieces reduce the theorem to the cases supported at the closed points and on their complements.
\end{statement}

\begin{statement}{Lemma 2.3}
Under the hypotheses of 1.1, if $L|_U=0$ and $M|_V=0$, then
\[
   \langle u_z,v_z\rangle=\langle u,v\rangle.
\]
The assertion is a direct form of the Lefschetz formula for the projections $\{x\}\to S$ and $\{y\}\to S$.
\end{statement}

The remaining case is reduced to a vanishing statement. After replacing the data by global curves and by constructible bounded complexes, one is led to prove the vanishing of a restriction morphism of the form
\[
   \R\Gamma(T-A,i_!E)	o \R\Gamma(B-t,i_!E|_{B-t})
\]
for a two-dimensional strict local scheme $T$, with divisors $A$, $B$, and $X$ meeting in the prescribed general position. This is the proposition below, after the standard reductions on coefficients.

\subsection*{3. Reduction to the tame case}
The proof first reduces to the case where $\Lambda$ is killed by a power of a prime $\ell\ne p$ and where the constructible sheaves in question are locally free with finite monodromy through finite $\ell$-groups. The reductions use dévissage, passage to finite étale covers, and the behavior of the fundamental group under strict localization. The geometrical part is to show that the general-position hypotheses are preserved by the finite covers used in this reduction.

Let $\widetilde A_i$ and $\widetilde B_j$ be the normalizations of the branches of $A$ and $B$ at the closed point. The tangent directions of the branches are represented by non-zero elements $\lambda_i$ and $\mu_j$ of the residue field. General position is the condition
\begin{equation}
   \lambda_i\ne0,
   \qquad
   \mu_j\ne0,
   \qquad
   \lambda_i\ne\mu_j
   \quad\text{for all }i,j.
   \tag{3.3.1}
\end{equation}
After passing to the finite covers, the new tangent directions $\lambda_i'$ and $\mu_j'$ satisfy
\[
   \lambda_i'=C\lambda_i,
   \qquad
   \mu_j'=C\mu_j
\]
for a single $C\in k^*$; hence the inequalities (3.3.1) are preserved.

\subsection*{4. End of the proof of Theorem 1.2}
It remains to establish the required vanishing under the coefficient hypotheses just stated. The assertion is the following.

\begin{statement}{Proposition 4.2}
Let $T$ be the strict localization at a closed point of a smooth two-dimensional $S$-scheme, with closed point $t$. Let $A$, $B$, and $X$ be closed one-dimensional subschemes of $T$. Let
\[
   i:T-(A\cup X)\hookrightarrow T-A
\]
be the inclusion. Assume $X$ is regular, and that $A$ is in general position with respect to $B$ and $X$; that is, the intersection of the tangent cone of $A$ at $t$ with the tangent cone of $B\cup X$ is reduced to $t$. Let $E\in\ob D^b(T-(A\cup X))$ have constructible locally constant cohomology, with monodromy factoring through an $\ell$-group. Then the restriction map
\[
   \R\Gamma(T-A,i_!E)	o \R\Gamma(B-t,i_!E|_{B-t})
\]
is zero.
\end{statement}

The proof rests on the following lemma.

\begin{statement}{Lemma 4.3}
Let $Z$ be a connected regular noetherian scheme of dimension $2$, with residual characteristics prime to $\ell$. Let $X$ and $Y$ be regular divisors meeting transversely at a point $z$, giving a diagram of inclusions
\[
\begin{array}{ccccc}
Z-X & \xrightarrow{j} & Z\\
\uparrow m & & \uparrow i\\
Y-z & \xrightarrow{q} & Y.
\end{array}
\]
Assume purity:
\[
   R^qj_*(\ZZ/\ell)=0\qquad(q>1).
\]
Let $L\in\ob D^b(Z-(X\cup Y),\Lambda)$ have constructible locally constant cohomology with monodromy through an $\ell$-group. Then:
\begin{enumerate}[label=(\alph*)]
\item the natural map $L\to m^*Rj_*L$ is an isomorphism, and $m^*Rj_*L$ has constructible locally constant cohomology;
\item if $Y$ is connected and $\bar s$ is a geometric point of $Y$ localized at a closed point $s\ne z$, then
\[
   \R\Gamma(Y,m^*Rj_*L)	o (Rj_*L)_{\bar s}
\]
is zero;
\item under the hypotheses of (b), if $C$ is a one-dimensional closed subscheme of $Z$ with $Y\cap C=s$, and $T$ is the strict localization of $C$ at $s$, then
\[
   \R\Gamma(Z-Y,i_!L)	o \R\Gamma(T-s,i_!L|_{T-s})
\]
is zero.
\end{enumerate}
\end{statement}

\begin{prooftext}[Proof of Lemma 4.3]
The localization triangle gives an isomorphism
\[
   i^*Rj_*i_!L\xrightarrow{\sim}Rq_*q^*i_!L=Rq_*L.
\]
By adjunction this yields $i_!L\to m_*m^*Rj_*L$, an isomorphism away from $z$. It remains to check that the stalk at $\bar z$ of $Rj_*i_!L$ is zero. By dévissage one reduces to $L=\ZZ/\ell$ with trivial action. Then the assertion is a compact-support calculation: in degrees $<1$ it follows from Abhyankar's lemma, and in degrees $>1$ from purity.

For (b), choose specialization maps from a geometric generic point of $Y$ to the strict localizations at $z$ and $s$. The resulting square of stalks shows that the restriction morphism factors through the stalk at $z$, which is zero. Part (c) follows by applying (b) to the commutative square of restriction maps from $Z-Y$ to the punctured strict localization of $C$ at $s$.
\end{prooftext}

\begin{prooftext}[Proof of Proposition 4.2]
Let $\varepsilon:T'\to T$ be the blow-up of the maximal ideal of $T$, and let $D=\varepsilon^{-1}(t)$ be the exceptional divisor. Let $A'$, $B'$, and $X'$ be the pure transforms of $A$, $B$, and $X$. Then $T'$ is regular, connected, and two-dimensional; $X'$ is a regular divisor meeting $D$ transversely in one point; and the general-position hypotheses mean that $A'$ is disjoint from $B'\cup X'$.

Pull $E$ back to $T'-(A'\cup D\cup X')$. It suffices to prove the vanishing of the restriction
\[
   \R\Gamma(T'-(A'\cup D),i_!E')	o
   \R\Gamma(B'-(B'\cap D),i_!E'|_{B'-(B'\cap D)}).
\]
The scheme $B'$ is a disjoint union of strict local one-dimensional schemes $B'_j$, each meeting $D$ in a single point $b_j$. Apply Lemma 4.3(c) to
\[
   Z=T'-A',\qquad X=X',\qquad Y=D-(D\cap A'),\qquad C=B'_j,
\]
using the relative purity theorem to verify the purity hypothesis. Each resulting restriction map is zero, and therefore the displayed restriction is zero. This proves Proposition 4.2, hence the required vanishing, and so Theorem 1.2.
\end{prooftext}

\section*{II. Equivariant correspondences}

\subsection*{5. Non-commutative traces}
Fix, for this section, a topos $T$ and a commutative unital ring $K$ of $T$. The rings considered below are associative and unital, but not necessarily commutative. A $K$-algebra means an associative unital ring $A$ of $T$ together with a unital homomorphism $K\to A$ whose image is contained in the center of $A$.

\subsubsection*{5.0. Introduction}
Let $E$ be a $K$-module which is locally a direct factor of a finite free module. There is then a local trace homomorphism
\begin{equation}
   \Tr_K:\Hom_K(E,E)\to K,
   \tag{5.0.1}
\end{equation}
defined by composing the evaluation map
\begin{equation}
   \Hom_K(E,K)\otimes E\to K,
   \qquad f\otimes x\mapsto f(x),
   \tag{5.0.2}
\end{equation}
with the inverse of the canonical tensor-product isomorphism
\begin{equation}
   \Hom_K(E,K)\otimes E\xrightarrow{\sim}\Hom_K(E,E),
   \qquad f\otimes y\mapsto (x\mapsto f(x)y).
   \tag{5.0.3}
\end{equation}
For any complex $E$ of $K$-modules the same two maps induce morphisms of complexes
\begin{align}
   \Hom_K^{\bullet}(E,K)\otimes E&\to K,                         \tag{5.0.4}\\
   \Hom_K^{\bullet}(E,K)\otimes E&\to \Hom_K^{\bullet}(E,E).       \tag{5.0.5}
\end{align}
After derivation, for $E\in\ob D^-(K)$, these give
\begin{align}
   \R\Hom_K(E,K)\otimes_K^L E&\to K,                              \tag{5.0.6}\\
   \R\Hom_K(E,K)\otimes_K^L E&\to \R\Hom_K(E,E).                   \tag{5.0.7}
\end{align}
If $E$ is perfect, the second map is an isomorphism, and composing its inverse with the first gives the global trace
\begin{equation}
   \Tr_K:\R\Hom_K(E,E)\to K.                                      \tag{5.0.8}
\end{equation}

The point is to examine to what extent $K$ may be replaced by a $K$-algebra $A$, especially the algebra of a finite group. The analogue of (5.0.3) causes no difficulty; the evaluation map does. If $E$ is a left $A$-module, the map
\[
   \Hom_A(E,A)\otimes E\to A,
   \qquad f\otimes x\mapsto f(x),
\]
does not in general factor through $\Hom_A(E,A)\otimes_AE$. It factors, after quotienting, as
\begin{equation}
   \Hom_A(E,A)\otimes_AE\to A_K,                                  \tag{5.0.9}
\end{equation}
where $A_K$ is the quotient of $A$ by the $K$-submodule locally generated by the commutators $ab-ba$. If $E$ is locally a direct factor of a finite free left $A$-module, then
\begin{equation}
   \Hom_A(E,A)\otimes_AE\xrightarrow{\sim}\Hom_A(E,E),              \tag{5.0.10}
\end{equation}
and hence one obtains
\begin{equation}
   \Tr_A:\Hom_A(E,E)\to A_K.                                      \tag{5.0.11}
\end{equation}
For $A=K$ this is (5.0.1). When $T$ is the punctual topos, it agrees with the trace of Stallings and Grothendieck. The rest of the section derives these maps for complexes and studies their behavior under change of scalars and under exterior tensor products; the resulting constructions will be used in 6 to define non-commutative Lefschetz--Verdier local terms.

\subsubsection*{5.1. Notation}
For a $K$-algebra $A$, let $A^{\circ}$ be the opposite algebra and let
\[
   A^e=A\otimes_KA^{\circ}
\]
be its enveloping algebra. Giving an $A$-bimodule is the same as giving a left, or right, $A^e$-module. If $E$ is a bimodule, recall that
\[
   H_0(A,E)=A\otimes_{A^e}E=E/[A,E]
\]
is the quotient of $E$ by the submodule generated locally by elements $ax-xa$. We shall sometimes write $E_A$ for this quotient. Thus $A_K$ is the quotient introduced in (5.0.9). If $A=K[G]$, where $G$ is an ordinary finite group, then $A_K$ is the free $K$-module on the set of conjugacy classes of $G$.

For $K$-algebras $A$ and $B$ we use the Cartan--Eilenberg notation ${}_AE$, $E_B$, and ${}_AE_B$ for left $A$-modules, right $B$-modules, and $A$-$B$-bimodules. We write $D(A)$ for the derived category of left $A$-modules and $D(A,B)$ for the derived category of $A$-$B$-bimodules.

\subsubsection*{5.2--5.4. Adjunction, evaluation, and tensor-product maps}
For $K$-algebras $A$ and $B$, in the situation $({}_AE,{}_BF_A,{}_BG)$, there is the canonical adjunction
\begin{equation}
\Hom_A(E,\Hom_B(F,G))\xrightarrow{\sim}\Hom_B(F\otimes_AE,G),
\qquad f\mapsto (y\otimes x\mapsto f(x)(y)).                    \tag{5.2.1}
\end{equation}
It derives to
\begin{equation}
\R\Hom_A(E,\R\Hom_B(F,G))\xrightarrow{\sim}
\R\Hom_B(F\otimes_A^LE,G),                                      \tag{5.2.2}
\end{equation}
for $E\in\ob D^-(A)$, $F\in\ob D^-(B,A^{\circ})$, and $G\in\ob D^+(B)$. One also obtains
\begin{equation}
\R\Hom_{A^e}(E,\R\Hom_B(F,G))\xrightarrow{\sim}
\R\Hom_B(F\otimes_A^LE,G).                                      \tag{5.2.3}
\end{equation}

For a $K$-algebra $A$ and data $({}_AE,{}_{A^{\circ}}F)$, there is an isomorphism
\begin{equation}
   E\otimes_{A^e}F\xrightarrow{\sim}(E\otimes_KF)\otimes_{A^e}A.
   \tag{5.3.1}
\end{equation}
For left $A$-modules $E,F$, the evaluation map
\begin{equation}
   \Hom_A(E,F)\otimes_AE\to F,
   \qquad f\otimes x\mapsto f(x),                               \tag{5.3.2}
\end{equation}
is $A^e$-linear after passing to coinvariants, and induces
\begin{equation}
   \Hom_A(E,F)\otimes_{A^e}E\to F\otimes_{A^e}A=F_A.              \tag{5.3.3}
\end{equation}
If $A$ is flat over $K$, this derives to
\begin{equation}
   E\otimes_{A^e}^LF\xrightarrow{\sim}(E\otimes_K^LF)\otimes_{A^e}^LA, \tag{5.3.4}
\end{equation}
and to the derived evaluation map
\begin{equation}
   \R\Hom_A(E,F)\otimes_{A^e}^LE\to F\otimes_{A^e}^LA.             \tag{5.3.5}
\end{equation}

There is also a canonical tensor-product map, for $({}_AE,{}_AF_{A^{\circ}},{}_AG_A)$,
\begin{equation}
   \Hom_A(E,F)\otimes_AG\to \Hom_A(E,F\otimes_AG),
   \qquad u\otimes y\mapsto (x\mapsto u(x)\otimes y).             \tag{5.4.1}
\end{equation}
Its derived form is
\begin{equation}
   \R\Hom_A(E,F)\otimes_A^LG\to \R\Hom_A(E,F\otimes_A^LG).         \tag{5.4.2}
\end{equation}
This morphism is not always an isomorphism, but it is one when $E$ is perfect. The examples intended for later use are $F=A$, variants of $A$ twisted by an endomorphism, and $F=M\otimes_KA^{\circ}$ with $M$ a dualizing complex.

\subsubsection*{5.5--5.7. Filtered and derived non-commutative traces}
Let $DF(A)$ be the derived category of finite filtered complexes of left $A$-modules. The preceding adjunction, evaluation, and tensor-product morphisms have filtered derived analogues, compatible with passage to associated graded objects.

Assume $A$ flat over $K$, and let $D(A)_{\mathrm{parf}}$ be the subcategory of complexes of finite perfect amplitude. For $E\in\ob D(A)_{\mathrm{parf}}$ and $F\in\ob D^{-}(A^{\circ})$, define
\begin{equation}
   \Tr_A:\R\Hom_A(E,F\otimes_A^LE)	o F\otimes_{A^e}^LA             \tag{5.6.1}
\end{equation}
by composing the inverse of (5.4.2), for $G=E$, with the evaluation map (5.3.5). On $H^0$ this gives
\begin{equation}
   \Tr_A:\Hom_A(E,F\otimes_AE)	o H^0(F)_A.                       \tag{5.6.2}
\end{equation}
For $F=A$ one obtains
\begin{align}
   \Tr_A:\R\Hom_A(E,E)&\to A_K,                                   \tag{5.6.3}\\
   \Tr_A:\Hom_A(E,E)&\to H^0(T,A_K).                              \tag{5.6.4}
\end{align}
The second map is local and elementary; the first is the global derived trace.

For filtered perfect complexes, the trace is compatible with the associated graded object. If $F$ is concentrated in filtration degree $0$ and $u$ is a morphism in $DF(A)$, then
\begin{equation}
   \Tr_A(u)=\sum_i\Tr_{A,\mathrm{gr}}(\operatorname{gr}^iu).        \tag{5.6.5}
\end{equation}
In particular, for $u:E\to F\otimes_A^LE$,
\begin{equation}
   \Tr(u[1])=-\Tr(u),                                              \tag{5.6.6}
\end{equation}
where $u[1]$ is the shifted morphism.

These traces simultaneously generalize the usual commutative traces of SGA 6 and the non-commutative traces of Stallings and Grothendieck. For $A=K$ they reduce to the standard trace maps. Moreover, for $u:E\to F$ and $v:F\to E$ between perfect $A$-complexes,
\begin{equation}
   \Tr_{A_K}(u\otimes v)=\Tr_A(vu)\in A_K.                         \tag{5.7.1}
\end{equation}
Homotopic endomorphisms have the same trace, because (5.6.4) is induced by a morphism of complexes.

\subsubsection*{5.8--5.10. Composition and restriction of scalars}
For perfect data $({}_AE,{}_BF_A,{}_BG)$, the trace maps are compatible with composition: there is a commutative square
\[
\begin{array}{ccc}
\R\Hom_A(E,F\otimes_AE)\otimes\R\Hom_B(F,G\otimes_BF) & \xrightarrow{\circ} & \R\Hom_A(E,G\otimes_BF\otimes_AE)\\
\downarrow \Tr_A\otimes\Tr_B && \downarrow \Tr_B\\
F_A\otimes G_B & \longrightarrow & (G\otimes_BF)_A.
\end{array}
\]
The analogous square with $A$ and $B$ interchanged is also commutative.

Let $A\to B$ be a morphism of $K$-algebras with $B$ flat over $A$. For $E\in\ob D^-(B)$, forgetting scalars gives
\begin{equation}
   \R\Hom_B(E,E)\to \R\Hom_A(E,E).                                 \tag{5.10.1}
\end{equation}
Together with (5.6.3) for $B$ and for $A$, this yields the commutative diagram
\[
\begin{array}{ccc}
\R\Hom_B(E,E) & \to & \R\Hom_A(E,E)\\
\downarrow && \downarrow\\
B_K & \to & A_K .
\end{array}
\]
Equivalently, for $u\in\Hom_B(E,E)$,
\begin{equation}
   \Tr_A(u)=\Tr_{B/A}\bigl(\Tr_B(u)\bigr),                         \tag{5.10.3}
\end{equation}
where $\Tr_{B/A}:B_K\to A_K$ is the trace associated with restriction of scalars.

\bigskip


\clearpage
\section*{SGA 5, Exposé III B after 5.10 through Exposé VI}
\addcontentsline{toc}{section}{SGA 5, Exposé III B after 5.10 through Exposé VI}

\bigskip

\section*{Exposé III B. Calculation of local terms -- continuation}

\subsection*{5.11. Twisted traces}

Let $A$ be a $K$-algebra, let $E\in\ob D^-(A)$, and let
$\rho:A\to A$ be an endomorphism of $K$-algebras. Denote by
$A_\rho$ the ring $A$, regarded as a left $A^e$-module through

\[
   a \otimes a' \longmapsto \rho(a)a'.
\]

There is a canonical isomorphism
\[
   \Hom_A(E,A_\rho\otimes_A E)\xrightarrow{\sim}\Hom_A(E,E),
\]
where the term on the right means the set of left $A$-module homomorphisms
$E\to E$ which are $\rho$-linear, that is, such that
\[
        u(ax)=\rho(a)u(x).
\]
For $u\in\Hom_A(E,A_\rho\otimes_A E)$ we shall sometimes write
$\Tr_A^\rho(u)$, or simply $\Tr_\rho(u)$, for the element
$\Tr_A(u)\in A_K$.

When $\rho$ is an inner automorphism, say $\rho(a)=gag^{-1}$, one checks that
$\Tr_\rho(u)$ depends only on the class of $g$ in $A_K^*$, and that for a
$\rho$-linear $u$ one has
\begin{equation}\tag{5.13.1}
   \Tr_\rho(u)=\text{the class of }g\text{ in }A_K\cdot \Tr(g^{-1}u),
\end{equation}
where $g^{-1}u$ is the $A$-linear endomorphism $x\mapsto g^{-1}u(x)$.

\subsection*{6. Equivariant correspondences and divisibility of local terms}

\subsubsection*{6.1. Notation}

We return to the notation of Part I and also fix a finite group $G$. We are given
$S$-schemes with a left action of $G$, denoted $\underline X$ and
$\underline Y$, together with stable closed subschemes
\[
  \underline A,\underline B\subset \underline X\times_S\underline Y.
\]
We write $X,Y,A,B$ for the underlying schemes, equivalently for the quotients by
$G$ in the sense used here. We assume that $A$ and $B$ satisfy the general-position
hypotheses of I, 1.2. We are also given objects
\[
      L_X\in\ob D(X),\qquad L_Y\in\ob D(Y),
\]
carrying actions of $G$ above the actions on $X$ and $Y$, and $G$-equivariant
homomorphisms
\[
      u:a_1^*L_X\longrightarrow a_2^!L_Y,
      \qquad
      v:b_2^*L_Y\longrightarrow b_1^!L_X.
\]

Then there are local terms
\[
      \langle u_z,v_z\rangle\in\Lambda
\]
as in I, 1.1.1. On the other hand, the global cup-product
$\langle \underline u,\underline v\rangle$ of I, 1.1.2 belongs to
\[
      H_c^4(X\times_S Y,K_{X\times_S Y})=\Lambda,
\]
and Theorem 1.2 applies to give
\[
      \langle u_z,v_z\rangle=\langle \underline u,\underline v\rangle.
\]

One may also consider the non-commutative traces of $u_z$ and $v_z$. More precisely,
$(L_X)_x$ and $(L_Y)_y$ are perfect complexes of $\Lambda[G]$-modules, and
$u_z,v_z$ are homomorphisms of $\Lambda[G]$-modules. Thus, by 5.6.3, one can form
\[
       \Tr_{\Lambda[G]}(u_zv_z)\in \Lambda[G]_{\Lambda},
\]
where $\Lambda[G]_{\Lambda}$ is the free $\Lambda$-module on the set $G/\sim$ of
conjugacy classes of $G$. Similarly the global cup-product gives an element
\[
       \langle\!\langle\underline u,\underline v\rangle\!\rangle
       \in H_c^4(X\times_S Y,\underline\Lambda)_{\Lambda[G]},
\]
where $\underline\Lambda$ denotes the constant sheaf with value
$\Lambda[G]_\Lambda$.

\begin{statement}{Theorem 6.2}
With the preceding notation, assume moreover that $A$ and $B$ are in general
position in the sense of I, 1.2. Then
\[
       \Tr_{\Lambda[G]}(u_zv_z)
       =\langle\!\langle\underline u,\underline v\rangle\!\rangle .
\]
In other words, the non-commutative local term coincides with the non-commutative
global term.
\end{statement}

The proof follows the same pattern as the proof of Theorem 1.2. By the same
reductions as in no. 2, one is first reduced to the case where $A$ and $B$ are
closed immersions; the matter is then to prove a vanishing analogous to (2.4.1).
The key lemma is the following.

\begin{statement}{Lemma 6.2.1}
Under the hypotheses of 6.2, suppose that $(L_X)|_U=0$ and $(L_Y)|_V=0$, with the
notation of 2.1. Then
\[
       \Tr_{\Lambda[G]}(u_zv_z)=\Tr_{\Lambda[G]}(\underline u\cdot\underline v)
\]
in $\Lambda[G]_{\Lambda}$.
\end{statement}

The proof of 6.2.1 is analogous to that of 2.3. One is reduced to the case in which
$L_X$ and $L_Y$ are concentrated at the closed points, and one applies the
non-commutative Lefschetz formula for the projections
$\{x\}\to S$ and $\{y\}\to S$.

To finish the proof of 6.2 one still has to establish the non-commutative analogue
of 2.5.

\begin{statement}{Proposition 6.2.2}
Under the hypotheses of 6.2, suppose that $a$ and $b$ are closed immersions and that
$(L_X)_x=0$, $(L_Y)_y=0$. Let $T$ be the strict localization of
$X\times_S Y$ at $z=(x,y)$, and identify $A$ and $B$ with closed subschemes of
$T$. Then the restriction map
\[
\Tr_{\Lambda[G]}:
R\Gamma_A\bigl(T,\RHom(DL_X\otimes L_Y,K_T)\bigr)
\longrightarrow
R\Gamma\bigl(B,\RHom(DL_X\otimes L_Y,K_T)|_B\bigr)
\]
is zero in $D(\Lambda[G]_{\Lambda})$.
\end{statement}

The proof occupies the rest of no. 6. As in no. 3, one is reduced to the tame case,
where the representations of inertia factor through $\ell$-groups. The reduction to
the tame case is analogous to 3.1--3.3. The remaining tame case is treated by means
of the purity theorem 4.3 and by an explicit calculation of local terms.

Theorem 6.2 has important consequences for the divisibility of local terms. Let
$\chi$ be a virtual representation of $G$ over $\Lambda$, i.e. an element of
$R_\Lambda(G)$. One may consider the local term $\chi(u_z,v_z)$ obtained by
applying $\chi$ to the non-commutative local term
$\Tr_{\Lambda[G]}(u_zv_z)$. More precisely, if
$\chi=\sum n_i[V_i]$, with the $V_i$ irreducible representations of $G$, put
\[
   \chi(u_z,v_z)=\sum n_i\Tr\bigl(u_zv_z\mid (L_X)_x\otimes_{\Lambda[G]}V_i\bigr).
\]
Then 6.2 implies the following.

\begin{statement}{Corollary 6.3}
With the notation of 6.2, let $\chi\in R_\Lambda(G)$, and write
\[
     e_\chi={1\over |G|}\sum_{g\in G}\chi(g^{-1})g
\]
for the central idempotent of $\Lambda[G]$ associated with $\chi$. Then
\[
     \chi(u_z,v_z)=e_\chi\cdot
       \langle\!\langle \underline u,\underline v\rangle\!\rangle .
\]
In particular, if $\chi$ is a representation of degree $1$, then
\[
     \chi(u_z,v_z)=\chi\bigl(\langle u_z,v_z\rangle^G\bigr),
\]
where $\langle u_z,v_z\rangle^G$ denotes the component of
$\Tr_{\Lambda[G]}(u_zv_z)$ on the trivial conjugacy class.
\end{statement}

When $G$ is cyclic, and one takes for $L_X$ and $L_Y$ the constant sheaves
$\Lambda[G]$, and for $u$ and $v$ the correspondences given by left and right
multiplication, Theorem 6.2 recovers Grothendieck's divisibility formula (XII, 4.5).
More generally, for every finite group $G$ and every virtual representation $\chi$,
formula 6.3 gives a canonical way of ``dividing'' ordinary local terms attached to
equivariant correspondences, as announced in the introduction.

\subsection*{Complementary bibliography}

\begin{enumerate}[label={[\arabic*]}]
\item H. Cartan and S. Eilenberg, \emph{Homological Algebra}, Princeton University Press, 1956.
\item A. Grothendieck, ``Formule de Lefschetz et rationalité des fonctions $L$'', Séminaire Bourbaki, exposé 279, 1964/65.
\item J.-L. Verdier, ``Catégories dérivées (état 0)'', in \emph{SGA $4\frac12$}, Lecture Notes in Mathematics 569, Springer, 1977.
\item R. P. Langlands, ``Modular forms and $\ell$-adic representations'', in \emph{Modular Functions of One Variable II}, Lecture Notes in Mathematics 349, Springer, 1973.
\item J. Stallings, ``Centerless groups -- an algebraic formulation of Gottlieb's theorem'', \emph{Topology} 4 (1965), 129--134.
\item M. Artin and J.-L. Verdier, ``Refinement of the Thomason--Trobaugh theorem'', manuscript.
\end{enumerate}

\newpage

\subsection*{6.7. Traces and change of scalars}

Let $X$ and $c$ be as in 6.5, and let $A\to B$ be a morphism of flat
$\Lambda$-algebras. With respect to extension and restriction of scalars, the maps
$\Tr_A$ of 6.5.3, 6.5.8, and 6.5.9 behave like the maps $\Tr_A$ of 5.6.1 and
5.6.2. The point is that, if $f$ is an $S$-morphism, the functors
\[
       f^*,\quad f_*,\quad f_!,\quad f^!
\]
commute with extension and restriction of scalars. Therefore the study of the
behavior of $\Tr_A$ in 6.5.3 with respect to extension and restriction of scalars is
reduced to the behavior of the tensor-product map
\[
\RHom_A(p_1^{*L},p_1^{*K A_X})\overset{L}{\otimes}_{A\,p_2^{*L}}
\longrightarrow
\RHom_A(p_1^{*L},p_1^{*K A_X})\overset{L}{\otimes}_{A\,p_2^{*L}}
\]
and of the evaluation map
\[
\RHom_A(L,K A_X)\overset{L}{\otimes}_A L\longrightarrow K A_{X,L_\natural}.
\]

Thus the commutativity of the squares 5.9.1 and 5.9.2 implies the commutativity of
\begin{equation}\tag{6.7.1}
\begin{array}{ccc}
\delta^{*c}{}_*\RHom_A(c_1^{*L},c_2^{!\,L}) & \xrightarrow{\Tr_A} & i_*^c K A_{X^c,L_\natural} \\[6pt]
\downarrow & & \downarrow \\[6pt]
\delta^{*c}{}_*\RHom_B(c_1^*(B\overset{L}{\otimes}_A L),c_2^{!}(B\overset{L}{\otimes}_A L)) & \xrightarrow{\Tr_B} & i_*^c K B_{X^c,L_\natural} .
\end{array}
\end{equation}
Here the left vertical map is extension of scalars and the right vertical map is
defined by the canonical morphism
\[
      A\overset{L}{\otimes}_{A^e}A\longrightarrow
      B\overset{L}{\otimes}_{B^e}B.
\]
It follows that, for $u\in\Hom_A(c_1^{*L},c_2^{!\,L})$,
\begin{equation}\tag{6.7.2}
      \Tr_B(B\overset{L}{\otimes}_A u)=\varphi\Tr_A(u),
\end{equation}
where
\[
  \varphi:H^0(X^c,K A_{X^c,L_\natural})\longrightarrow
  H^0(X^c,K B_{X^c,L_\natural})
\]
is the canonical map.

Assume that $A$ is of finite type and flat as a $\Lambda$-module, that $B$ is of
finite type and flat as a left $A$-module, and that $A_\natural$ and $B_\natural$ are
flat over $\Lambda$. We then have $\Tr_{B/A}:B_\natural\to A_\natural$, and we again
write
\begin{equation}\tag{6.7.3}
     \Tr_{B/A}:K B_{X,\natural}\longrightarrow K A_{X,\natural}
\end{equation}
for the map obtained by applying $K_X\overset{L}{\otimes}-$, as in 6.5.2. It follows
from 5.10.6 and 5.10.9 that, for $L\in\ob D_{ctf}^{-}(X)$, the square
\begin{equation}\tag{6.7.4}
\begin{array}{ccc}
\delta^{*c}{}_*\RHom_B(c_1^{*L},c_2^{!\,L}) & \xrightarrow{\Tr_B} & i_*^c K B_{X^c,\natural} \\[6pt]
\downarrow & & \downarrow\scriptstyle \Tr_{B/A} \\[6pt]
\delta^{*c}{}_*\RHom_A(c_1^{*L},c_2^{!\,L}) & \xrightarrow{\Tr_A} & i_*^c K A_{X^c,\natural}
\end{array}
\end{equation}
is commutative, the left vertical map being restriction of scalars. Consequently, for
$u\in\Hom_B(c_1^{*L},c_2^{!\,L})$, one has in
$H^0(X^c,K A_{X^c,\natural})$
\begin{equation}\tag{6.7.5}
      \Tr_A(u)=\Tr_{B/A}(\Tr_B(u)).
\end{equation}

\subsection*{6.8. External products}

Let $A$ and $B$ be flat $\Lambda$-algebras, let $X$ and $Y$ be $S$-schemes, and put
$Z=X\times_S Y$. Let
\[
      c:C\to X^2,\qquad d:D\to Y^2
\]
be correspondences, and let
\[
      e=c\times_S d:E=C\times_S D\to Z^2
\]
be their product. Then $Z^e=X^c\times_S Y^d$. Put $R=A\otimes_\Lambda B$.
As in III, 5.3, one defines the external tensor product
\begin{equation}\tag{6.8.1}
\Hom_A(c_1^{*L},c_2^{!\,L})\otimes
\Hom_B(d_1^{*M},d_2^{!\,M})
\xrightarrow{\overset{L}{\otimes}_S}
\Hom_R(e_1^{*}(L\overset{L}{\otimes}_SM),e_2^{!}(L\overset{L}{\otimes}_SM)).
\end{equation}
From the commutativity of 5.12.1 and 5.12.3 it follows easily that, for
$u\in\Hom_A(c_1^{*L},c_2^{!\,L})$ and
$v\in\Hom_B(d_1^{*M},d_2^{!\,M})$, one has in
$H^0(Z^e,K R_{Z^e,L_\natural})$
\begin{equation}\tag{6.8.2}
      \Tr_R(u\overset{L}{\otimes}_S v)=\Tr_A(u)\,\Tr_B(v),
\end{equation}
where the product on the right is defined by the canonical map
\begin{equation}\tag{6.8.3}
H^0(X^c,K A_{X^c,L_\natural})\otimes
H^0(Y^d,K B_{Y^d,L_\natural})
\longrightarrow
H^0(Z^e,K R_{Z^e,L_\natural}).
\end{equation}

\subsection*{6.9. $G$-schemes and equivariant cohomological correspondences}

Let $G$ be a finite group. If $X$ is a $G$-$S$-scheme, that is, an $S$-scheme endowed
with an action of $G$ by $S$-automorphisms, and if $A$ is a $\Lambda$-algebra, we
write $D(G,A^X)$ for the derived category of sheaves of left $G$-$A$-modules on
$X$ for the étale topology. For the notions of $G$-sheaves of sets, rings, modules,
etc., we refer to XV, 1. We use the subscripts $c$ and $ctf$ for the full
subcategories consisting of complexes which, as complexes of $A$-modules, have
constructible cohomology, respectively finite Tor-dimension and constructible
cohomology.

Similarly, if $B$ is a $\Lambda$-algebra, we write $D(G,A_B^X)$ for the derived
category of $G$-$(A,B)$-bimodules on $X$. If $G$ acts trivially on $X$, then a
left $G$-$A$-module on $X$ is just a left $A[G]$-module on $X$ without operators.
Thus $D(G,A^X)=D(A[G]^X)$. Since $G$ is finite one also has
$D_c(A^X)=D_c(A[G]^X)$, but only an inclusion
\[
     D_{ctf}(A[G]^X)\subset D_{ctf}(G,A^X).
\]

The duality formalism of SGA 4, XVII--XVIII extends without difficulty to
$G$-$S$-schemes. Since the functor forgetting the $G$-action is faithful, it is enough
to extend naturally the definitions of the six operations
\[
   \overset{L}{\otimes},\quad \RHom,\quad f^*,\quad f_*,\quad f_!,\quad f^!
\]
and of the canonical morphisms considered there (base change, Künneth, duality,
etc.). This extension causes no difficulty; for the definition of $f_!$ and $f^!$ one
uses equivariant compactifications. We leave to the reader the paraphrase of the
definitions and results of 6.2--6.8 in the setting of $G$-$S$-schemes.

If, for $i=1,2$, $X_i$ is a $G$-$S$-scheme, if
$c:C\to X=X_1\times_S X_2$ is a $G$-$S$-morphism, and if
$L_i\in\ob D_{ctf}(G,A_{X_i}^{X_i})$, the elements of
\[
     \Hom_A(c_1^{*L_1},c_2^{!\,L_2})
\]
(where $\Hom_A$ means $\Hom$ in the relevant equivariant derived category) will be
called \emph{equivariant cohomological correspondences from $L_1$ to $L_2$ with
support in $c$}. Notice that, when $G$ acts trivially on the $X_i$ and
$L_i\in\ob D_{ctf}(A[G]^{X_i})$, the isomorphism 6.4.1 in $D(G,X)$ refines to
an isomorphism
\[
D_{A[G]}L_1\overset{L}{\otimes}_{S,A[G]}L_2
\xrightarrow{\sim}
\RHom_{A[G]}(p_1^{*L_1},p_2^{!\,L_2}).
\]
Consequently, in an equivariant situation of type 6.5, with $G$ acting trivially on
the spaces and with $L\in\ob D_{ctf}(A[G]^X)\subset D_{ctf}(G,A^X)$, one can, for
$u\in\Hom_{A[G]}(c_1^{*L},c_2^{!\,L})$, define a local term
\[
     \Tr_{A[G]}(u)\in H^0(X^c,K A[G]_{X^c,L_\natural}),
\]
finer in the evident sense than the local term
\[
     \Tr_A(u)\in H^0(G,X^c,K A_{X^c,L_\natural})
\]
obtained by using only the finite Tor-dimension of $L$ relative to $A$.

\subsection*{6.10. Values on conjugacy classes}

Let $X$ and $c$ be as in 6.5, and let $L\in\ob D_{ctf}(A[G]^X)$. For
$u\in\Hom_{A[G]}(c_1^{*L},c_2^{!\,L})$, 6.5.9 gives a local trace
\begin{equation}\tag{6.10.1}
\Tr_{A[G]}(u)\in H^0(X^c,K A[G]_{X^c,\natural})
 = H^0(X^c,K_{X^c}\otimes A[G_\natural]),
\end{equation}
where $G_\natural$ is the set of conjugacy classes of $G$. For every
$x\in G_\natural$ one may consider the value
$\Tr_{A[G]}(u)(x)\in H^0(X^c,K_{X^c})$, obtained by projecting onto the component
indexed by $x$. From 6.7.5 and 5.11.2 one obtains
\begin{equation}\tag{6.10.2}
    \Tr_A(u)=|G|\,\Tr_{A[G]}(u)(e),
\end{equation}
where $e$ denotes the identity element of $G$.

More generally, let $g\in G$, let $Z_g$ be the centralizer of $g$, and let
$gu\in\Hom_{A[Z_g]}(c_1^{*L},c_2^{!\,L})$ be the correspondence obtained by
composing $u$ with left multiplication by $g$ on $c_1^{*L}$. Then
$L\in\ob D_{ctf}(A[Z_g]^X)$, and 6.10.2 gives
\begin{equation}\tag{6.10.3}
    \Tr_A(gu)=|Z_g|\,\Tr_{A[Z_g]}(gu)(e).
\end{equation}
By the same argument as in 5.11.3, 6.7.4 gives
\begin{equation}\tag{6.10.4}
    \Tr_{A[Z_g]}(gu)(e)=\Tr_{A[G]}(u)(g^{-1}),
\end{equation}
and hence, if desired,
\begin{equation}\tag{6.10.5}
    \Tr_A(gu)=|Z_g|\,\Tr_{A[G]}(u)(g^{-1}).
\end{equation}

\subsection*{6.11--6.13. Direct images and $\ell$-adic passage to the limit}

Let $f:X\to Y$ be a proper $G$-$S$-morphism, and let
\begin{equation}\tag{6.11.1}
\begin{array}{ccc}
X\times_S X & \xleftarrow{\ c\ } & C \\[4pt]
\downarrow\scriptstyle f\times_S f & & \downarrow \\[6pt]
Y\times_S Y & \xleftarrow{\ d\ } & D
\end{array}
\end{equation}
be a commutative equivariant square of $G$-$S$-schemes, with $c$ and $d$ proper.
Assume that $G$ acts trivially on $Y$ and on $D$. Write
\begin{equation}\tag{6.11.2}
       f^{c/d}\text{, or }f^c:X^c\longrightarrow Y^d,
\end{equation}
for the morphism induced on schemes of fixed points. Let
$L\in\ob D_{ctf}(G,X)$ and let
$u\in\Hom(c_1^{*L},c_2^{!\,L})$ be an equivariant correspondence with support in
$c$. Its direct image is an equivariant correspondence
\[
     f_*u\in\Hom_{A[G]}(d_1^{*}f_*L,d_2^{!}f_*L).
\]
For $g\in G$, write $Z_g$ for the centralizer of $g$, and let
\begin{equation}\tag{6.11.3}
      gu\in\Hom((gc)_1^{*L},c_2^{!\,L})
\end{equation}
be the $Z_g$-equivariant correspondence with support in
\begin{equation}\tag{6.11.4}
      gc=(gc_1=c_1g,c_2):C\longrightarrow X\times_S X
\end{equation}
defined as the composite
\[
      g^*(c_1^{*L})\xrightarrow{\alpha_g}c_1^{*L}\xrightarrow{u}c_2^{!\,L},
\]
where $\alpha_g$ is the map giving the action of $G$ on $c_1^{*L}$. It is clear that
\begin{equation}\tag{6.11.5}
      f_*(gu)=g(f_*u).
\end{equation}

\begin{statement}{Proposition 6.12}
Under the hypotheses of 6.11, assume that
$f_*L\in\ob D_{ctf}(A[G]^Y)$. Then, for every $g\in G$,
\begin{equation}\tag{6.12.1}
 f_*^{gc}\Tr_A(gu)=|Z_g|\,\Tr_{A[G]}(f_*u)(g^{-1}),
\end{equation}
where
\[
      f_*^{gc}:H^0(X^{gc},K_{X^{gc}})\longrightarrow H^0(Y^d,K_{Y^d})
\]
is the Gysin morphism defined by the adjunction map
$f_*^{gc}K_{X^{gc}}=f_*^{gc}f^{gc!}K_{Y^d}\to K_{Y^d}$.
\end{statement}

Indeed, applying 6.10.5 to $f_*u$ and using 6.11.5 gives
\[
     \Tr_A(f_*(gu))=|Z_g|\,\Tr_{A[G]}(f_*u)(g^{-1}),
\]
and 6.12.1 follows from the Lefschetz formula III, 4.5.1.

Let now $\ell\ne p$ be a prime. Put $\Lambda_n=\bZ/\ell^n\bZ$ for $n\ge1$.
Under the hypotheses of 6.11, let $L=(L_n)$ be a projective system of objects of
$D_{ctf}(G,X,\Lambda_n)$ such that
\[
      L_m\overset{L}{\otimes}_{\Lambda_m}\Lambda_n\xrightarrow{\sim}L_n
      \qquad (m\ge n),
\]
and let $u=(u_n)$ be a projective system of equivariant correspondences
$u_n\in\Hom(c_1^{*L_n},c_2^{!\,L_n})$, compatible under extension of scalars. By
compatibility of local traces with extension of scalars, the local terms
\[
      \Tr_{\Lambda_n}(gu_n)\in H^0(X^{gc},K\Lambda_{n,X^{gc}})
\]
define an element
\[
      \Tr_{\bZ_\ell}(gu)\in H^0(X^{gc},K_{\ell,X^{gc}})
      :=\varprojlim_n H^0(X^{gc},K\Lambda_{n,X^{gc}}),
\]
and the terms $f_*^{gc}\Tr_{\Lambda_n}(gu_n)=\Tr_{\Lambda_n}(g f_*u_n)$ define an
element
\[
      f_*\Tr_{\bZ_\ell}(gu)\in H^0(Y^d,K_{\ell,Y^d})
      :=\varprojlim_n H^0(Y^d,K\Lambda_{n,Y^d}).
\]
Assume moreover that $f_*M_n\in\ob D_{ctf}(\Lambda_n[G]^Y)$ for all $n$. Then the
local traces $\Tr_{\Lambda_n[G]}(f_*u_n)$ likewise define
\[
      \Tr_{\bZ_\ell[G]}(f_*u)
      \in H^0(Y^d,K\bZ_\ell[G]_{Y^d,\natural})
      :=\varprojlim_n H^0(Y^d,K\Lambda_n[G]_{Y^d,\natural}).
\]
It follows from 6.12 that
\begin{equation}\tag{6.13.1}
      f_*^{gc}\Tr_{\bZ_\ell}(gu)
      =|Z_g|\,\Tr_{\bZ_\ell[G]}(f_*u)(g^{-1}).
\end{equation}

\subsection*{6.14--6.17. Open immersions and divisibility}

Let $X$ be a $G$-$S$-scheme, let $j:U\to X$ be an equivariant open immersion, and let
$c:C\to X\times_S X$ be a morphism of $G$-$S$-schemes such that
$c_2=p_2c:C\to X$ is quasi-finite and of finite Tor-dimension. Assume that
$c_1^{-1}(U)\subset c_2^{-1}(U)$. The trace morphism
\[
      \Tr_{c_2}:c_{2!}\Lambda_C\longrightarrow \Lambda_X
\]
defines an equivariant filtered correspondence
\begin{equation}\tag{6.14.2}
      \underline c\in\Hom(c_1^*\Lambda_X,c_2^{!}\Lambda_X),
\end{equation}
where $\Lambda_X$ is endowed with the filtration defined by the submodule
$i_!\Lambda_U$, namely $F^n\Lambda_X=\Lambda_X$ for $n\le0$,
$F^1\Lambda_X=i_!\Lambda_U$, and $F^n\Lambda_X=0$ for $n\ge1$. Its associated
graded object consists of two equivariant correspondences
\begin{align}
\gr^0\underline c&=\underline c_{X-U,X}
   \in\Hom(c_1^*\Lambda_{X-U,X},c_2^{!}\Lambda_{X-U,X}),\tag{6.14.3}\
\gr^1\underline c&=\underline c_{U,X}
   \in\Hom(c_1^*\Lambda_{U,X},c_2^{!}\Lambda_{X-U,X}),
\end{align}
where $\Lambda_{X-U,X}=j_*\Lambda_{X-U}$ and
$\Lambda_{U,X}=i_!\Lambda_U$. Moreover
\begin{equation}\tag{6.14.4}
        \underline c_{X-U,U}=j_*\underline c_{X-U},
\end{equation}
where $\underline c_{X-U}$ is the correspondence defined by
$\Tr_{c_2,X-U}$.

Forgetting the action of $G$, one has local terms
\begin{equation}\tag{6.14.5}
\Tr_\Lambda(\underline c),\quad
\Tr_\Lambda(\underline c_{U,X}),\quad
\Tr_\Lambda(\underline c_{X-U,X})\in H^0(X^c,K_{X^c}),
\end{equation}
which, by additivity of traces (III, 4.13.1), satisfy
\begin{equation}\tag{6.14.6}
\Tr_\Lambda(\underline c)=\Tr_\Lambda(\underline c_{U,X})
+\Tr_\Lambda(\underline c_{X-U,X}).
\end{equation}
Furthermore, by 6.14.4,
\begin{equation}\tag{6.14.7}
\Tr_\Lambda(\underline c_{X-U,X})=j_*^c\Tr_\Lambda(\underline c_{X-U}),
\end{equation}
where $j_*^c$ is the direct-image morphism attached to
$j^c:(X-U)^c\to X^c$.

Recall that if $X$ is smooth, $c$ is a closed immersion, and $X^c$ is finite, then
$\Tr_\Lambda(\underline c)$ is the image in $\Lambda$ of the intersection multiplicity
$(C\cdot X)$ of $C$ with the diagonal. If moreover $X-U$ is smooth, then
$\Tr_\Lambda(\underline c_{X-U})$ is the image in $\Lambda$ of
$(c_2^{-1}(X-U)\cdot (X-U))$. Also, if $X-U$ is proper over $S$ and
$c|_{X-U}$ is the diagonal, the Lefschetz formula III, 4.7 gives that
\[
     \int_{X-U}\Tr(\underline c_{X-U})\in\Lambda
\]
is the Euler-Poincaré characteristic of $X-U$ with values in $\Lambda$.

Now let a diagram 6.11.1 be given, with $f,c,d$ proper, and let
$U\to X$ be an equivariant open immersion. Assume that $c_2$ is quasi-finite and
of finite Tor-dimension, that $c_1^{-1}(U)\subset c_2^{-1}(U)$, that $G$ acts
trivially on $C$ and $D$, and that
\[
     f|_U:U\longrightarrow f(U)=V
\]
is a Galois étale covering with group $G$. This last hypothesis implies that
$Rf_*(\Lambda_{U,X})\in\ob D_{ctf}(A[G]^Y)$; indeed it is concentrated in degree
zero and is the extension by zero of a sheaf on $V$ locally isomorphic to
$\Lambda[G]$. Thus one can apply 6.12 to
$L=\Lambda_{U,X}$ and $u=\underline c_{U,X}$, and, using 6.14.6 and 6.14.7, one
gets:

\begin{statement}{Proposition 6.16}
Under the hypotheses of 6.15, for every $g\in G$,
\begin{equation}\tag{6.16.1}
 f_*^{gc}\Tr_\Lambda(\underline{gc})
 -(fj)_*^{gc}\Tr_\Lambda(\underline{gc}_{X-U})
 =|Z_g|\,\Tr_{A[G]}(f_*\underline c_{U,X})(g^{-1}),
\end{equation}
where $j:X-U\to X$ is the inclusion and $Z_g$ is the centralizer of $g$.
\end{statement}

Applying 6.16 to the projective system of correspondences
$\underline c_\ell=(\underline c^{\Lambda_n})$, with $\Lambda_n=\bZ/\ell^n$ and
$\ell\ne p$, and passing to the limit in 6.14.6 and 6.14.7, gives:

\begin{statement}{Corollary 6.17}
Under the hypotheses of 6.15, for every $g\in G$,
\begin{equation}\tag{6.17.1}
 f_*^{gc}\Tr_{\bZ_\ell}(\underline{gc}_\ell)
 -(fj)_*^{gc}\Tr_{\bZ_\ell}(\underline{gc}_{\ell,X-U})
 =|Z_g|\,\Tr_{\bZ_\ell[G]}(f_*\underline c_{\ell,U,X})(g^{-1}),
\end{equation}
with the notation of 6.16.
\end{statement}

Taking into account compatibility of local traces with localization and the remarks
made in 6.14, one sees that 6.17 gives an affirmative answer, in a more general
formulation, to Grothendieck's divisibility conjecture (XII, 4.5) in the case where
the endomorphism $\varphi$ of $G$ considered there is the identity.

In fact, inspired by 5.13, it is not difficult to formulate analogous divisibility
statements to 6.12, 6.13, 6.16, and 6.17 in the case where $c$ is not
$G$-equivariant but $c_2$ is, and $c_1$ satisfies
\[
       c_1(xg)=c_1(x)\varphi(g)
\]
for every point $x$ of $G$, with a fixed endomorphism $\varphi$ of $G$. The
formulae then read as above, with $Z_g$ replaced by the $\varphi$-centralizer of
$g^{-1}$ (5.13.5).

\subsection*{6.18--6.24. Invariants and the comparison with Grothendieck local terms}

Put
\begin{equation}\tag{6.18.1}
      N_G=\sum_{g\in G}g\in\Lambda[G].
\end{equation}
Let $X$ be an $S$-scheme. For $L\in\ob D(A[G]^X)$, left multiplication by $N_G$
induces a morphism in $D(X)$
\begin{equation}\tag{6.18.2}
      N_G:\Lambda\overset{L}{\otimes}_{A[G]}L
      \longrightarrow \RHom_{A[G]}(\Lambda,L),
\end{equation}
where $\Lambda$ is regarded as a $\Lambda[G]$-algebra through the natural
augmentation sending every $g$ to $1$.

\begin{statement}{Lemma 6.18.3}
If $L\in\ob D_{ctf}(A[G]^X)$, the morphism 6.18.2 is an isomorphism.
\end{statement}

By dévissage one is reduced to the case in which $L$ is perfect, then to
$L=\Lambda[G]$. The assertion then follows from the fact that $N_G$ is a basis of
the invariant module of the left $\Lambda[G]$-module $\Lambda[G]$.

\begin{statement}{Proposition 6.19}
In the situation of 6.11, let $M\in\ob D_{ctf}(Y)$, let
$v\in\Hom(d_1^{*M},d_2^{!\,M})$, and let
$a:M\to f_*L$ be a morphism in $D(A[G]^Y)$, where $M$ is regarded as an object of
$D(A[G]^Y)$ through the augmentation $\Lambda[G]\to\Lambda$. Assume that the square
\[
\begin{array}{ccc}
d_1^{*M} & \xrightarrow{v} & d_2^{!\,M} \\[6pt]
\downarrow\scriptstyle d_1^{*a} & & \downarrow\scriptstyle d_2^{!a} \\[6pt]
d_1^{*}f_*L & \xrightarrow{f_*u} & d_2^{!}f_*L
\end{array}
\]
is commutative. Assume moreover that $f_*M\in\ob D_{ctf}(A[G]^Y)$ and that $a$
induces an isomorphism
\[
       M\xrightarrow{\sim}\RHom_{A[G]}(\Lambda,f_*L).
\]
Then
\begin{equation}\tag{6.19.1}
       \Tr_\Lambda(v)=\sum_{g\in G_\natural}\Tr_{A[G]}(f_*u)(g^{-1}).
\end{equation}
\end{statement}

Indeed, by 6.18.3, composing $a$ with $N_G^{-1}$ gives an isomorphism
\[
       M\xrightarrow{\sim}\Lambda\overset{L}{\otimes}_{A[G]}f_*L
\]
under which $v$ is identified with $\Lambda\overset{L}{\otimes}_{A[G]}f_*u$. By
6.7.2, $\Tr_\Lambda(v)$ is the image of $\Tr_{A[G]}(f_*u)$ under the map induced
by the augmentation $\Lambda[G]\to\Lambda$. This map sends each conjugacy class
to $1$, which proves the formula.

Let $X_i$ for $i=1,2$ be $G$-$S$-schemes, put $X=X_1\times_S X_2$, and let
$c:C\to X$ be a $G$-$S$-morphism such that $c_2$ is quasi-finite and of finite
Tor-dimension. For $L_i\in\ob D^+(G,X_i)$, the trace morphism
$\Tr_{c_2}:c_{2!}c_2^{*L_2}\to L_2$ defines by adjunction a map, again denoted
\begin{equation}\tag{6.20.1}
      \Tr_{c_2}:c_2^{*L_2}\longrightarrow c_2^{!\,L_2},
\end{equation}
and hence every equivariant morphism
$u\in\Hom(c_1^{*L_1},c_2^{*L_2})$ determines an equivariant correspondence
\begin{equation}\tag{6.20.2}
      u^!=\Tr_{c_2}\circ u\in\Hom(c_1^{*L_1},c_2^{!\,L_2}).
\end{equation}

\begin{statement}{Proposition 6.21}
Under the hypotheses of 6.11, assume moreover that an open immersion
$i:V\hookrightarrow Y$ is given such that $(fc_1)^{-1}(V)$ is connected, that
$d_1^{-1}(V)\subset d_2^{-1}(V)$, and that $f$ induces a principal Galois covering
with group $G$,
\[
      U=f^{-1}(V)\xrightarrow{f_V}V.
\]
Assume also that $M=i_!E$ with $E\in\ob D(V)$, and that one has an isomorphism
\[
      \alpha:f_V^*E\xrightarrow{\sim}\Lambda_X\otimes_\Lambda E_0
\]
in $D(G,X)$, with $E_0\in\ob D(\Lambda[G])$ perfect over $\Lambda$. Let
$v\in\Hom(d_1^{*M},d_2^{*M})$ be as above, and let
$v_0\in\Hom_{\Lambda[G]}(E_0,E_0)$ be the endomorphism defined, via $\alpha$, by
$f^*v|_{c_1^{-1}(U)}$. If $\underline c_{U,X}$ denotes the correspondence of
6.14.3, then, with the notation 6.20.2,
\begin{equation}\tag{6.23.1}
\Tr_\Lambda(v^!)=
\sum_{g\in G_\natural}
\Tr_{A[G]}(f_*\underline c_{U,X})(g^{-1})\,\Tr_\Lambda(gv_0).
\end{equation}
\end{statement}

Because of the central character of $\Tr_\Lambda$, the right hand side does not
depend on the choice of $\alpha$. The projection formula
$f_*f^*\Lambda_Y\otimes M\xrightarrow{\sim}f_*f^*M$ and 6.15 show that $M$
satisfies the hypotheses needed above. The isomorphism $\alpha$ identifies
$(f^*v)^!$ with the external tensor product $\underline c_{U,X}\otimes v_0$.
Thus $f_*f^*M$ is identified with $f_*(\Lambda_{U,X})\otimes E_0$, and
$f_*(f^*v)^!$ with $f_*\underline c_{U,X}\otimes v_0$. For $g\in G$, the
$Z_g$-equivariant correspondence $g f_*(f^*v)^!$ is identified with
$g f_*\underline c_{U,X}\otimes g v_0$. Since 6.10.4 gives
\[
\Tr_{A[G]}(f_*(f^*v)^!)(g^{-1})=
\Tr_{A[Z_g]}(g f_*(f^*v)^!)(e),
\]
it is enough, using 6.22.1, to prove
\[
\Tr_{A[Z_g]}(g f_*\underline c_{U,X}\otimes gv_0)(e)=
\Tr_{A[G]}(f_*\underline c_{U,X})(g^{-1})\,\Tr_\Lambda(gv_0),
\]
or equivalently, by 6.10.4,
\begin{equation}\tag{6.23.2}
\Tr_{A[Z_g]}(g f_*\underline c_{U,X}\otimes gv_0)(e)=
\Tr_{A[Z_g]}(g f_*\underline c_{U,X})(e)\,\Tr_\Lambda(gv_0).
\end{equation}
Replacing $G$ by $Z_g$, one may suppose $g=e$, and then, returning to the definition
of traces, it suffices to establish the following lemma.

\begin{statement}{Lemma 6.23.3}
Let $P$ be a $G_V$-module locally a direct factor of a free module of finite type, put
$P'=i_!P$, and let $Q$ be a complex of $\Lambda[G]$-modules perfect over
$\Lambda$. Then $P\overset{L}{\otimes}_\Lambda Q$, with the diagonal action of $G$,
is perfect over $\Lambda[G]$, and the square
\[
\begin{array}{ccc}
\delta^*\RHom_{\Lambda[G]}(p_1^{*P'},p_2^{!P'})
\overset{L}{\otimes}\RHom_{\Lambda[G]}(Q,Q)&\xrightarrow{(1)}&K_Y\\[8pt]
\downarrow\scriptstyle(*) && \parallel\\[8pt]
\delta^*\RHom_{\Lambda[G]}(p_1^{*P'}\overset{L}{\otimes}Q,
 p_2^{!P'}\overset{L}{\otimes}Q)&\xrightarrow{(3)}&K_Y
\end{array}
\]
is commutative. Here $\delta:Y\to Y\times_S Y$ is the diagonal; (1) is the tensor
product of
\[
\delta^*\RHom_{\Lambda[G]}(p_1^{*P'},p_2^{!P'})\xrightarrow{6.5.3}
K\Lambda[G]_{Y,\natural}\xrightarrow{x\mapsto x(e)}K_Y
\]
and
\[
\RHom_{\Lambda[G]}(Q,Q)\xrightarrow{\Tr_\Lambda}\Lambda;
\]
$(*)$ is induced from the external tensor product by restriction along the diagonal
$\Lambda[G]\to\Lambda[G]\otimes\Lambda[G]$; and (3) is 6.5.3 followed by
$x\mapsto x(e)$.
\end{statement}

The first assertion is clear: one is reduced to $P=\Lambda[G]_V$, and then
$P\overset{L}{\otimes}_\Lambda Q$ is the tensor product of $P$ by $Q$ with trivial
$G$-action. For the second assertion, 6.5.4 and 6.5.5 give
\[
\delta^*\RHom_{\Lambda[G]}(p_1^{*P'},p_2^{!P'})
\xrightarrow{\sim}
D_{\Lambda[G]}(i_!P)\overset{L}{\otimes}_{\Lambda[G]}i_!P
\xrightarrow{\sim}
i_!(D_{\Lambda[G]}P\overset{L}{\otimes}_{\Lambda[G]}P),
\]
so one is reduced to checking the induced square on $V$; thus one may assume
$V=Y$. The square becomes
\[
\begin{array}{ccc}
\RHom_{\Lambda[G]}(P,K_Y\overset{L}{\otimes}P)
\overset{L}{\otimes}\RHom_{\Lambda[G]}(Q,Q)&\xrightarrow{(1)}&K_Y\\[8pt]
\downarrow && \parallel\\[8pt]
\RHom_{\Lambda[G]}(P\overset{L}{\otimes}Q,
K_Y\overset{L}{\otimes}P\overset{L}{\otimes}Q)&\xrightarrow{(3)}&K_Y.
\end{array}
\]
Here (1) is given by $\Tr_{\Lambda[G]}(-)(e)\otimes\Tr_\Lambda$, and (3) by
$\Tr_{\Lambda[G]}(-)(e)$. By descent one is reduced to
$P=\Lambda[G]_Y$, and then one may suppose that $G$ acts trivially on $Q$. The
assertion is then an immediate consequence of 5.12.5.

If $\Lambda$ is killed by a power of a prime $\ell\ne p$, the right hand side of
6.23.1 may be rewritten as
\[
\sum_{g\in G_\natural}
\Tr_{\bZ_\ell[G]}(f_*\underline c_{\ell,U,X})(g^{-1})\,\Tr_\Lambda(gv_0),
\]
with the notation of 6.17. Formula 6.17.1 shows that Grothendieck's local terms in
XII, 6.2 coincide with Verdier's local terms, at least when the endomorphism
$\varphi$ of $G$ considered there is the identity; as already observed, this
restriction is not serious. If $Y$ is proper over $S$, combining 6.23.1 with the
Lefschetz formula for $Y\to S$ (III, 4.7) gives a formula generalizing XII, 6.3,
and hence in particular the Lefschetz formula for Frobenius proved in SGA
$4\frac12$, Rapport, and in Part I of the present exposé.

\subsection*{Bibliography}

\begin{enumerate}[label={[\arabic*]}]
\item H. Cartan and S. Eilenberg, \emph{Homological Algebra}, Princeton University Press, 1956.
\item A. Grothendieck, \emph{Formule de Lefschetz et rationalité des fonctions $L$}, Séminaire Bourbaki 1964/65, no. 279, in \emph{Dix exposés sur la cohomologie des schémas}, North Holland, 1968.
\item L. Illusie, \emph{Complexe cotangent et déformations I, II}, Lecture Notes in Mathematics 239 and 283, Springer, 1971--1972.
\item R. P. Langlands, \emph{Modular forms and $\ell$-adic representations}, in \emph{Modular Functions of One Variable II}, Lecture Notes in Mathematics 349, Springer, 1973.
\item J. Stallings, \emph{Centerless groups -- an algebraic formulation of Gottlieb's theorem}, Topology 4 (1965), 129--134.
\item J.-L. Verdier, \emph{The Lefschetz fixed point formula in étale cohomology}, in \emph{Proceedings of a Conference on Local Fields}, Springer, 1967.
\end{enumerate}

SGA $4\frac12$: \emph{Cohomologie étale}, by P. Deligne, Lecture Notes in Mathematics
569, Springer, 1977. In SGA $4\frac12$: \emph{Th. Finitude} refers to the finiteness
theorems in $\ell$-adic cohomology; \emph{Cycle} to the cohomology class associated
with a cycle; and \emph{Rapport} to the report on the trace formula.

\newpage

\section*{Exposé V. $J$-adic projective systems}
\begin{center}
\large J. P. Jouanolou
\end{center}

\subsection*{1. Generalities on abelian $A$-categories}

\subsubsection*{1.1. Definitions and first properties}

Let $A$ be a commutative ring. An \emph{$A$-category} $\cC$ is a category for which,
for each pair $(X,Y)$ of objects, the set $\Hom(X,Y)$ is endowed with an
$A$-module structure, compatibly with composition. Thus for three objects
$X,Y,Z$ the composition map
\[
       \Hom(Y,Z)\times\Hom(X,Y)\longrightarrow\Hom(X,Z)
\]
is $A$-bilinear.

An \emph{abelian $A$-category} is an additive $A$-category in which every morphism
has a kernel and a cokernel and in which every morphism factors canonically as an
epimorphism followed by a monomorphism. More precisely, one assumes an equivalence
between $\cC$ and a full subcategory of the category of sheaves of $A$-modules on a
site, so that finite inverse and direct limits, kernels, cokernels, images, and
coimages are representable and are computed sheafwise.

An \emph{$A$-functor}, respectively an additive $A$-functor, between $A$-categories
is a functor, respectively an additive functor, whose action on Hom-sets is
$A$-linear. An \emph{$(A,B)$-functor} from an $A$-category to a $B$-category consists
of a ring homomorphism $A\to B$ and a $B$-functor to the target. When $A=B$ and the
map is the identity, this is simply an $A$-functor.

Let $\cC$ be an abelian $A$-category. An object $M$ of $\cC$ is said to be of
\emph{finite type} (respectively \emph{noetherian}) if, for every filtered inductive
system $(N_i)_{i\in I}$ of subobjects of $M$, the canonical morphism
\[
       \varinjlim_i N_i\longrightarrow M
\]
has the expected finiteness behavior: in the finite-type case the union stabilizes
as soon as it is all of $M$, and in the noetherian case every increasing sequence of
subobjects is stationary. An object is \emph{coherent} if it is of finite type and if,
for every map $u:N\to M$ with $N$ of finite type, $\Ker(u)$ is of finite type. For
ordinary modules these are the usual notions.

If $J$ is an ideal of $A$, the \emph{canonical $J$-filtration} of an object $M$ is the
decreasing filtration by the images $J^nM$ of
$J^n\otimes_A M\to M$. An object is called \emph{$J$-adic} if this filtration is
$J$-adic: for every $n$, the canonical map
$J\otimes_A J^nM\to J^{n+1}M$ is an epimorphism and the transition maps
$M/J^{n+1}M\to M/J^nM$ are epimorphisms.

Throughout the exposé, the ideal $J$ is assumed to satisfy the Artin--Rees condition
(AR) of 1.1.8(ii). We shall use without further comment the elementary facts that
the full subcategory of noetherian objects is thick and abelian: it is stable under
subobjects, quotients, and extensions.

\subsubsection*{1.2--1.3. Projective systems and AR-zero systems}

Let $\Hom(\bN^\circ,\cC)$ denote the category of projective systems indexed by
$\bN$ with values in $\cC$. A system $X=(X_n)$ is called \emph{AR-zero} if for every
$n$ there is an integer $m\ge n$ such that the transition map
$X_m\to X_n$ is zero. The AR-zero systems form a thick subcategory, and in fact a
localizing subcategory, of $\Hom(\bN^\circ,\cC)$ in the sense of Gabriel.

The quotient category
\[
     \Hom(\bN^\circ,\cC)/\Hom(\bN^\circ,\cC)_{\AR\text{-zero}}
\]
is therefore abelian, and the canonical functor is exact. If $\cC'$ is a full thick
abelian subcategory of $\cC$, one writes
$\Hom_{\AR}(\bN^\circ,\cC')$ for the full subcategory of systems which are
AR-isomorphic to systems with values in $\cC'$. The usual properties of quotient
categories then give the corresponding thick abelian subcategory in the quotient.

\subsubsection*{1.4. Examples and constructions}

For an object $M$ of $\cC$, let $\kappa(M)$ denote the constant projective system
with value $M$. The functor
\[
       \kappa:\cC\longrightarrow \Hom(\bN^\circ,\cC)
\]
is fully faithful and exact.

For an integer $r\ge0$, the translation functor $[r]$ on projective systems is
defined by
\[
      (X[r])_n=X_{n+r}.
\]
It is exact, and there is a canonical monomorphism $X[r]\to X$ whose cokernel is
AR-zero.

If $\cC$ is complete, the projective limit
\[
      \varprojlim X=\varprojlim_n X_n
\]
is an exact-left functor. For ordinary countable systems one has the standard exact
sequence involving $R^1\varprojlim$, and the higher derived functors vanish under
the hypotheses used below.

The separated $J$-adic completion of an object $M$ is
\[
      \widehat M=\varprojlim_n M/J^nM
\]
when it exists. If $\cC$ is complete and countable products are exact, completion is
exact on noetherian objects. A projective system $X$ is \emph{strict $J$-adic} if
$X_n$ is killed by $J^{n+1}$, the transition maps are epimorphisms, and
$X_m/J^{n+1}X_m\xrightarrow{\sim}X_n$ for $m\ge n$. Such a system is noetherian,
respectively artinian, if its components are noetherian, respectively artinian.

The following elementary lemmas will be used repeatedly. A noetherian object has all
subobjects noetherian. If $M$ is separated and complete for the $J$-adic topology,
then the system $(M/J^nM)_n$ is strict $J$-adic, and its image in the quotient by
AR-zero systems is isomorphic to the image of the constant system $\kappa(M)$.

\subsection*{2. Mittag-Leffler--Artin--Rees and calculus of fractions}

In this section $\cC$ is an abelian $A$-category and $J\subset A$ satisfies (AR).
A projective system $X=(X_n)$ satisfies the \emph{Mittag-Leffler--Artin--Rees
condition}, abbreviated (MLAR), if for every $n$ there is $m\ge n$ such that for
all $p\ge m$ the image of $X_p\to X_n$ is equal to the image of $X_m\to X_n$.
In other words, the family of images
\[
       \Imn(X_m\to X_n)_{m\ge n}
\]
is stationary. A stricter version asks in addition that this stable image be a
$J$-adic object of $\cC$.

A system satisfying (MLAR) is AR-isomorphic to a strict projective system whose
transition maps are epimorphisms. The condition is stable under short exact
sequences in the familiar way. A $J$-adic projective system satisfies (MLAR) if and
only if it is strict. Every noetherian projective system satisfies (MLAR).

Let $X$ be a strict $J$-adic system. Its kernel is
\[
      \Ker(X)=(\Ker(X_n\to X_0))_n.
\]
This is a strict subsystem, and the quotient is the constant system $\kappa(X_0)$.
If $Y$ is a strict subsystem of $X$, we say that $Y$ is \emph{$J$-adapted} to $X$ if
for every $n$,
\[
      Y_n=X_n\cap Y_0\quad\text{inside }X_0.
\]
For such a subsystem, $Y$ is $J$-adic exactly when the associated graded system is a
module over $\gr_J(A)$ in the evident sense.

The associated graded object of a strict $J$-adic system is
\[
      \gr(X)=\bigoplus_{n\ge0} X_n/X_{n+1},
\]
viewed as a graded module over $\gr_J(A)$. The system $X$ is noetherian,
respectively artinian, if and only if $\gr(X)$ is noetherian, respectively artinian,
as a graded $\gr_J(A)$-module.

Localizing subcategories in the sense of Gabriel will be used to express the
calculus of fractions. If $\cL$ is a localizing subcategory of an abelian category
$\cA$, a morphism is a $\cL$-equivalence if its kernel and cokernel belong to
$\cL$. The quotient $\cA/\cL$ is the localization with respect to these equivalences,
and the canonical functor $\cA\to\cA/\cL$ is exact and universal for exact functors
annihilating $\cL$.

Applied to the localizing subcategory of AR-zero systems, this gives the category in
which AR-isomorphisms are inverted. In particular, for a strict $J$-adic system $X$,
the functor sending a $J$-adapted strict subsystem $Y\subset X$ to its image in the
quotient is an equivalence from the category of such subsystems to the category of
subobjects of the image of $X$ in the quotient.

\subsection*{3. $J$-adic and AR-$J$-adic projective systems}

A projective system is called \emph{$J$-adic} if it is isomorphic to a strict
$J$-adic system. A system is called \emph{AR-$J$-adic} if it is AR-isomorphic to a
strict $J$-adic system. We write
\[
J\text{-}\adn(\cC),\quad J\text{-}\adf(\cC),\quad
\AR\text{-}J\text{-}\adn(\cC),\quad \AR\text{-}J\text{-}\adf(\cC)
\]
for the corresponding noetherian and artinian subcategories.

A $J$-adic system remains $J$-adic after any translation $[r]$. Moreover, every
AR-$J$-adic noetherian system is represented, in the quotient by AR-zero systems,
by a translate $Y[r]$ of an actual $J$-adic noetherian system $Y$.

Assume now that $\cC$ is complete and that countable products are exact. For any
projective system $X$, the functor $\varprojlim$ is left exact, and its derived
functors satisfy
\[
       R^q\varprojlim X=0\quad(q\ge2).
\]
If $X$ is AR-zero, then $R^q\varprojlim X=0$ for all $q\ge0$. If $X$ is
AR-$J$-adic noetherian, then
\[
       R^1\varprojlim X=0.
\]
Consequently $\varprojlim X$ is separated and complete for the $J$-adic topology,
and the canonical morphisms
\[
      X_n\longrightarrow \varprojlim X/J^{n+1}\varprojlim X
\]
are isomorphisms.

The $J$-adic completion functor
\[
      \Lambda_J(M)=\varprojlim_n M/J^nM
\]
is additive and left exact. Its left derived functors satisfy
\[
      L_i\Lambda_J(M)=0\qquad(i\ge2).
\]
For $M$ noetherian, one has $L_1\Lambda_J(M)=0$, and
$\Lambda_J(M)\cong\varprojlim_n M/J^nM$. If $\cC$ is complete with exact countable
products and $M$ is noetherian, the canonical map
\[
      \widehat A\otimes_A M\longrightarrow \Lambda_J(M)
\]
is an isomorphism; this is the familiar exactness of completion for finite modules
in this categorical setting.

\subsection*{4. Filtrations and gradings}

Let $X=(X_n)$ be a projective system. A filtration of $X$ is a compatible family of
decreasing filtrations $F^\bullet X_n$. It is called \emph{$J$-good} if, for every
$n$, the filtration on $X_n$ is $J$-good, i.e. for some $n_0$,
\[
       J^pF^{n_0}X_n=F^{n_0+p}X_n\qquad(p\ge0).
\]
The associated graded system is
\[
       \gr(X)=(\gr(X_n))_n,
       \qquad \gr(X_n)=\bigoplus_{p\ge0}F^pX_n/F^{p+1}X_n.
\]
If $X$ is strict $J$-adic and has a $J$-good filtration, then $\gr(X)$ is a strict
$\gr_J(A)$-adic projective system. If moreover $X$ is noetherian, so is
$\gr(X)$.

The Artin--Rees lemma takes the following form. Let $M$ be a noetherian object of
an abelian $A$-category and $N\subset M$ a subobject. Give $M$ its canonical
$J$-filtration. Then the induced filtration on $N$ is $J$-adic: for every $n$ there
exists $m$ such that
\[
        N\cap J^mM\subset J^nN.
\]
The proof is by applying (MLAR) to the projective system
$(N/(N\cap J^nM))_n$.

For a noetherian object $M$, the canonical map
\[
      J^n\otimes_A M\longrightarrow J^nM
\]
is an epimorphism. Moreover, the system $(J^nM)_n$ is AR-zero if and only if $M$
is killed by a power of $J$. If an AR-$J$-adic noetherian system carries a $J$-good
filtration, its associated graded object is AR-$\gr_J(A)$-adic and noetherian. In a
short exact sequence of $J$-adic noetherian systems with $J$-good filtrations, the
associated sequence of graded systems is exact. Finally, for a noetherian strict
$J$-adic system $X$, there exists $r\ge0$ such that, for every strict subsystem
$Y\subset X$, the graded object $\gr(Y[r])$ is a graded $\gr_J(A)$-module.

The Nakayama lemma also has the expected form. If $M$ is separated and complete for
the $J$-adic topology and $M/JM=0$, then $M=0$. Hence a morphism between separated
and complete objects is an epimorphism if and only if it is so modulo $J$. Similarly,
$M$ is of finite type or noetherian if and only if $M/JM$ has the corresponding
property over $A/J$, with the usual finiteness condition on the kernels of
$J^n\otimes_A M\to J^nM$ in the noetherian criterion.

\subsection*{5. Noetherian AR-$J$-adic systems}

In this section $\cC$ is an abelian $A$-category, $J$ is an ideal satisfying (AR),
and $\cC$ has countable direct sums.

A projective system $X=(X_n)$ is \emph{AR-$J$-adic noetherian} (respectively
artinian) if it is AR-$J$-adic and its components are noetherian (respectively
artinian). We write
\[
     \AR\text{-}J\text{-}\adn(\cC),
     \qquad
     \AR\text{-}J\text{-}\adf(\cC)
\]
for the corresponding full subcategories of the AR quotient. The canonical functors
\[
      p_{\AR}^n:J\text{-}\adn(\cC)\to\AR\text{-}J\text{-}\adn(\cC),
      \qquad
      p_{\AR}^f:J\text{-}\adf(\cC)\to\AR\text{-}J\text{-}\adf(\cC)
\]
are equivalences of categories.

We shall need Hilbert's theorem in the following form. Let $A$ be a commutative
ring, $\cC$ an abelian $A$-category, and $M$ an object of finite type,
respectively noetherian. For every finitely generated $A$-algebra
$B=\bigoplus_{n\ge0}B_n$ graded by $\bN$, the graded $(B,\cC)$-module
$B\otimes_AM$ is of finite type, respectively noetherian. The proof reduces at once
to $B=A[T]$, where submodules of $A[T]\otimes_AM$ correspond to increasing arrays
of subobjects of $M$; the assertion is then precisely the finite-type or noetherian
condition on $M$.

If $\cC$ has countable direct sums, the same conclusion holds for an arbitrary
finitely generated $A$-algebra $B$, not necessarily graded. It follows in particular
that if $A$ is noetherian and $J\subset A$ is an ideal, then, for every noetherian
$J$-adic system $X$, the graded object $\gr(X)$ is a noetherian
$(\gr_J(A),\cC)$-module. Indeed $\gr_J(A)$ is a finitely generated $A$-algebra and
there is an epimorphism
\[
     \gr_J(A)\otimes_A X_0\longrightarrow \gr(X).
\]

Let $\cE_\cC$ be the full subcategory of $\underline{\Hom}(\bN^\circ,\cC)$ whose
objects are AR-$J$-adic noetherian systems.

\begin{statement}{Proposition 5.2.1}
The category $\cE_\cC$ is stable under kernels and cokernels in
$\underline{\Hom}(\bN^\circ,\cC)$. Thus $\cE_\cC$ is an abelian category and the
inclusion
\[
       \cE_\cC\longrightarrow \underline{\Hom}(\bN^\circ,\cC)
\]
is exact.
\end{statement}

Let $u:X\to Y$ be a morphism in $\cE_\cC$. One proves successively that
$\Coker(u)$, $\Imn(u)$, and $\Ker(u)$ are AR-$J$-adic and noetherian. For the
image one uses the following lemma.

\begin{statement}{Lemma 5.2.1.1}
Let
\[
       0\to L\to M\to N\to0
\]
be exact in $\underline{\Hom}(\bN^\circ,\cC)$. If $M$ and $N$ are AR-$J$-adic
noetherian, then $L$ is AR-$J$-adic noetherian.
\end{statement}

It is enough to prove the AR-$J$-adic assertion. One reduces as in the proof of
3.2.4 to the case where $M$ and $N$ are strict, and even $J$-adic noetherian. Then
3.1.3(ii) shows that $L$ is strict. It is adapted to $J$, and its graded object, as
a graded subobject of $\gr(M)$, is noetherian. The assertion follows from the
Artin--Rees lemma 4.2.6. Applying this first to
$0\to\Imn(u)\to Y\to\Coker(u)\to0$ and then to
$0\to\Ker(u)\to X\to\Imn(u)\to0$ gives the proposition.

\begin{statement}{Proposition 5.2.2}
The objects of the abelian category $\AR\text{-}J\text{-}\adn(\cC)$ are noetherian.
\end{statement}

It is enough to show that, for a $J$-adic noetherian object $X$, any increasing
sequence of subobjects in the AR category is stationary. Such a sequence may be
represented by morphisms
\[
      v_n:Y^n\longrightarrow X
\]
which are monomorphisms in the AR category and factor through one another there.
The assertion that $v_n$ factors through $v_{n+1}$ means that for some $r$ there is
a morphism $f:Y^n[r]\to Y^{n+1}$ with
\[
      v_n\circ a=v_{n+1}\circ f,
\]
where $a:Y^n[r]\to Y^n$ is the canonical strict morphism. Hence the image of
$v_n$ is contained in that of $v_{n+1}$. Replacing $Y^n$ by the image, one is
reduced to an increasing sequence of strict subsystems of $X$. Since $\gr(X)$ is
noetherian, the increasing sequence of the corresponding graded systems is
stationary, and so is the original sequence.

In summary: $\cE_\cC$ is a thick abelian subcategory of
$\underline{\Hom}(\bN^\circ,\cC)$; passage to the quotient by AR-zero systems
induces an equivalence
\[
\cE_\cC/(\cE_\cC\cap\underline{\Hom}(\bN^\circ,\cC)_{\AR\text{-zero}})
\xrightarrow{\sim}
\AR\text{-}J\text{-}\adn(\cC);
\]
and the objects of $\AR\text{-}J\text{-}\adn(\cC)$ are noetherian.

Theorem 5.2.3, stated in this language, says that both
$J\text{-}\adn(\cC)$ and $\AR\text{-}J\text{-}\adn(\cC)$ are abelian and noetherian
categories. For $\AR\text{-}J\text{-}\adn(\cC)$ this is precisely what has just been
proved; the category $J\text{-}\adn(\cC)$ is equivalent to it.

\begin{statement}{Proposition 5.2.4. Stability under extensions}
Let
\[
       0\to X\to Y\to Z\to0
\]
be exact in $\underline{\Hom}(\bN^\circ,\cC)$, with $J^{n+1}Y_n=0$ for all $n$.
Assume that $Z$ is AR-$J$-adic and that $X$ is AR-$J$-adic artinian. Then $Y$ is
AR-$J$-adic.
\end{statement}

As in the proof of 3.2.4, one reduces to the case in which $Y$ and $Z$ are strict
and $Z$ is $J$-adic. Then $X$ is strict. For some $r$ the projective system
\[
      (X_{n+r}/J^{n+1}X_{n+r})_{n\in\bN}
\]
is $J$-adic artinian. Since $X$ is a quotient of it and is AR-isomorphic to it, one
is reduced to the following lemma.

\begin{statement}{Lemma 5.2.4.1}
Let
\[
       0\to T\to X\to Y\to Z\to0
\]
be exact in $\underline{\Hom}(\bN^\circ,\cC)$, with $J^{n+1}Y_n=0$ for all $n$.
If $X$ is $J$-adic artinian, $Z$ is $J$-adic, and $T$ is AR-zero, then $Y$ is
AR-$J$-adic.
\end{statement}

For fixed $n$, the evident system
\[
       (Y_{n+k}/J^{n+1}Y_{n+k})_{k\in\bN}
\]
is essentially constant. Indeed, let $S_n^k$ be the kernel of
\[
       X_{n+k}/J^{n+1}X_{n+k}\longrightarrow
       Y_{n+k}/J^{n+1}Y_{n+k}.
\]
In the exact commutative diagram
\[
\begin{array}{ccccccccc}
0&\to&S_n^k&\to&X_{n+k}/J^{n+1}X_{n+k}&\to&Y_{n+k}/J^{n+1}Y_{n+k}&\to&Z_{n+k}/J^{n+1}Z_{n+k}\to0\\
 &&\downarrow&&\downarrow\wr&&\downarrow&&\downarrow\wr\\
0&\to&T_n&\to&X_n&\to&Y_n&\to&Z_n\to0,
\end{array}
\]
the $S_n^k$ are subobjects of the artinian object $T_n$ and form a decreasing
sequence. Hence they are eventually constant, which gives the essential constancy.

Define a system $Y'$ by
\[
       Y'_n=\varprojlim_r Y_{n+r}/J^{n+1}Y_{n+r}.
\]
Then $Y'$ is $J$-adic. Passing to the limit gives an exact commutative diagram
\[
\begin{array}{ccccccccc}
0&\to&S&\to&X&\to&Y'&\to&Z\to0\\
 &&\downarrow&&\downarrow\id&&\downarrow&&\downarrow\id\\
0&\to&T&\to&X&\to&Y&\to&Z\to0.
\end{array}
\]
Since $T$ is AR-zero, the snake lemma shows that $Y'\to Y$ is an AR-isomorphism.
Thus $Y$ is AR-$J$-adic.

\begin{statement}{Corollary 5.2.5}
Let
\[
       0\to X\to Y\to Z\to0
\]
be exact in $\underline{\Hom}(\bN^\circ,\cC)$, with $J^{n+1}Y_n=0$ for all $n$.
Assume that $X$ is AR-$J$-adic artinian, respectively artinian and noetherian, and
that $Z$ is AR-$J$-adic artinian, respectively noetherian. Then $Y$ is
AR-$J$-adic artinian, respectively noetherian.
\end{statement}

\subsection*{5.3. AR-$J$-adic noetherian systems and $\partial$-functors}

Let $\cC$ and $\cD$ be two abelian $A$-categories and let
$F:\cC\to\cD$ be an additive functor. It extends termwise to an additive functor
\[
       F:\underline{\Hom}(\bN^\circ,\cC)	o
       \underline{\Hom}(\bN^\circ,\cD)
\]
which commutes with translations.

\begin{statement}{Proposition 5.3.1}
Let $\cC$ and $\cD$ be abelian $A$-categories, and let
$T=(T^i)_{i\in\bN}$ be an exact $A$-linear $\partial$-functor from $\cC$ to $\cD$.
Assume:
\begin{enumerate}[label=(\roman*)]
\item for every noetherian object $M$ of $\cC$ killed by a power of $J$, and every
$n$, the object $T^n(M)$ is noetherian;
\item for every finitely generated graded $A$-algebra $B$ killed by $J$, and every
noetherian graded $(B,\cC)$-module $N$, the graded $(B,\cD)$-modules $T^n(N)$ are
noetherian.
\end{enumerate}
Then, for every AR-$J$-adic noetherian object $X$ of
$\underline{\Hom}(\bN^\circ,\cC)$, the objects $T^n(X)$ are AR-$J$-adic and
noetherian.
\end{statement}

Since the $T^n$ commute with translations, it suffices to treat the case where $X$
is $J$-adic noetherian. Then $\gr_J(X)$ is a noetherian graded
$\gr_J(A)$-module. By hypothesis (ii) and a variant of a corollary of Shih's theorem
(Appendix A.3), for each $n$ one has: $T^n(X)$ satisfies (MLAR), and
$\gr_J T^n(X)$ is a noetherian graded $(\gr_J(A),\cD)$-module. The Artin--Rees
lemma 4.2.6 completes the proof.

It follows that the $\partial$-functor
\[
        T=(T^i)_{i\in\bN}:
        \underline{\Hom}(\bN^\circ,\cC)	o
        \underline{\Hom}(\bN^\circ,\cD)
\]
induces an exact $A$-linear $\partial$-functor
\[
        \cE_\cC\longrightarrow\cE_\cD,
\]
commuting with translations, and hence, after passing to the quotient, an exact
$A$-linear $\partial$-functor, again denoted $(T^i)$,
\[
        \AR\text{-}J\text{-}\adn(\cC)\longrightarrow
        \AR\text{-}J\text{-}\adn(\cD).
\]

\subsection*{Appendix. Shih's theorem}

The purpose of this appendix is to prove Shih's theorem in the form used in the proof
of 5.3.1. Since this form is more general than the one stated in EGA $\mathcal O_{\mathrm{III}}$,
a direct proof, adapted from that of loc. cit., is required.

Fix an integer $r_0\ge1$. The data, for every $r\ge r_0$, of a bigraded object
$E_r=(E_r^{pq})_{p,q\in\bZ}$, endowed with a differential $d_r$ of bidegree
$(r,1-r)$, together with isomorphisms
\[
        H(E_r)\xrightarrow{\sim}E_{r+1},
\]
will be called a \emph{spectral filter}. The notion of morphism of spectral filters
is evident.

\subsubsection*{A.1. Spectral sequences attached to $\partial$-functors}

Let $\cC$ be an abelian category and let $X$ be an object of $\cC$ with a decreasing
filtration $(F^pX)_{p\in\bZ}$. For every exact $\partial$-functor from $\cC$ to an
abelian category $\cD$, the method of exact couples gives a spectral sequence
\[
        \mathbb E(X):E_1^{pq}=T^{p+q}(F^pX/F^{p+1}X)
        \Rightarrow E^n=T^n(X),
\]
whose filtration on $E^n$ is
\[
        F^p(E^n)=\Imn(T^n(F^pX)\to T^n(X))
        =\Ker(T^n(X)\to T^n(X/F^pX)).
\]
If the filtration on $X$ is finite, this spectral sequence is weakly convergent
(EGA $\mathcal O_{\mathrm{III}}$, 11.1.3).

The construction is functorial. If $f:X\to Y$ is a morphism of filtered objects,
there is a morphism of spectral sequences
\[
        \mathbb E(f):\mathbb E(X)\to\mathbb E(Y)
\]
which is $T^{p+q}(\gr^p f)$ on $E_1^{pq}$ and $T^n(f)$ on the abutment. If
$u:X\to X$ is an endomorphism with $u(F^pX)\subset F^{p+h}X$ for an integer $h$
independent of $p$, then $\mathbb E(u)$ is an endomorphism of bidegree $(h,-h)$ of
the spectral filter, commuting with the differentials.

\subsubsection*{A.2. The case of projective systems}

Let $X=(X_n)_{n\in\bZ}$ be a projective system of objects of $\cC$, with
$X_i=0$ for $i<0$. If $T$ is an exact $\partial$-functor from $\cC$ to $\cD$, it
extends naturally to a $\partial$-functor on
$\underline{\Hom}(\bZ^\circ,\cC)$. Giving $X$ its natural filtration as a projective
system yields a spectral sequence
\[
       \mathbb E(X)=\mathbb E(X_n)_{n\in\bZ},
\]
where $\mathbb E(X_n)$ is the spectral sequence associated with $T$ and the induced
filtration on $X_n$.

If $X$ is strict, the filtration of the abutment $E^n(X)$ in this spectral sequence
is the natural filtration of the projective system. Indeed,
\[
 F^p(E^n(X))=\Imn(T^n(F^pX)\to T^n(X)),
\]
and $(F^pX)_k=X_k$ for $k\ge p$ and $0$ for $k<p$.

Suppose a ring $S$ acts on $X$ compatibly with its natural filtration. Then, for each
$r\ge1$, the bigraded object $E_r(X)$ is canonically a graded
$(\gr^\bullet S,\cD)$-module; the same holds for
\[
       \cE_r^{(n)}(X)=(\cE_r^{p,q}(X))_{p+q=n}.
\]
Moreover $S$ acts on $T^n(X)$, and whenever $T^n(X)$ satisfies the Mittag-Leffler
condition, the strict graded object $\gr_S^\bullet T^n(X)$ is a graded
$(\gr^\bullet S,\cD)$-module.

\begin{statement}{Theorem A.3}
Suppose that, for an integer $n$, the graded $(\gr^\bullet S,\cD)$-modules
\[
      \cE_1^{(n)}(X)=T^n(\gr_S^\bullet X),
      \qquad
      \cE_1^{(n-1)}(X)=T^{n-1}(\gr_S^\bullet X)
\]
are noetherian. Then:
\begin{enumerate}[label=(\roman*)]
\item $T^n(X)$ satisfies (MLAR);
\item $\gr_S^\bullet T^n(X)$ is a noetherian graded
$(\gr^\bullet S,\cD)$-module.
\end{enumerate}
\end{statement}

For every pair $(m,r)$ one has
\[
      \cE_r^{(m)}(X)=(\cE_r^{p,m-p}(X))_{p\in\bN},
\]
a graded submodule of $\cE_1^{(m)}(X)$. Thus, for fixed $m$, the
$\cE_r^{(m)}(X)$ form an increasing sequence indexed by $r$. Since
$\cE_1^{(n)}(X)$ and $\cE_1^{(n-1)}(X)$ are noetherian, the corresponding sequences
are stationary. Hence for some $r_0$ and all $r\ge r_0$,
\[
 \cE_r^{(n)}(X)=\cE_{r_0}^{(n)}(X),
 \qquad
 \cE_r^{(n-1)}(X)=\cE_{r_0}^{(n-1)}(X).
\]
An equivalent characterization of the Mittag-Leffler condition then shows that the
spectral sequence $\mathbb E(X)$ is weakly convergent in degrees $n$ and $n-1$. It
follows that the filtration of $E^n(X)=T^n(X)$ is separated and that
$\cE_\infty^{(n)}(X)$ is a subobject of the noetherian module
$\cE_{r_0}^{(n)}(X)$.

To prove that $T^n(X)$ satisfies (ML), it is enough to prove that for each $p$ the
sequence
\[
\Imn(T^n(F^pX)\to T^n(X/F^{p+r}X))_{r\ge1}
\]
is stationary. Weak convergence identifies these images, for $r$ large, with the
kernels of maps to $E_\infty^{(n)}(X)$; since these kernels lie in a noetherian
object, they are stationary. Assertion (ii) follows because
$\gr_S^\bullet T^n(X)$ is a subquotient of $\cE_\infty^{(n)}(X)$.

\subsubsection*{A.4. Application to $\partial$-functors}

Let $\cC$ and $\cD$ be abelian $A$-categories, and let
$T=(T^i)_{i\in\bN}$ be an exact $A$-linear $\partial$-functor from $\cC$ to $\cD$.
Let $X$ be an object of $\underline{\Hom}(\bN^\circ,\cC)$, and let $S$ be a graded
$A$-algebra acting on $X$ compatibly with the transition maps. Assume that $X$ is
AR-$J$-adic noetherian; that $T^n(M)$ is noetherian for each noetherian $M$ killed
by a power of $J$; and that, for every finitely generated graded $A$-algebra $B$
killed by $J$ and every noetherian graded $(B,\cC)$-module $N$, the graded
$(B,\cD)$-modules $T^n(N)$ are noetherian.

\begin{statement}{Theorem A.4.1}
Under these hypotheses, the objects $T^n(X)$ are AR-$J$-adic and noetherian.
\end{statement}

By translating one is reduced to the case where $X$ is $J$-adic noetherian. Then
$\gr_J^\bullet(X)$ is a noetherian graded
$(\gr_J^\bullet(A),\cC)$-module. Hypothesis (iii) makes
$\cE_1^{(n)}(X)$ and $\cE_1^{(n-1)}(X)$ noetherian, and Theorem A.3 gives the
conclusion. This completes Exposé V.

\newpage

\section*{Exposé VI. $\ell$-adic cohomology}

\subsection*{1. Categories of constructible $\ell$-adic sheaves}

\subsubsection*{1.1. Generalities}

Let $X$ be a scheme and let $\ell$ be a prime number. We write
\[
       \bZ_\ell\text{-}\fc(X)
\]
for the category of constructible $\bZ_\ell$-sheaves on $X$, defined in SGA 5, V.
It is an abelian category. We write
\[
       \bQ_\ell\text{-}\fc(X)
\]
for the category of constructible $\bQ_\ell$-sheaves, i.e. the quotient of the
preceding category by the thick subcategory of torsion sheaves. We have also defined
the categories
\[
       \AR\text{-}\bZ_\ell\text{-}\fc(X),
       \qquad
       \AR\text{-}\bQ_\ell\text{-}\fc(X).
\]

Recall that giving a constructible $\bZ_\ell$-sheaf $\cF$ is equivalent to giving a
projective system $(\cF_n)_{n\in\bN}$ of constructible
$\bZ/\ell^{n+1}\bZ$-modules on $X$ such that, for every $n$, the transition morphism
$\cF_{n+1}\to\cF_n$ factors through an isomorphism
\[
       \cF_{n+1}/\ell^{n+1}\cF_{n+1}\xrightarrow{\sim}\cF_n.
\]
For two such sheaves $\cF=(\cF_n)$ and $\cG=(\cG_n)$,
\[
       \Hom(\cF,\cG)=\varprojlim_n\Hom(\cF_n,\cG_n).
\]
For each $n$, the canonical map
\[
       \Hom(\cF,\cG)\longrightarrow\Hom(\cF_n,\cG_n)
\]
is surjective, and its kernel consists of morphisms which factor through a torsion
object.

The ring $\bZ_\ell$ acts in the evident way on every component of a constructible
$\ell$-adic sheaf, and hence on the sheaf itself. Every morphism is compatible with
these actions. Thus the category of constructible $\ell$-adic sheaves is canonically
an abelian $\bZ_\ell$-category. We shall call these objects constructible
$\bZ_\ell$-sheaves. This should not be confused with a module over the sheaf of
rings $(\bZ_\ell)_X$, a notion not used here. The category will also be denoted
$\bZ_\ell\text{-}\fc(X)$.

\subsubsection*{1.2. Twisted constant constructible $\bZ_\ell$-sheaves}

A constructible $\bZ_\ell$-sheaf
\[
       \cF=(\cF_n)_{n\in\bN}
\]
is called \emph{twisted constant} if and only if all the components $\cF_n$ are
locally constant on $X_{\mathrm{et}}$. The category of twisted constant constructible
$\bZ_\ell$-sheaves is the full subcategory of $\bZ_\ell\text{-}\fc(X)$ generated by
these objects.

Assume that $\ell$ is invertible on $X$. For each $n$, raising to the $\ell$th power
gives an endomorphism of $\mathbb G_m$, and the sheaf
\[
       \bZ/\ell^{n+1}\bZ(1)
\]
is the sheaf $\mu_{\ell^{n+1}}$ of $\ell^{n+1}$st roots of unity. The projective
system of these sheaves defines a twisted constant $\bZ_\ell$-sheaf denoted
$\bZ_\ell(1)$.

Let $\cF=(\cF_n)$ be a twisted constant constructible $\bZ_\ell$-sheaf. For every
geometric point $\bar x$ of $X$, the geometric fiber $\cF_{\bar x}$ is a finite type
$\bZ_\ell$-module endowed with a continuous action of the fundamental group
$\pi_1(X,\bar x)$. The fiber functor
\[
       \cF\longmapsto\cF_{\bar x}
\]
is an equivalence between twisted constant constructible $\bZ_\ell$-sheaves on $X$
and finite type $\bZ_\ell$-modules with continuous action of $\pi_1(X,\bar x)$.

For a finite type $\bZ_\ell$-module $M$ we write $M_X$ for the corresponding twisted
constant sheaf. In particular,
\[
       \bZ_\ell(r)=\bZ_\ell(1)^{\otimes r}\qquad(r\in\bZ),
\]
where for $r<0$ one sets $\bZ_\ell(r)=\bZ_\ell(-r)^{\vee}$. If $\cF$ is a twisted
constant $\bZ_\ell$-sheaf, define
\[
       H^i(X,\cF)=\varprojlim_n H^i(X,\cF_n).
\]
These groups are finite type $\bZ_\ell$-modules when $X$ is proper over a
separably closed field.

\subsubsection*{1.3. Invertible $\ell$-adic sheaves}

Let $\cL$ be a constructible $\bZ_\ell$-sheaf. We say that $\cL$ is \emph{invertible}
if, for every $n$, the component $\cL_n$ is an invertible
$(\bZ/\ell^{n+1}\bZ)_X$-module, i.e. locally free of rank one. Equivalently,
$\cL$ is locally free in the sense of the preceding paragraph and the geometric
fibers of the $\cL_n$ are non-zero cyclic abelian groups.

For an invertible $\cL$ put
\[
      \cL^{-1}=\cH om(\cL,\bZ_\ell).
\]
There is then an isomorphism
\[
      \cL\otimes\cL^{-1}\xrightarrow{\sim}\bZ_\ell,
\]
obtained componentwise. Following EGA $\mathcal O_I$, 5.4.4, set
$\cL^{\otimes0}=\bZ_\ell$ and, for $n\ge1$,
$\cL^{\otimes(-n)}=(\cL^{-1})^{\otimes n}$. With this notation there are canonical
functorial isomorphisms
\[
      \cL^{\otimes m}\otimes\cL^{\otimes n}\xrightarrow{\sim}
      \cL^{\otimes(m+n)}
\]
for all integers $m,n$, satisfying the usual associativity rules. Thus
$(\cL^{\otimes i})_{i\in\bN}$ and $(\cL^{\otimes i})_{i\in\bZ}$ are graded
$\bZ_\ell$-algebras. The sheaf $\bZ_\ell(1)$ is invertible, and the convention above
gives $\bZ_\ell(r)=\bZ_\ell(1)^{\otimes r}$ for all $r\in\bZ$.

\subsubsection*{1.4. Categories built from constructible $\bZ_\ell$-sheaves}

Let $A$ be a finite $\bZ_\ell$-algebra. Define the category
\[
       A\text{-}\fc(X)
\]
of constructible $A$-sheaves as follows. An object is a pair $(\cF,u)$ consisting of
a constructible $\bZ_\ell$-sheaf $\cF$ and a homomorphism of $\bZ_\ell$-algebras
\[
       u:A\longrightarrow \End(\cF).
\]
A morphism $(\cF,u)\to(\cF',u')$ is a morphism $f:\cF\to\cF'$ of constructible
$\bZ_\ell$-sheaves such that, for every $a\in A$, the square
\[
\begin{array}{ccc}
\cF&\xrightarrow{f}&\cF'\\
\downarrow\scriptstyle u(a)&&\downarrow\scriptstyle u'(a)\\
\cF&\xrightarrow{f}&\cF'
\end{array}
\]
commutes. The category $A\text{-}\fc(X)$ is abelian, and the forgetful functor to
$\bZ_\ell\text{-}\fc(X)$ is exact and faithful. The dual of a constructible
$A$-sheaf $\cF$ is
\[
       \cF^\vee=\cH om(\cF,A_X).
\]
If $f:X\to Y$ is a morphism of schemes, inverse image and direct image extend in
the evident way to the corresponding categories of constructible $A$-sheaves.

Let $p:\bZ_\ell\text{-}\fc(X)\to\bQ_\ell\text{-}\fc(X)$ be the quotient functor. For a
constructible $\bZ_\ell$-sheaf $\cF$, the functor obtained by composing $p$ with
$\cG\mapsto\cF\otimes\cG$ makes multiplication by non-zero elements of
$\bZ_\ell$ invertible, and hence defines a right-exact functor
\[
       \bQ_\ell\text{-}\fc(X)\longrightarrow\bQ_\ell\text{-}\fc(X).
\]
Varying the first argument as well gives a right-exact bifunctor, called tensor
product,
\[
       \bQ_\ell\text{-}\fc(X)\times\bQ_\ell\text{-}\fc(X)
       \longrightarrow\bQ_\ell\text{-}\fc(X).
\]
If $X$ is connected and $\cF,\cG$ are twisted constant constructible
$\bQ_\ell$-sheaves corresponding to representations $V,W$ of $\pi_1(X)$, then
$\cF\otimes\cG$ is again locally constant and corresponds to the diagonal
representation on $V\otimes W$.

\subsubsection*{1.5. Constructible $\ell$-adic sheaves over the spectrum of a field}

If $k$ is separably closed, then
$\bZ_\ell\text{-}\fc(\Spec k)$ is identified with the category of finite type
$\bZ_\ell$-modules, and $\bQ_\ell\text{-}\fc(\Spec k)$ with finite-dimensional
$\bQ_\ell$-vector spaces.

If $k$ is a field and $\bar k$ is a separable closure, the geometric-fiber functor at
$\Spec(\bar k)\to\Spec(k)$ gives an equivalence
\[
\bZ_\ell\text{-}\fc(\Spec k)\xrightarrow{\sim}
\{\text{finite type }\bZ_\ell\text{-modules with continuous action of }
\Gal(\bar k/k)\}.
\]
After inverting torsion, this identifies $\bQ_\ell\text{-}\fc(\Spec k)$ with the
category of finite-dimensional continuous $\ell$-adic representations of
$\Gal(\bar k/k)$.

If $k$ is finite with $q$ elements, let
\[
        F_k\in\Gal(\bar k/k)
\]
denote the geometric Frobenius. A continuous $\ell$-adic representation of
$\Gal(\bar k/k)$ is determined by the action of $F_k$. If $V$ is a finite-dimensional
$\bQ_\ell$-vector space with continuous action of $\Gal(\bar k/k)$, the eigenvalues
of $F_k$ on $V$ are algebraic numbers whose inverses are integral over $\bZ$.

Let
\[
       \bZ_\ell\text{-}\mathrm{modn}(\pi_1)
\]
be the category of finite type $\bZ_\ell$-modules endowed with a continuous action of
$\pi_1$.

\begin{statement}{Lemma 1.2.4.2}
The preceding projective-limit functor
\[
\varprojlim_{\bN}:(\ell\text{-}\bZ)\text{-}\mathrm{ad}(\mathrm{Abf}(\pi_1))
\longrightarrow
\bZ_\ell\text{-}\mathrm{modn}(\pi_1)
\]
is an equivalence of categories.
\end{statement}

It is clearly essentially surjective: a finite type $\bZ_\ell$-module $E$ with a
continuous action of $\pi_1$ is obtained from the projective system
\[
       (E/\ell^{n+1}E)_{n\in\bN}
\]
with the evident actions of $\pi_1$.

\subsection*{2. Formalism of $\ell$-adic cohomology}

\subsubsection*{2.1. Inverse image}

Let $f:X\to Y$ be a morphism of locally noetherian preschemes. The inverse-image
functor
\[
      f^*:\mathrm{Ab}(Y)\longrightarrow \mathrm{Ab}(X)
\]
is exact and therefore extends in the evident way to an exact functor, denoted by the
same symbol,
\[
      f^*:\underline{\Hom}(\bN^\circ,\mathrm{Ab}(Y))
      \longrightarrow\underline{\Hom}(\bN^\circ,\mathrm{Ab}(X)).
\]
It plainly carries constructible $\bZ_\ell$-sheaves to constructible
$\bZ_\ell$-sheaves and locally AR-zero systems to locally AR-zero systems. With the
notation of 1.5.4 it therefore induces inverse-image functors
\[
\begin{array}{rcl}
 f^*:\bZ_\ell\text{-}\fc(Y)&\longrightarrow&\bZ_\ell\text{-}\fc(X),\\[3pt]
 f^*:\bQ_\ell\text{-}\fc(Y)&\longrightarrow&\bQ_\ell\text{-}\fc(X).
\end{array}
\]
These functors are exact, commute with tensor products, and preserve twisted
constant sheaves. If $A$ is a finite $\bZ_\ell$-algebra, the same construction gives
inverse-image functors for constructible $A$-sheaves.

\subsubsection*{2.2. Direct image}

Let again $f:X\to Y$ be a morphism of locally noetherian preschemes. The direct-image
functor on étale sheaves is left exact and so has right derived functors. Applying it
term by term to a projective system gives a left exact functor
\[
      f_*:\underline{\Hom}(\bN^\circ,\mathrm{Ab}(X))
      \longrightarrow\underline{\Hom}(\bN^\circ,\mathrm{Ab}(Y)).
\]
The corresponding derived functors are denoted $R^if_*$. When $f$ is proper and of
finite type, the finiteness theorem for étale cohomology shows that $R^if_*$ carries
constructible $\bZ_\ell$-sheaves to constructible $\bZ_\ell$-sheaves. Thus one obtains
functors
\[
      R^if_*:\bZ_\ell\text{-}\fc(X)\longrightarrow \bZ_\ell\text{-}\fc(Y)
\]
and, after passage to the quotient by torsion,
\[
      R^if_*:\bQ_\ell\text{-}\fc(X)\longrightarrow \bQ_\ell\text{-}\fc(Y).
\]
They are compatible with formation of fibers in the usual sense: at a geometric point
$\bar y$ of $Y$, the fiber of $R^if_*\cF$ is the $i$th cohomology of the geometric
fiber $X_{\bar y}$ with coefficients in the pullback of $\cF$.

\subsubsection*{2.3. Tensor products and internal Hom}

The tensor product of projective systems is defined level by level. If
$\cF=(\cF_n)$ and $\cG=(\cG_n)$ are constructible $\bZ_\ell$-sheaves, set
\[
      (\cF\otimes\cG)_n=\cF_n\otimes_{\bZ/\ell^{n+1}\bZ}\cG_n.
\]
This again defines a constructible $\bZ_\ell$-sheaf. The functor is right exact in
each variable; its left derived functors are denoted $\Tor_i^{\bZ_\ell}(\cF,\cG)$.
For constructible $\bQ_\ell$-sheaves the tensor product is exact.

The internal Hom is defined similarly by
\[
      \cH om(\cF,\cG)_n=
      \cH om_{\bZ/\ell^{n+1}\bZ}(\cF_n,\cG_n),
\]
with the transition morphisms induced in the evident way. Its right derived functors
are written $\mathcal Ext^i(\cF,\cG)$. The familiar adjunctions between tensor
product and internal Hom remain valid in the $\ell$-adic setting.

\subsubsection*{2.4. Cohomology}

For a locally noetherian scheme $X$ and a constructible $\bZ_\ell$-sheaf
$\cF=(\cF_n)$, define
\[
      H^i(X,\cF)=\varprojlim_n H^i(X,\cF_n).
\]
When $X$ is proper over a separably closed field, these groups are finite
$\bZ_\ell$-modules. Passing to $\bQ_\ell$ gives
\[
      H^i(X,\cF\otimes\bQ_\ell)
      =H^i(X,\cF)\otimes_{\bZ_\ell}\bQ_\ell.
\]
If $X$ is not proper, one uses cohomology with proper supports, defined by
compactification. For separated morphisms of finite type, the construction is
independent of the chosen compactification.

\subsubsection*{2.5. Finite cohomological dimension}

Let $X$ be of finite type over a separably closed field and let $\ell$ be invertible
on $X$. Then $X$ has finite $\ell$-cohomological dimension. More precisely, if
$\dim X=d$, then for every torsion constructible sheaf $F$ killed by a power of
$\ell$, one has
\[
       H^i(X,F)=0\qquad (i>2d),
\]
and similarly for cohomology with proper supports. This is the usual cohomological
dimension theorem of SGA 4, XIV.

\subsubsection*{2.6. Cohomology with proper supports}

Let $f:X\to S$ be a separated morphism of finite type between locally noetherian
schemes, and let $\cF$ be a constructible $\bZ_\ell$-sheaf on $X$. The cohomology
with proper supports
\[
      H_c^i(X/S,\cF)
\]
is the constructible $\bZ_\ell$-sheaf on $S$ obtained by compactifying $f$ and
applying the corresponding direct image with extension by zero. If
$j:U\to X$ is an open immersion and $i:Z\to X$ is the complementary closed
immersion, then for every constructible $\bZ_\ell$-sheaf $\cF$ on $U$ there is the
usual long exact localization sequence
\[
\cdots\to R^if_!(j_!\cF)\to R^if_!(Rj_*\cF)\to
R^if_!(i_*i^*Rj_*\cF)\to\cdots .
\]

\subsection*{3. $\ell$-adic $L$- and $Z$-functions}

\subsubsection*{3.1. Definitions}

Let $\bF_q$ be a finite field with $q$ elements and let $X$ be a scheme of finite
type over $\bF_q$. Let $|X|$ be the set of closed points of $X$. For
$x\in|X|$, write $d(x)$ for the degree of the residue field $k(x)$ over $\bF_q$, and
put $N(x)=q^{d(x)}$.

Let $\cF$ be a constructible $\bQ_\ell$-sheaf on $X$. For each closed point $x$, the
geometric fiber $\cF_{\bar x}$ is a finite-dimensional $\bQ_\ell$-vector space with
a continuous action of $\Gal(\overline{k(x)}/k(x))$. Let $F_x$ be geometric
Frobenius. If the eigenvalues of $F_x$ on $\cF_{\bar x}$ are
$\alpha_1,\ldots,\alpha_n$, then
\[
  \det(1-F_xt^{d(x)}\mid \cF_{\bar x})=\prod_{i=1}^n(1-\alpha_i t^{d(x)}).
\]
The associated $Z$-function is
\[
      Z(X,\cF;t)=\prod_{x\in|X|}
      {1\over\det(1-F_xt^{d(x)}\mid\cF_{\bar x})}\in\bQ_\ell[[t]].
\]
For $\cF=\bQ_\ell$ one writes simply $Z(X;t)$.

\subsubsection*{3.2. Lefschetz trace formula}

\begin{statement}{Theorem 3.2.1. Lefschetz trace formula}
Let $X$ be proper and smooth over $\bF_q$ and let $\cF$ be a constructible
$\bQ_\ell$-sheaf. Then
\[
Z(X,\cF;t)=\prod_{i=0}^{2\dim X}
\det(1-F^*t\mid H^i(X\otimes\overline{\bF}_q,\cF))^{(-1)^{i+1}}.
\]
\end{statement}

The proof uses the Lefschetz formula for correspondences on proper smooth varieties,
as developed in SGA 5, III, together with the Artin--Verdier theory of the
$L$-functor. The essential point is that the global trace of Frobenius on
cohomology is equal to the sum of its local terms at the fixed points, which are the
rational points over finite extensions of $\bF_q$.

\subsubsection*{3.3. The Weil conjectures}

Let $X$ be proper and smooth over a finite field $\bF_q$, and let $\ell$ be different
from the characteristic of $\bF_q$. The Weil conjectures assert:
\begin{enumerate}[label=(\roman*)]
\item the characteristic polynomials
\[
   \det(1-F^*t\mid H^i(X\otimes\overline{\bF}_q,\bQ_\ell))
\]
have integral coefficients and are independent of $\ell$;
\item the eigenvalues of $F^*$ on $H^i(X\otimes\overline{\bF}_q,\bQ_\ell)$ are
algebraic integers all of whose complex absolute values are $q^{i/2}$;
\item Poincaré duality gives perfect pairings
\[
H^i(X\otimes\overline{\bF}_q,\bQ_\ell)\times
H^{2d-i}(X\otimes\overline{\bF}_q,\bQ_\ell(d))\longrightarrow \bQ_\ell.
\]
\end{enumerate}
These conjectures are proved by Deligne. The proof uses the trace formula and the
theory of weights.

\subsection*{4. Cohomology of $\ell$-adic completions}

Let $X$ be proper over $\Spec(\bZ[1/\ell])$. For each $n\ge1$, put
\[
      X_n=X\otimes_{\bZ[1/\ell]}\bZ[1/\ell,\zeta_{\ell^n}],
\]
where $\zeta_{\ell^n}$ is a primitive $\ell^n$th root of unity. The cohomology groups
\[
      H^i(X_n,\bZ/\ell^n\bZ)
\]
form a projective system, and its limit
\[
      H^i_{\ell\text{-}\mathrm{comp}}(X)=
      \varprojlim_n H^i(X_n,\bZ/\ell^n\bZ)
\]
is called the $i$th completed $\ell$-adic cohomology group of $X$. These are finite
type $\bZ_\ell$-modules and coincide, in the proper smooth case, with the usual
$\ell$-adic cohomology groups. With coefficients in a constructible
$\bZ_\ell$-sheaf $\cF$, define similarly
\[
      H^i_{\ell\text{-}\mathrm{comp}}(X,\cF)=
      \varprojlim_n H^i(X_n,\cF/\ell^n\cF).
\]

There is also an Iwasawa spectral sequence. If
$G=\Gal(\bQ(\zeta_{\ell^\infty})/\bQ)\cong\bZ_\ell^\times$, then $G$ acts on
$H^i_{\ell\text{-}\mathrm{comp}}(X)$ by transport of structure. These groups may
therefore be regarded as modules over the Iwasawa algebra
\[
      \Lambda=\bZ_\ell[[G]]=
      \varprojlim_n\bZ_\ell[G/\ell^n\bZ_\ell^\times].
\]
The resulting $\Lambda$-modules $H^i_{\ell\text{-}\mathrm{comp}}(X)$ are of finite
type. The proof uses the Hochschild--Serre spectral sequence for the covers
$X_n\to X$ and the finiteness theorem for cohomology of proper schemes.

\subsection*{5. Applications}

\begin{statement}{Theorem 5.1.1. Rationality of $Z$-functions}
Let $X$ be of finite type over $\bF_q$ and let $\cF$ be a constructible
$\bQ_\ell$-sheaf on $X$. Then $Z(X,\cF;t)$ is a rational function of $t$.
\end{statement}

When $X$ is proper and smooth this follows from the Lefschetz trace formula 3.2.1.
The general case is deduced by dévissage, stratification, and the long exact
cohomology sequence.

\begin{statement}{Theorem 5.2.1. Unit theorem}
Let $K$ be a number field, let $\cO_K$ be its ring of integers, and let $\ell$ be a
prime. Then $\cO_K^\times$ is finitely generated, and its rank is
\[
       \operatorname{rank}(\cO_K^\times)=r_1+r_2-1,
\]
where $r_1$ is the number of real places and $r_2$ the number of complex places of
$K$.
\end{statement}

This classical theorem can be interpreted in terms of $\ell$-adic cohomology. The
Iwasawa spectral sequence for $\Spec(\cO_K[1/\ell])$ gives a description of the unit
group in terms of Iwasawa modules.

\begin{statement}{Theorem 5.3.1. Galois structure of cohomology}
Let $k$ be a number field and let $X$ be a smooth projective variety over $k$. Let
$\bar k$ be an algebraic closure of $k$. Then the representation of
$\Gal(\bar k/k)$ on
\[
      H^i_{\mathrm{et}}(X\otimes_k\bar k,\bQ_\ell)
\]
is semisimple, and its composition factors are independent of $\ell$ for $\ell$
different from the residual characteristics.
\end{statement}

The semisimplicity follows from Deligne's theory of weights. The independence of
$\ell$ is a consequence of the Weil conjectures and the theory of algebraic cycles.
This completes Exposé VI.

\bigskip
\begin{center}
\rule{0.45\textwidth}{0.4pt}\\[0.5em]
\end{center}


\clearpage
\section*{SGA 5, Exposé VII through Exposé XV}
\addcontentsline{toc}{section}{SGA 5, Exposé VII through Exposé XV}

\bigskip
\section*{Expos\'e VII -- Chern Classes, Self-Intersection, Flag Schemes, and Blow-Ups}
\addcontentsline{toc}{section}{Expos\'e VII -- continuation}

\noindent We resume in the middle of the proof of the projective bundle formula, at the continuation of 2.2.2, and continue through the end of the expos\'e.

\subsection*{2.2. Projective bundles: continuation}

Let $p:P\to S$ be the projective bundle associated with a locally free $\cO_S$-module $E$ of rank $r+1$, and let
\[
   \xi=c_1(\cO_P(1))\in H^2(P,\mu_n).
\]
The remaining point in the proof of the projective bundle theorem is to show that, for $0\le i\le r$, the class $\xi^i$ is a free element over the cohomology ring of $S$. It is enough to verify this for $\xi^r$, since the lower powers are then handled by the usual triangular argument. This is exactly the compatibility of the trace morphism with the fundamental class in the diagram
\[
\begin{array}{ccc}
H^{2r}(P,\mu_n^{\otimes r}) & \xrightarrow{\Tr} & A \\
\cup p^*(x)\uparrow && \uparrow x \\
H^0(P,\ZZ/n\ZZ) & \xrightarrow{\sim} & A,
\end{array}
\]
where $p^*$ is induced by a closed point of the fiber.

\begin{statement}{Corollary 2.2.4}
Let $L$ be a complex of $A$-modules. The morphism
\[
   R\Gamma_*(S,L)\longrightarrow R\Gamma_*(P,p^*L)
\]
induced by the canonical splitting of the projective bundle formula is an isomorphism after applying the corresponding projector.
\end{statement}

Indeed, it suffices to apply $R\Gamma_*$ to the isomorphism furnished by the projective bundle formula.

For later use set
\[
   A^*(X)=\bigoplus_i H^{2i}(X,\mu_n^{\otimes i}).
\]
This notation will be used systematically.

\begin{statement}{Corollary 2.2.6}
The pullback homomorphism
\[
   p^*:A^*(S)\longrightarrow A^*(P)
\]
is an injective homomorphism of unital rings, and $1,\xi,\ldots,\xi^r$ form a basis of the $A^*(S)$-module $A^*(P)$. In particular $A^*(P)$ is free over $A^*(S)$ of rank equal to the rank of $E$.
\end{statement}

This is the form of the projective bundle theorem used in the construction of Chern classes and in the splitting principle below.

\subsection*{3. Chern classes}

Let $S$ be a scheme and let $E$ be a locally free $\cO_S$-module of constant rank $r$. Put $P=\mathbb P(E)$ and let $\xi=c_1(\cO_P(1))$. The projective bundle formula shows that there is a unique relation
\[
   \xi^r-c_1\xi^{r-1}+c_2\xi^{r-2}-\cdots+(-1)^r c_r=0,
\]
with
\[
   c_i\in H^{2i}(S,\mu_n^{\otimes i}).
\]

If $E$ has non-constant rank, the scheme $S$ is partitioned into open-and-closed subsets $S_r$ on which $E$ has rank $r$. Since cohomology decomposes over this disjoint union, the preceding construction defines the classes on each component and hence on $S$.

\begin{statement}{Definition 3.2}
For a locally free $\cO_S$-module $E$, the element $c_i(E)\in H^{2i}(S,\mu_n^{\otimes i})$ obtained by the preceding construction is called the $i$th Chern class of $E$. If the rank of $E$ is bounded, the total Chern class is
\[
   c(E)=\sum_i c_i(E)\in A^*(S).
\]
\end{statement}

It is immediate from the definition that $c_0(E)=1$, and if $E$ has rank at most $d$, then $c_i(E)=0$ for $i>d$.

\begin{statement}{Proposition 3.4}
The Chern classes satisfy the following properties.
\begin{enumerate}[label=(\roman*)]
\item Normalization: if $L$ is invertible and $[L]$ is the image of the class of $L$ under the boundary map
\[
   H^1(S,\cO_S^*)\to H^2(S,\mu_n),
\]
then $c(L)=1+[L]$.
\item Base change: for $f:S'\to S$, one has $f^*c_i(E)=c_i(f^*E)$.
\item Additivity: for an exact sequence of locally free modules
\[
   0\to E'\to E\to E''\to0,
\]
one has
\[
   c_n(E)=\sum_{i+j=n}c_i(E')c_j(E'').
\]
\end{enumerate}
\end{statement}

\begin{prooftext}
For line bundles the projective bundle is $S$ itself and $\cO_P(1)$ is identified with the dual of the given line bundle, giving the normalization. Base change follows immediately from the functoriality of $\mathbb P(E)$, of $\cO(1)$, and of the defining relation.

For additivity one first reduces, by partitioning $S$ into open-and-closed pieces, to the case where all ranks are constant. Then one applies the splitting principle. Let $h:Z\to S$ be the product of the flag schemes of $E'$ and $E''$. The pullbacks $h^*E'$ and $h^*E''$ have finite filtrations with invertible quotients, say $L_i$ and $M_j$. The pullback $h^*E$ has a filtration with consecutive quotients the $L_i$ and $M_j$. The lemma below gives
\[
   c(h^*E)=\prod_i(1+[L_i])\prod_j(1+[M_j])
       =c(h^*E')c(h^*E'').
\]
The morphism $h$ is a composite of projective-bundle morphisms, so $h^*:A^*(S)\to A^*(Z)$ is injective by Corollary 2.2.6. The desired identity follows.
\end{prooftext}

\begin{statement}{Lemma 3.6}
Let $E$ be locally free and endowed with a finite filtration
\[
   0=E_0\subset E_1\subset\cdots\subset E_r=E
\]
whose successive quotients $L_i=E_i/E_{i-1}$ are invertible. Then
\[
   c(E)=\prod_{i=1}^r(1+[L_i]).
\]
\end{statement}

The proof is the standard flag argument. On $P=\mathbb P(E)$ the filtration defines closed subschemes
\[
   S\simeq Y_1\subset Y_2\subset\cdots\subset Y_r=P.
\]
The immersion $Y_i\hookrightarrow Y_{i+1}$ is defined by an invertible ideal identified with
\[
   p^*(L_i^\vee)\otimes\cO_{Y_{i+1}}(-1).
\]
Using the projection formula for the Gysin morphism and the identity $u_*(1_Z)=-[J]$ for an immersion defined by an invertible ideal $J$, one proves by induction the vanishing relation
\[
   (a_i)_*a_i^*\!iggl(\prod_{j\le i}p^*(L_j^\vee)\otimes\cO_P(-1)\biggr)=0.
\]
For $i=r$ this is precisely
\[
   \prod_{j=1}^r(\xi+c_1(p^*L_j))=0,
\]
which is the desired expression of the defining Chern polynomial.

\begin{statement}{Corollary 3.5}
For every exact sequence of locally free sheaves $0\to E'\to E\to E''\to0$,
\[
   c(E)=c(E')c(E'')
\]
in $A^*(S)$.
\end{statement}

\subsubsection*{3.8. Comparison with classical Chern classes}

If $X$ is locally of finite type over $\Spec\CC$, the analytic space $X(\CC)$ carries the usual topological theory of Chern classes. For a continuous complex vector bundle $V$ on $X(\CC)$ one has classes
\[
   c_i(V)\in H^{2i}(X(\CC),\ZZ).
\]
If $E$ is an algebraic locally free sheaf on $X$, its analytification gives such a bundle, and reduction modulo $n$ gives classes in
\[
   H^{2i}(X(\CC),\ZZ/n\ZZ).
\]
The comparison isomorphism between analytic and \etale cohomology identifies these with the \etale Chern classes just defined. In particular, the square
\[
\begin{array}{ccc}
H^1(X,\cO_X^*) & \longrightarrow & H^1(X(\CC),\cO_{X_{\mathrm{an}}}^*) \\
\partial\downarrow && \downarrow c_1 \\
H^2(X,\mu_n) & \longrightarrow & H^2(X(\CC),\ZZ/n\ZZ)
\end{array}
\]
commutes. This follows from the two Kummer-type exact sequences: algebraically,
\[
   0\to\mu_n\to\cO_X^*\xrightarrow{n}\cO_X^*\to0,
\]
and analytically,
\[
   0\to\ZZ\to\cO_{X_{\mathrm{an}}}\xrightarrow{\exp}\cO_{X_{\mathrm{an}}}^*\to0.
\]

Passing to the projective limit over $n=\ell^r$ defines the $\ell$-adic Chern classes
\[
   c_i(E)\in H^{2i}(X,\ZZ_\ell(i)),
\]
for every prime $\ell$ invertible on $X$.

\subsection*{4. The self-intersection formula and applications}

Let $i:Y\hookrightarrow X$ be a regular closed immersion of codimension $d$ between smooth schemes, and let $N=N_{Y/X}$ be the normal bundle. The cohomological Gysin morphism
\[
   i_*:H^m(Y,M)\to H^{m+2d}(X,M(d))
\]
is characterized by the fundamental class of $Y$ in $X$ and by compatibility with pullback, cup-product, and trace.

\begin{statement}{Self-intersection formula}
For $x\in H^*(Y,M)$ one has
\[
   i^*i_*(x)=c_d(N)\cup x.
\]
\end{statement}

The proof reduces by deformation to the normal bundle. In the vector-bundle case the zero-section $s:S\to N$ has normal bundle $N$, and the identity follows from the computation of the top Chern class through the projective completion
\[
   \widehat N=\mathbb P(N\oplus\cO_S).
\]
The projective bundle formula and the explicit description of the class of the zero-section give the equality.

One application is the Euler--Poincare class. If $X$ is projective over an algebraically closed field $k$, the diagonal class in $X\times X$ has a self-intersection whose degree is
\[
   \chi_\ell(X)=\sum_i (-1)^i\dim_{\QQ_\ell}H^i(X,\QQ_\ell).
\]
When $X$ is projective, this class is represented in the Chow ring by the self-intersection of the diagonal cycle. Hence it is independent of the auxiliary prime $\ell\ne\operatorname{char}(k)$; one may speak simply of the Euler--Poincare characteristic $\chi(X)$.

\subsection*{5. Flag schemes}

Let $E$ be a locally free $\cO_S$-module of rank $n$ and let
\[
   n=p_1+\cdots+p_k
\]
be a partition. The flag scheme $D=D(p_1,\ldots,p_k;E)$ represents filtrations of $E$ with successive quotients of ranks $p_i$. It is built by iterating Grassmannian or projective bundle constructions. Consequently the pullback
\[
   A^*(S)\longrightarrow A^*(D)
\]
is injective, and $A^*(D)$ is a finite free $A^*(S)$-module with an explicit basis in Schubert monomials in the Chern classes of the tautological quotients.

The expos\'e follows Grothendieck's treatment of Chern classes in the Chevalley seminar. The key lemma used in flag calculations says that, for an exact sequence
\[
   0\to L\to E\to F\to0
\]
with $L$ invertible, the induced sequence of symmetric algebras
\[
   0\to L\otimes S(E)\to S(E)\to S(F)\to0
\]
is exact as a sequence of graded $\cO_S$-modules. This identifies each step in the flag tower with a regular immersion defined by an invertible ideal, allowing the Gysin formula to be applied inductively.

\subsection*{6. Group schemes}

The same formalism applies to smooth group schemes and their homogeneous spaces. If a smooth group scheme $G$ acts transitively on a smooth scheme $X$ with stabilizer $H$, the quotient geometry gives a comparison between the cohomology of $X$, of $G$, and of $H$. When the quotient is built from flag varieties or from projective bundles, the projective-bundle theorem and Chern-class formulas compute its cohomology ring.

The point for later use is not merely the calculation itself but the functorial behavior: characteristic classes of the associated vector bundles are compatible with change of group, restriction to subgroups, and passage to quotient varieties. These compatibilities feed into the trace and Lefschetz formulas where equivariance is present.

\subsection*{7. Complete intersections}

Let $X$ be a smooth projective scheme and let $Y\subset X$ be a smooth complete intersection cut out by a regular sequence, or more generally by a regular section of a locally free sheaf. The normal bundle of $Y$ in $X$ is then identified with the restriction of that locally free sheaf. The self-intersection formula gives
\[
   i^*i_*(1)=c_d(N_{Y/X}).
\]

\begin{statement}{Lefschetz theorem 7.1}
For a smooth hyperplane section or suitable complete intersection, the pullback on cohomology is an isomorphism in the usual stable range and is injective at the middle degree boundary.
\end{statement}

The expos\'e treats this as the \etale analogue of the classical Lefschetz theorem. It also records the expected stronger Lefschetz statements suggested by the transcendental theory. The point is that Chern classes and Gysin morphisms supply the algebraic incarnation of the topological cycle classes used in the classical proof.

In particular, the Euler--Poincare characteristic of complete intersections is computed by applying Gauss--Bonnet to the total Chern class of the tangent bundle, with the normal exact sequence
\[
   0\to T_Y\to T_X|_Y\to N_{Y/X}\to0
\]
and the multiplicativity of total Chern classes.

\subsection*{8. Blown-up varieties}

Let $X$ be regular and let $Y\subset X$ be a regular closed subscheme of codimension $d$. Let
\[
   p:X'\to X
\]
be the blow-up of $X$ along $Y$, and let $Y'=p^{-1}(Y)$ be the exceptional divisor. Then
\[
   Y'\simeq \mathbb P(N_{Y/X}),
\]
where $N_{Y/X}$ is the normal bundle. Let $q:Y'\to Y$ be the projection and let $j:Y'\hookrightarrow X'$ be the exceptional divisor immersion.

The cohomology of $X'$ is determined by that of $X$ and $Y$. The additive decomposition is
\[
   H^*(X')\simeq p^*H^*(X)\oplus
      \bigoplus_{i=1}^{d-1} j_*\bigl(q^*H^{*-2i}(Y)(-i)\,\xi^{i-1}\bigr),
\]
where $\xi=c_1(\cO_{Y'}(1))$. The proof follows from the localization exact sequence for the pair $(X',Y')$, the corresponding sequence for $(X,Y)$, and the projective bundle formula for $Y'\to Y$.

The multiplicative structure is determined by three rules. First,
\[
   p^*(a)p^*(b)=p^*(ab).
\]
Second, the projection formula gives
\[
   p^*(a)\,j_*(z)=j_*(j^*p^*(a)\,z)=j_*(q^*i^*(a)\,z).
\]
Third, products of exceptional terms are reduced to the self-intersection formula for $j$:
\[
   j^*j_*(z)=c_1(\cO_{Y'}(-1))\,z=-\xi z.
\]
Together with the Chern polynomial relation of the projective bundle $\mathbb P(N_{Y/X})$, these rules give the full ring structure.

\subsection*{9. Chow ring of a blow-up and self-intersection in the Chow ring}

The last part of the expos\'e repeats the preceding calculation inside the Chow ring. Let $X$ be noetherian regular, let $Y\subset X$ be regular of codimension $d$, let $X'$ be the blow-up of $X$ along $Y$, and let $Y'$ be the exceptional divisor. Again $Y'\simeq\mathbb P(N_{Y/X})$.

\begin{statement}{Theorem 9.1}
With the preceding notation, the homomorphism
\[
   A^*(X)\oplus \bigoplus_{i=1}^{d-1}A^{*-i}(Y)\longrightarrow A^*(X')
\]
which sends $(a,b_1,\ldots,b_{d-1})$ to
\[
   p^*a+\sum_{i=1}^{d-1}j_*(q^*b_i\,\xi^{i-1})
\]
is an isomorphism of abelian groups.
\end{statement}

The proof is the Chow-theoretic version of the cohomological calculation. It uses the deformation to the normal cone and the projective bundle formula in Chow groups. The key computational lemma is the following.

\begin{statement}{Lemma 9.1}
Let $S$ be smooth and quasi-projective over a separably closed field, let $E$ be locally free of rank $r+1$, and let $p:P=\mathbb P(E)\to S$. Let $F$ be defined by
\[
   0\to F\to p^*E\to\cO_P(1)\to0,
\]
and put $\xi=c_1(\cO_P(1))$. Then, for $y\in C^*(S)$,
\[
   y=p_!\bigl(p^*y\,c_r(F)\bigr).
\]
If $E=N\oplus\cO_S$ and $s:S\to P$ is the zero-section in the projective completion $\widehat N$, then
\[
   s_!(y)=p^*y\,\xi^r,
   \qquad
   s^*s_!(y)=y\,c_r(N^\vee),
\]
and for $z\in C^*(\widehat N)$,
\[
   s^*(z)=p_!\bigl(z\,c_r(F)\bigr).
\]
\end{statement}

This lemma is a direct analogue of the cohomological computation in the normal bundle. It reduces the self-intersection formula to the projective bundle calculation and the behavior of the zero-section.

\begin{statement}{The self-intersection formula in the Chow ring}
Let $i:Y\hookrightarrow X$ be a regular closed immersion of codimension $d$, with normal bundle $N$. For $y\in A^*(Y)$ one has
\[
   i^*i_*(y)=c_d(N)\,y
\]
in $A^*(Y)$.
\end{statement}

This is Mumford's Chow-theoretic self-intersection formula. It is obtained by applying the preceding blow-up computation to the deformation to the normal cone, where the assertion becomes the explicit calculation for the zero-section.

\subsubsection*{9.3. Useful relations}

Inside the Chow ring of the blow-up, the exceptional divisor class and the Chern classes of the normal bundle satisfy the relation
\[
   \xi^d-c_1(N)\xi^{d-1}+\cdots+(-1)^d c_d(N)=0
\]
on $Y'=\mathbb P(N)$. The exceptional divisor satisfies
\[
   j^*j_*(1)=-\xi.
\]
Together with the projection formula, these relations determine all products involving exceptional classes.

\subsubsection*{9.4. Key lemma}

The key lemma asserts that the kernel and cokernel appearing in the comparison of $A^*(X')$ with the direct sum of $A^*(X)$ and the shifted $A^*(Y)$ are killed successively by projective bundle trace maps. More concretely, for the filtration by powers of the exceptional divisor, the associated graded pieces are the corresponding graded pieces of
\[
   A^*(Y')[1,\xi,\ldots,\xi^{d-1}],
\]
modulo the Chern polynomial. The projective bundle formula identifies these pieces with $A^*(Y)$.

\subsubsection*{9.5. End of the proof}

The proof of the blow-up theorem follows by combining the key lemma with the localization exact sequences
\[
   A^*(Y)\to A^*(X)\to A^*(X\setminus Y)\to0
\]
and
\[
   A^*(Y')\to A^*(X')\to A^*(X'\setminus Y')\to0,
\]
noting that $X'\setminus Y'$ is canonically isomorphic to $X\setminus Y$. The resulting diagram has exact rows. The projective bundle formula computes the left vertical term, and the five lemma gives the decomposition.

\subsubsection*{9.6. The key formula}

The fundamental identity may be written in the compact form
\[
   j_*\bigl(q^*(y)\xi^m\bigr)\cdot j_*\bigl(q^*(z)\xi^n\bigr)
     =-j_*\bigl(q^*(yz)\xi^{m+n+1}\bigr),
\]
with the understanding that powers $\xi^a$ for $a\ge d$ are reduced by the Chern polynomial of $N$. This formula is simply the projection formula and the equality $j^*j_*(1)=-\xi$.

\subsubsection*{9.8. Proof of Proposition 9.6}

Proposition 9.6 is proved by applying the preceding formula repeatedly and using the fact that $q^*:A^*(Y)\to A^*(Y')$ is injective. The exceptional contribution is expressed in the basis $1,\xi,\ldots,\xi^{d-1}$ over $A^*(Y)$; uniqueness of the expression follows from the projective bundle theorem. The compatibility with pushforward to $X$ comes from
\[
   p_*j_*=i_*q_*.
\]

\subsubsection*{9.9. Structure theorem for the Chow group of the blow-up}

The structure theorem states that $A^*(X')$ is the direct sum of the pullback of $A^*(X)$ and the exceptional summands generated by $j_*(q^*A^*(Y)\xi^i)$ for $0\le i\le d-2$. Multiplication is described by the formulas above. Equivalently, $A^*(X')$ is the $A^*(X)$-algebra generated by the exceptional divisor class, with relations induced by the Chern polynomial of $N$ and the restriction of cycles from $X$ to $Y$.

\subsubsection*{9.10. Interpretation and Riemann--Roch}

The final interpretation identifies the blow-up formula with the Riemann--Roch behavior of the exceptional divisor and the normal bundle. The Chern character and Todd class transform as dictated by the exact sequences
\[
   0\to T_{X'}\to p^*T_X\to \text{exceptional quotient}\to0
\]
and
\[
   0\to T_Y\to T_X|_Y\to N\to0.
\]
Thus the blow-up formula is compatible with the expected Riemann--Roch identity: pushforward of characteristic classes through $p$ recovers the classes on $X$, while the exceptional summands account for the correction terms supported on $Y$.

\subsection*{Bibliography of Expos\'e VII}

The sources cited include Grothendieck's theory of Chern classes, the Chevalley seminar on Chow rings and applications, Mumford's proof of the self-intersection formula in the Chow ring, and the Lefschetz theorems used in the cohomological applications.

\newpage

\section*{Expos\'e VIII -- Class Groups of Abelian and Triangulated Categories. Perfect Complexes}
\addcontentsline{toc}{section}{Expos\'e VIII -- class groups and perfect complexes}

\begin{center}
\textbf{Class groups of abelian and triangulated categories. Perfect complexes.}\\
\emph{by A. Grothendieck, written by I. Bucur}
\end{center}

This expos\'e recalls, in a rapid form, the notions and results needed later concerning Grothendieck groups, pseudo-coherent complexes, and perfect complexes. A more complete treatment, with detailed proofs, is announced for the first expos\'es of the seminar on the Riemann--Roch theorem for schemes, later SGA 6.

\subsection*{1. Abelian categories}

Let $\cC$ be an abelian category and $\cD$ a full subcategory of $\cC$. A function $f$ from $\cD$ to an abelian group $G$ is called \emph{additive} if, for every exact sequence
\[
   0\to E'\to E\to E''\to0
\]
in $\cC$ whose three terms belong to $\cD$, one has
\[
   f(E)=f(E')+f(E'').
\]
From the definition one immediately obtains $f(0)=0$, invariance under isomorphism, and the corresponding formula for exact sequences viewed in the opposite category. These elementary properties are the universal motivation for the Grothendieck group.

\subsection*{2. The Grothendieck group of an abelian category}

Let $\cC$ be an abelian category. Let $F(\cC)$ be the free abelian group on isomorphism classes $[E]$ of objects of $\cC$. The Grothendieck group $K(\cC)$ is the quotient of $F(\cC)$ by the subgroup generated by the elements
\[
   [E]-[E']-[E'']
\]
for all exact sequences $0\to E'\to E\to E''\to0$. The canonical map
\[
   k_\cC:\operatorname{Ob}\cC/\simeq\,\longrightarrow K(\cC)
\]
is universal for additive functions out of $\cC$: if $f$ is additive with values in an abelian group $G$, then there is a unique homomorphism $g:K(\cC)\to G$ such that $f=g\circ k_\cC$.

If $u:\cC\to\cC'$ is exact, then $u$ induces a homomorphism
\[
   K(u):K(\cC)\to K(\cC'),
\]
characterized by
\[
   K(u)([E])=[u(E)].
\]
This is the functoriality of the Grothendieck group in the abelian setting.

\subsection*{3. Triangulated categories}

Let $\cC$ be a triangulated category. The Grothendieck group $K(\cC)$ is defined as the free abelian group on isomorphism classes of objects, modulo the relations
\[
   [X]-[Y]+[Z]=0
\]
for every distinguished triangle
\[
   X\to Y\to Z\to X[1].
\]
Thus $[Y]=[X]+[Z]$. An exact functor of triangulated categories $u:\cC\to\cC'$ induces
\[
   K(u):K(\cC)\to K(\cC')
\]
by the same rule $[E]\mapsto [u(E)]$.

If $\cC=\D^b(\cA)$ is the bounded derived category of an abelian category, the canonical functor from $\cA$ to $\D^b(\cA)$ induces an isomorphism
\[
   K(\cA)\simeq K(\D^b(\cA)).
\]
The class of a bounded complex $K^\bullet$ is
\[
   [K^\bullet]=\sum_i(-1)^i[H^i(K^\bullet)],
\]
or equivalently the alternating sum of its terms whenever this makes sense.

\subsection*{4. Thick subcategories and localization}

A full subcategory $\cC'$ of a triangulated category $\cC$ is \emph{thick} if it is triangulated and closed under direct summands: any direct factor in $\cC$ of an object of $\cC'$ again belongs to $\cC'$. If $\cC'$ is thick, the quotient category
\[
   \cC/\cC'
\]
is triangulated and comes with an exact functor $Q:\cC\to\cC/\cC'$ which is universal among exact functors from $\cC$ annihilating $\cC'$.

There is a natural exact sequence
\[
   K(\cC')\longrightarrow K(\cC)\xrightarrow{K(Q)}K(\cC/\cC')\longrightarrow0.
\]
This is the categorical form of localization. It will later be used to pass from complexes with torsion, or with prescribed support, to quotient categories where those complexes vanish.

\subsection*{5. Pseudo-coherent complexes}

Let $\cC$ be an abelian category and let $\cC_0$ be a full subcategory satisfying stability conditions under kernels, cokernels, and the subquotients needed to perform homological constructions. We denote by
\[
   \D^\pm_{\cC_0}(\cC)
\]
the essential image of $\D^\pm(\cC_0)$ in $\D(\cC)$. Its objects are called $\cC_0$-pseudo-coherent complexes. Equivalently, a complex of $\cC$ is $\cC_0$-pseudo-coherent if it is isomorphic in $\D(\cC)$ to a complex bounded above whose terms lie in $\cC_0$.

This subcategory is triangulated. Indeed, if two terms in a distinguished triangle
\[
   U\to V\to W\to U[1]
\]
are pseudo-coherent, represent them by complexes $X^\bullet$ and $Y^\bullet$ with terms in $\cC_0$. After replacing $X^\bullet$ by a quasi-isomorphic complex, the morphism is represented by an actual morphism of complexes. Its mapping cone has terms in $\cC_0$ by the stability assumptions, and represents the third object $W$.

The intersection of $\D^\pm_{\cC_0}(\cC)$ with $\D^b(\cC)$ is denoted
\[
   \D^{\pm,b}_{\cC_0}(\cC).
\]
An object $M$ of $\cC$ is pseudo-coherent if it is pseudo-coherent as a complex concentrated in degree zero; equivalently, it admits a resolution on the left by objects of $\cC_0$.

\begin{statement}{Proposition 5.1}
A complex $X^\bullet$ is $\cC_0$-pseudo-coherent if and only if it is isomorphic to a bounded-above complex whose terms are $\cC_0$-pseudo-coherent objects.
\end{statement}

The proof is deferred to SGA 6, where the theory is treated systematically.

\subsection*{6. Perfect complexes}

An object of $\D(\cC)$ is called $\cC_0$-perfect if it is isomorphic to an object of $\D^b(\cC_0)$, that is, to a bounded complex all of whose terms lie in $\cC_0$. The full subcategory of perfect complexes is denoted
\[
   \D^{\mathrm{parf}}_{\cC_0}(\cC).
\]
It is the essential image of $\D^b(\cC_0)$ in $\D(\cC)$. It is a triangulated subcategory, stable under direct summands under the standard hypotheses. The Grothendieck group of perfect complexes is therefore defined by the preceding triangulated construction.

When $\cC$ is a category of modules and $\cC_0$ is the category of finite free or finite projective modules, this recovers the usual perfect complexes. In geometry, perfect complexes are precisely the complexes locally quasi-isomorphic to bounded complexes of finite locally free modules.

\subsection*{7. An important special case}

Let $A$ be a ring and let $\cC$ be the category of $A$-modules. Let $\cC_0$ be the category of finite projective $A$-modules. Then a complex is perfect if and only if it is locally, for the relevant topology, quasi-isomorphic to a bounded complex of finite projective modules. The resulting Grothendieck group is the $K$-group used in Riemann--Roch formulas.

This special case is central because cohomology complexes with constructible coefficients often become perfect after imposing finite Tor-dimension conditions. In that situation, the trace of an endomorphism is computed in the Grothendieck group and then evaluated by a rank, character, or local term.

\subsection*{8. Tor of complexes}

Let $\cC_\ell$ denote a category of modules over a coefficient ring such as $\ZZ_\ell$ or $\ZZ/\ell^n\ZZ$. A complex $X$ of $\D(\cC_\ell)$ has \emph{finite Tor-dimension} if it is bounded and if there exists an integer $n_0$ such that
\[
   \operatorname{Tor}^i(X,M)=0
\]
for every object $M$ and every $i>n_0$ after the usual shift convention. Equivalently, tensoring with $X$ has uniformly bounded cohomological amplitude.

Finite Tor-dimension is the condition which permits one to define derived tensor products without leaving bounded categories. It is also the condition which makes duality behave as expected: a perfect complex has finite Tor-dimension, and its dual is again perfect.

\subsection*{9. Reminders on linear representations of finite groups}

Let $G$ be a finite group and $A$ a commutative ring. The category of finite projective $A[G]$-modules has a Grothendieck group $K(A[G])$. If $A$ is a field of characteristic not dividing $|G|$, this group is controlled by characters. In the integral $\ell$-adic setting one uses projective $\ZZ_\ell[G]$-modules and compares them after extending scalars to $\QQ_\ell$.

Induction and restriction are exact functors and therefore induce maps on Grothendieck groups:
\[
   \operatorname{Ind}_H^G:K(A[H])\to K(A[G]),
   \qquad
   \operatorname{Res}_H^G:K(A[G])\to K(A[H]).
\]
Duality sends a module $M$ to
\[
   M^*=\Hom_A(M,A),
\]
with the contragredient action. If the character of $M$ is $\chi$, then that of $M^*$ is $g\mapsto\chi(g^{-1})$.

\subsection*{4. The Swan representation}

We now turn to the local representation attached to a ramified Galois cover of curves. Let $C$ be a smooth, proper, connected curve over an algebraically closed field $k$, with generic point $\eta$ and function field $K=k(\eta)$. Let
\[
   K\subset K',\qquad [K':K]=n
\]
be a finite Galois extension with group $G$, and let $C'$ be the normalization of $C$ in $K'$. The canonical map is
\[
   p:C'\to C.
\]
The group $G$ acts on $C'$. For $y'\in C'$ write $G_{y'}$ for its stabilizer.

If $\Pi$ is a uniformizer at $y'$ and $s\in G_{y'}$, $s\ne e$, define
\[
   v_{y'}(s)=v_{y'}(s(\Pi)-\Pi).
\]
This is the multiplicity of $y'$ as a fixed point of the automorphism $s:C'\to C'$. The equality follows from the local intersection calculation: the diagonal and the graph of $s$ meet with multiplicity $v_{y'}(s(\Pi)-\Pi)$.

\subsubsection*{4.2. Definition of the Swan representation}

The Swan character attached to $y'$ is the class function
\[
   \operatorname{sw}_{y'}(s)=
   \begin{cases}
      -1-v_{y'}(s), & s\ne e,\\[0.3em]
      \sum_{s\ne e}v_{y'}(s), & s=e.
   \end{cases}
\]
The theorem of Artin--Serre--Swan asserts that, for a prime $\ell\ne\operatorname{char}(k)$, there is a finite projective $\ZZ_\ell[G_{y'}]$-module $\Sw_{y'}$, unique up to isomorphism, having this character.

One also considers
\[
   \wt{\Sw}_{y'}=\Sw_{y'}\otimes_{\ZZ_\ell}\QQ_\ell.
\]
Its character is
\[
   \widetilde{\operatorname{sw}}_{y'}(s)=
   \begin{cases}
      -1, & s\ne e,\\[0.3em]
      |G_{y'}|-1, & s=e.
   \end{cases}
\]
Finally, the Artin character is
\[
   \operatorname{art}_{y'}=\operatorname{sw}_{y'}+\mathbf 1_{G_{y'}},
\]
where $\mathbf 1_{G_{y'}}$ denotes the character of the augmentation representation. Explicitly,
\[
   \operatorname{art}_{y'}(s)=
   \begin{cases}
      -v_{y'}(s), & s\ne e,\\[0.3em]
      \sum_{s\ne e}v_{y'}(s)+|G_{y'}|, & s=e.
   \end{cases}
\]
Let $y=p(y')$. By induction from $G_{y'}$ to $G$ set
\[
\begin{aligned}
   \Sw_y &= \Sw_{y'}\otimes_{\ZZ_\ell[G_{y'}]}\ZZ_\ell[G],\\
   \wt{\Sw}_y &= \wt{\Sw}_{y'}\otimes_{\ZZ_\ell[G_{y'}]}\ZZ_\ell[G],\\
   \Art_y &= \Art_{y'}\otimes_{\ZZ_\ell[G_{y'}]}\ZZ_\ell[G].
\end{aligned}
\]
Their characters are denoted $\operatorname{sw}_y$, $\widetilde{\operatorname{sw}}_y$, and $\operatorname{art}_y$.

\subsubsection*{4.3. Independence of the point above $y$}

The induced modules just defined depend, up to isomorphism, only on $y\in C$, not on the chosen point $y'$ above $y$. If $y''=g y'$, then the situation at $y''$ is obtained from that at $y'$ by the automorphism $g$ of $C'$ and by conjugation $s\mapsto gsg^{-1}$ on $G$. Extension of operators through this inner automorphism sends the corresponding module to an isomorphic one. Hence $\Sw_y$, $\wt{\Sw}_y$, and $\Art_y$ are well-defined.

\begin{statement}{Proposition 4.4}
With the preceding notation, and for $M^\bullet\in\D^*(\ZZ_\ell[G])$, one has canonical isomorphisms
\[
   \Sw_{y'}^*=\Sw_{y'},
\]
and
\[
   \Hom_{\ZZ_\ell[G]}(\Sw_{y'},M^\bullet)
      \simeq \Sw_{y'}\otimes_{\ZZ_\ell[G]}M^\bullet.
\]
\end{statement}

\begin{prooftext}
Finite projective $\ZZ_\ell[G]$-modules are determined by their characters in this setting. If $\sigma$ is the character of $\Sw_{y'}$, then the dual has character $\sigma^*(s)=\sigma(s^{-1})$. But
\[
   v_{y'}(s^{-1}(\Pi)-\Pi)
    =v_{y'}(s^{-1}(s(\Pi)-\Pi))
    =v_{y'}(s(\Pi)-\Pi),
\]
so the Swan character is invariant under inversion. This gives self-duality; the Hom--tensor formula follows.
\end{prooftext}

\subsubsection*{4.5. Explicit character formulas}

For later use the character of the induced Swan representation can be written explicitly. For $s\in G$,
\[
   \operatorname{sw}_y(s)=
   \sum_{p(y')=y}\operatorname{sw}_{y'}(s),
\]
with terms understood to be zero unless $s$ fixes the corresponding point. Thus
\[
   \operatorname{sw}_y(s)=
   \begin{cases}
      \sum_{p(y')=y}(1-v_{y'}(s)), & s\ne e,\\
      0, & s=e,
   \end{cases}
\]
under this convention. Similarly
\[
   \operatorname{art}_y(s)=
   \begin{cases}
      \sum_{p(y')=y}(-v_{y'}(s)), & s\ne e,\\
      0, & s=e,
   \end{cases}
\]
and
\[
   \widetilde{\operatorname{sw}}_y(s)=
   \begin{cases}
      \sum_{p(y')=y}(-1), & s\ne e,\\
      0, & s=e.
   \end{cases}
\]
These formulas are obtained from the character formula for induced representations, choosing representatives for the right cosets of $G_{y'}$ in $G$.

\subsubsection*{4.6. The non-connected case}

The invariants $\Sw_y$, $\wt{\Sw}_y$, and $\Art_y$ are defined under slightly more general hypotheses. Let $C'$ be a smooth algebraic curve over $k$, not necessarily connected, with an action of a finite group $G$. Suppose there is a dense open subset $U$ of $C=C'/G$ such that $C'|_U$ is a principal $G$-cover of $U$.

For a point $y\in C$, choose $y'\in C'$ above $y$. Let $C'_{y'}$ be the irreducible component of $C'$ containing $y'$, and $C_y$ the component of $C$ containing $y$. Let $G'_{y}$ be the stabilizer of the component $C'_{y'}$. Then the extension of function fields
\[
   k(\eta_y)\subset k(\eta'_{y})
\]
is Galois with group $G'_y$, and $C'_{y'}$ is the normalization of $C_y$ in that extension. Applying the connected construction to $C'_{y'}\to C_y$ gives projective $\ZZ_\ell[G'_y]$-modules. The desired $\ZZ_\ell[G]$-modules $\Sw_y$, $\wt{\Sw}_y$, and $\Art_y$ are their inductions from $G'_y$ to $G$.

\begin{statement}{Remark 4.6}
The modules $\Sw_y$, $\wt{\Sw}_y$, and $\Art_y$ are purely local invariants. They are determined by the action of the stabilizer on the completed local ring $\wh{\cO}_{C',y'}$. Moreover $\Sw_y=0$ if and only if the cover is tamely ramified at $y'$; equivalently, the stabilizer $G_{y'}$ has order prime to $\operatorname{char}(k)$.
\end{statement}

\newpage

\section*{Expos\'e XI (end) -- Euler--Poincare Characteristic and Weil Formulas}
\addcontentsline{toc}{section}{Expos\'e XI (end) -- Euler--Poincare characteristic}

\subsection*{5. The Weil formula}

This final part of Expos\'e XI compares the trace formalism with the classical formula of Weil. The situation is that of a smooth proper curve $C$ over an algebraically closed field, an open dense subset $U\subset C$, a finite Galois covering $p:C'\to C$ with group $G$, and a coefficient ring $\Lambda$ such as $\ZZ_\ell$ or a finite quotient. The sheaves under consideration are constructible on $C$, locally constant on $U$, and controlled near the boundary by the local Galois data of the covering.

The Euler--Poincare characteristic is written
\[
   \chi(C,F)=\sum_i(-1)^i[H^i(C,F)]
\]
in the Grothendieck group of $\Lambda$-modules, or in the corresponding Grothendieck group of $\Lambda[G]$-modules when an equivariant structure is present. The goal is to express the difference between the global Euler characteristic and the rank contribution in terms of local ramification classes.

\begin{statement}{Proposition 5.1}
Let $F$ be a constructible sheaf on $C$, locally constant on $U$, and suppose that its pullback to $U'$ is represented by a finite projective $\Lambda[G]$-module. Then the equivariant Euler characteristic of $F$ is the sum of the unramified global term and local correction terms at the points of $C\setminus U$.
\end{statement}

The proof is a formal consequence of excision and induction. One chooses a finite Galois cover which trivializes the sheaf on $U$, writes the cohomology on $U$ in terms of the cohomology upstairs together with the $G$-module of fibers, and compares compactly supported cohomology with cohomology on the compact curve. The boundary contributions are precisely the Swan and Artin representations defined in Expos\'e VIII.

\subsection*{6. Definition of the local terms $\varepsilon_x^\Delta(F)$}

Let $x\in C\setminus U$. The local term $\varepsilon_x^\Delta(F)$ is an element of a Grothendieck group, attached to the action of the decomposition group at a point $x'$ above $x$. It measures the defect between the representation carried by the generic fiber of $F$ and the representation seen after extending across $x$. In concrete terms it is built from the Swan representation and the fiber $F_{\bar\eta}$:
\[
   \varepsilon_x^\Delta(F)
      = [\Sw_x\otimes F_{\bar\eta}]+\text{tame correction}-\text{invariant correction}.
\]
The precise expression depends on whether one writes the formula in $K(\Lambda)$, in $K(\Lambda[G])$, or after applying a character. The important point is that the local term is determined by the completed local ring and the local monodromy representation.

The notation $\Delta$ emphasizes that this term is the local contribution at the diagonal in the Lefschetz trace formula: it is the same obstruction which appears when one computes the intersection of the graph of an endomorphism with the diagonal.

\subsection*{7. Euler--Poincare formula}

\begin{statement}{Theorem 7.1}
Let $C$ be a smooth proper connected curve over an algebraically closed field, let $U\subset C$ be a dense open subset, and let $F$ be a constructible $\Lambda$-sheaf on $C$, locally constant on $U$. Then
\[
   \chi(C,F)=\operatorname{rk}(F)\chi(C,\Lambda)
      -\sum_{x\in C\setminus U} a_x(F),
\]
where $a_x(F)$ is the Artin conductor, expressed in $K$-theoretic form by the local term $\varepsilon_x^\Delta(F)$.
\end{statement}

\begin{prooftext}
The proof proceeds in four steps.

First, for a bounded-below complex $\Psi$ of sheaves of $\Lambda[G]$-modules on $C$, one uses the projection formula and the exactness of induction to identify the equivariant Euler characteristic of $p_*\Psi$ with the induced class from the cohomology of $\Psi$ upstairs.

Second, if $\Phi$ is a bounded-below complex on $C$ which is zero outside $U$, its compactly supported cohomology is computed by extension by zero and the preceding equivariant comparison.

Third, the local boundary triangle
\[
   j_!j^*F\to F\to i_*i^*F\to
\]
expresses $\chi(C,F)$ as a sum of the compactly supported contribution over $U$ and the stalk contributions at $C\setminus U$.

Finally, by choosing a cover $U'\to U$ which trivializes the sheaf and extending to a finite cover $C'\to C$, the local terms are identified with the Swan and Artin modules. This yields the stated conductor formula.
\end{prooftext}

\begin{statement}{Corollary 7.9}
If $F$ is locally constant on all of $C$, then all local conductor terms vanish and
\[
   \chi(C,F)=\operatorname{rk}(F)\chi(C,\Lambda).
\]
\end{statement}

\begin{statement}{Corollary 7.12}
For a constructible sheaf $F$ on a smooth curve, the Euler characteristic depends only on the generic rank, the Euler characteristic of the curve, and the local ramification invariants at the finitely many points where $F$ is not locally constant.
\end{statement}

\begin{statement}{Lemma 7.14}
The local conductor class is additive in distinguished triangles and in short exact sequences of constructible sheaves. It is invariant under replacing a local trivializing cover by a dominating cover.
\end{statement}

The proof is a diagram chase in the Grothendieck group, using exactness of restriction, induction, and the additivity of the Swan representation. This is the form which makes the Euler--Poincare formula independent of auxiliary choices.

\subsection*{References for Expos\'e XI}

The references are to the preceding theory of local representations, to the Grothendieck--Ogg--Shafarevich formula, and to the trace formula for sheaves on curves.

\newpage

\section*{Expos\'e XII -- Nielsen--Wecken and Lefschetz Formulas in Algebraic Geometry}
\addcontentsline{toc}{section}{Expos\'e XII -- Nielsen--Wecken and Lefschetz formulas}

\begin{center}
\textbf{Nielsen--Wecken and Lefschetz formulas in algebraic geometry}
\end{center}

\subsection*{1. General setting}

The expos\'e studies fixed point formulas for correspondences and endomorphisms in algebraic geometry, with particular attention to local terms. The input consists of a scheme $X$, an endomorphism or correspondence $f$, a constructible sheaf or complex $F$, and sometimes a finite group $G$ acting on a covering $X'\to X$. One wants a formula expressing the trace of the induced endomorphism on cohomology as a sum over fixed point classes.

In topology this is the Nielsen--Wecken refinement of the Lefschetz fixed point formula. In algebraic geometry, the role of a fixed point class is played by a connected component of the intersection of the graph of the correspondence with the diagonal, or by its lift to a finite Galois covering.

\subsection*{2. Fixed points and lifted fixed points}

Let $f:X\to X$ be an endomorphism and let $p:X'\to X$ be a finite covering with group $G$. A lift $f':X'\to X'$ need not commute with the $G$-action; it may instead satisfy
\[
   f' g = \varphi(g) f'
\]
for an endomorphism $\varphi:G\to G$. For $g\in G$ one considers the fixed points of
\[
   g^{-1}f':X'\to X'.
\]
The local terms attached to these fixed points are then assembled into a class function on $G$, or into an element of a Grothendieck group after applying traces.

This is the algebraic form of separating fixed points into Nielsen classes. The summation over $g\in G$ recovers the ordinary Lefschetz number downstairs.

\subsection*{3. The local Nielsen--Wecken invariant}

Let $x$ be a fixed point downstairs, and choose points $x'$ above it. The local invariant depends on the intersection multiplicity of the graph of $g^{-1}f'$ with the diagonal at $x'$, and on the trace of the induced endomorphism on the stalk of the sheaf. In the case of curves, if $R$ is the completed local ring, $\pi$ a uniformizer, and $h$ the induced endomorphism, the fixed point multiplicity is
\[
   \nu(h)=v(h(\pi)-\pi).
\]
This is the same local intersection calculation used earlier for the Swan representation.

\subsection*{4. Generalities on local terms}

For a prime $\ell\ne\operatorname{char}(k)$, the local term associated with $x$ is denoted
\[
   S_{x,\ell}(f',f,G,\varphi).
\]
It is a rational expression obtained by summing the local traces over the lifted fixed points.

\begin{statement}{Conjecture 4.5}
For every $\ell\ne\operatorname{char}(k)$, the values of $S_{x,\ell}(f',f,G,\varphi)$ are $\ell$-adic integers.
\end{statement}

The expos\'e proves this conjecture in dimension one by reducing the assertion to an explicit calculation in complete discrete valuation rings.

\begin{statement}{Proposition 4.6}
Let $R$ be a discrete valuation ring with maximal ideal $\mathfrak m$, uniformizer $\pi$, algebraically closed residue field $k$, and valuation $v$. Let $f'_R:R\to R$ be an endomorphism and let $G$ be a finite group of automorphisms of $R$. Assume that $f'_R$ and all elements of $G$ induce the identity on $k$, and that an endomorphism $\varphi:G\to G$ is given with the compatibility
\[
   f'_R\circ g = \varphi(g)\circ f'_R.
\]
Then the local expression defining the Nielsen--Wecken term is an $\ell$-adic integer.
\end{statement}

The proof reduces to the case $g=e$ and $\varphi=\mathrm{id}$, by factoring through the kernel on which the action is relevant. Let
\[
   r=\nu(f'_R)=v(f'_R(\pi)-\pi).
\]
If $r\ge2$, the following lemma controls the valuations of the iterates and of the terms obtained after applying elements of $G$.

\begin{statement}{Lemma 4.7}
Under the hypotheses of Proposition 4.6, assume in addition that $\varphi=\mathrm{id}_G$ and $r\ge2$. Then the valuation $r$ satisfies the stability conditions required to compute the local fixed point multiplicities, and the resulting local term is integral.
\end{statement}

The verification is a direct calculation in $R$. One writes $f'_R(\pi)=\pi+u\pi^r+\cdots$ with $u$ a unit and tracks the effect of automorphisms $g\in G$ on the leading term. The compatibility with $f'_R$ forces the same leading valuation for the relevant differences. Thus the denominator which appears in the formal local term is cancelled by the multiplicity.

\begin{statement}{Corollary 4.8}
The rational local numbers obtained from Proposition 4.6 have $\ell$-adic integral images for every $\ell\ne\operatorname{char}(k)$.
\end{statement}

\begin{statement}{Proposition 4.9}
Conjecture 4.5 is true when $\dim X=1$.
\end{statement}

The proof introduces the endomorphism on the completed local rings at the fixed points, reduces to the complete DVR calculation, and sums over the finite set of lifted fixed points.

\subsection*{5. Generalization of the Nielsen--Wecken formula}

\begin{statement}{Theorem 5.1}
In the preceding equivariant situation, the trace of the induced endomorphism on the cohomology of $X$ is equal to the sum of the local Nielsen--Wecken terms over the fixed point classes:
\[
   \operatorname{Tr}(f^*\mid R\Gamma(X,F))
      =\sum_x S_{x,\ell}(f',f,G,\varphi;F).
\]
\end{statement}

\begin{prooftext}
One first reduces modulo $\ell^n$ for every $n\ge1$, where the finite coefficient Lefschetz formula is available. The equivariant decomposition of the covering separates the trace into contributions indexed by $g\in G$. Passing to the inverse limit and using the integrality proved in dimension one gives the asserted $\ell$-adic formula. Compatibility of the local terms with reduction modulo $\ell^n$ ensures that the resulting expression is independent of $n$.
\end{prooftext}

\subsubsection*{5.10. Special cases}

If the covering is unramified near the relevant fixed points, the Swan terms vanish and the local formula reduces to the ordinary local Lefschetz multiplicity. If all multiplicities
\[
   \nu_{x'}(g^{-1}f')
\]
are equal to $1$ at the lifted fixed points, then each local contribution is simply the trace on the corresponding stalk, weighted by the number of lifted fixed points in the Nielsen class.

\subsection*{6. Application to a Lefschetz formula}

Let $X$ be proper and let $f:X\to X$ be an endomorphism with isolated fixed points, or more generally with a fixed point scheme satisfying the usual cleanliness hypotheses. Let $F$ be a constructible complex with an endomorphism compatible with $f$. The Lefschetz formula reads
\[
   \sum_i(-1)^i\operatorname{Tr}(f^*\mid H^i(X,F))
      =\sum_{x\in \operatorname{Fix}(f)}LT_x(f,F).
\]
The local term $LT_x(f,F)$ is computed by the Nielsen--Wecken construction after choosing a suitable covering and summing over the lifts.

\subsection*{6.4. Remarks on local terms}

If $g\in G$, conjugation by $g$ transports the fixed point data at $x'$ to that at $gx'$, and the local terms transform accordingly. Thus the local contribution depends only on the conjugacy class of the local stabilizer data.

Let $x'$ be a fixed point of a lift and let $G_{x'}$ be its stabilizer. This stabilizer is stable under $\varphi$, and the local term may be computed entirely in the complete local ring at $x'$ together with the induced endomorphism of $G_{x'}$. In the tamely ramified case, the Swan contribution vanishes; the local term is then the tame trace multiplied by the geometric fixed point multiplicity.

\begin{statement}{Lemma 6.4.5}
If the ramification at the fixed point is tame and the fixed point multiplicity is simple, the local term is the ordinary trace of the induced endomorphism on the stalk.
\end{statement}

When the sheaf $F$ is unramified at all fixed points, the formula simplifies to the standard Lefschetz trace formula with the local terms given by stalk traces and intersection multiplicities.

\subsection*{7. Comments on validity hypotheses}

The final comments emphasize that the formula depends on finiteness, properness, and constructibility assumptions. The fixed point locus must be proper over the base, or else cohomology with proper supports must be used. The sheaf must have finite Tor-amplitude in the relevant coefficient category, so that traces are defined. When these hypotheses hold, the local terms are stable under finite coefficient reduction and hence pass to the $\ell$-adic limit.

\subsection*{7. Complements: end of Expos\'e XII}

The end of the expos\'e adds compatibility statements for the local terms. They are invariant under replacing the Galois cover by a dominating cover, additive in distinguished triangles of coefficients, and compatible with the specialization maps used in the proof of the curve case. These facts allow the formula to be applied without making the auxiliary covering part of the final data.

The bibliography of Expos\'e XII includes the topological Nielsen--Wecken references, the Grothendieck--Verdier trace formalism, and the preceding SGA expos\'es on local terms and Euler--Poincare formulas.

\newpage

\section*{Expos\'e XV -- Frobenius Morphism and Rationality of $L$-Functions}
\addcontentsline{toc}{section}{Expos\'e XV -- Frobenius and rationality of $L$-functions}

\begin{center}
\textbf{Frobenius morphism and rationality of $L$-functions}
\end{center}

\subsection*{1. Frobenius morphism}

\subsubsection*{1.1. Frobenius endomorphism}

Let $X$ be a scheme of characteristic $p>0$. The absolute Frobenius endomorphism
\[
   F_X:X\to X
\]
is the identity on the underlying topological space and raises functions to their $p$th powers:
\[
   a\longmapsto a^p.
\]
If $X$ is defined over $\FF_q$, one also considers the $q$th Frobenius, obtained by iterating $F_X$ so that functions are sent to $a^q$.

For a geometric point $\bar x$ over a closed point $x\in |X|$, the Frobenius element acts on the geometric stalk of a constructible sheaf. The convention in the trace formula is to use geometric Frobenius, inverse to the arithmetic Frobenius on residue fields.

\subsubsection*{1.2. Relative Frobenius}

Let $X\to S$ be a morphism of schemes in characteristic $p$. The absolute Frobenius maps of $X$ and $S$ form a commutative square only after twisting the structure morphism. The relative Frobenius is the morphism
\[
   F_{X/S}:X\to X^{(p)}=X\times_{S,F_S}S.
\]
When $S=\Spec k$ for a perfect field $k$, the Frobenius twist $X^{(p)}$ is another $k$-scheme, often identified with $X$ after choosing the Frobenius automorphism of $k$.

The relative Frobenius is functorial in $X/S$, compatible with products, and finite under the usual finite type hypotheses. It is the geometric morphism which relates point-counting over finite fields to traces on cohomology.

\subsubsection*{1.3. Effect on a module}

If $M$ is a module over a ring $A$ of characteristic $p$, the Frobenius pullback is
\[
   F_A^*M=A\otimes_{F_A,A}M.
\]
For sheaves this construction is local on $X$. A semilinear Frobenius structure on a sheaf is an isomorphism between the pullback by Frobenius and the sheaf itself, or dually an action of Frobenius on the stalks. Such structures are the coefficient data in the cohomological interpretation of $L$-functions.

\subsubsection*{1.4. Frobenius and cohomology}

Let $X$ be separated of finite type over $\FF_q$ and let $F$ be a constructible $\ell$-adic sheaf. The geometric Frobenius acts on
\[
   H^i_c(X_{\overline{\FF}_q},F).
\]
The Lefschetz trace formula gives
\[
   \sum_i(-1)^i\operatorname{Tr}\bigl(F^n\mid H^i_c(X_{\overline{\FF}_q},F)\bigr)
     =\sum_{x\in X(\FF_{q^n})}\operatorname{Tr}\bigl(F_x^n\mid F_{\bar x}\bigr).
\]
This identity is the bridge between Euler products and determinants of Frobenius on cohomology.

\subsection*{2. Rationality of $L$-functions}

\subsubsection*{2.1. The theorem}

Let $X$ be separated of finite type over $\FF_q$ and let $F$ be a constructible $\QQ_\ell$-sheaf on $X$. Define
\[
   L(X,F;t)=\prod_{x\in |X|}
      \det(1-F_x t^{\deg x}\mid F_{\bar x})^{-1}.
\]

\begin{statement}{Theorem}
The function $L(X,F;t)$ is rational. More precisely,
\[
   L(X,F;t)=\prod_i
      \det(1-Ft\mid H^i_c(X_{\overline{\FF}_q},F))^{(-1)^{i+1}}.
\]
\end{statement}

This is Grothendieck's cohomological rationality theorem for $L$-functions. The proof is based on the trace formula and on a reduction which permits one to apply the formula to enough powers of Frobenius.

\subsubsection*{2.2. A reduction lemma}

The reduction lemma says that it suffices to prove the theorem after replacing $X$ by a stratification on which $F$ is locally constant and after replacing $F$ by a sheaf trivialized by a finite cover. Additivity of compactly supported cohomology and multiplicativity of Euler products in distinguished triangles reduce the problem to these cases.

One also reduces to curves in the key step. Given a closed point contribution, one chooses curves passing through the point and uses the Lefschetz formula on those curves. The local computations of Expos\'es XI and XII control the terms which arise in this passage.

\subsubsection*{2.3. Proof of the theorem}

The logarithm of the Euler product is
\[
   \log L(X,F;t)=\sum_{n\ge1}\frac{t^n}{n}
       \sum_{x\in X(\FF_{q^n})}\operatorname{Tr}(F_x^n\mid F_{\bar x}).
\]
By the Lefschetz trace formula the inner sum is
\[
   \sum_i(-1)^i\operatorname{Tr}(F^n\mid H^i_c(X_{\overline{\FF}_q},F)).
\]
For a finite-dimensional vector space $V$ with endomorphism $T$ one has
\[
   \exp\left(\sum_{n\ge1}\operatorname{Tr}(T^n)\frac{t^n}{n}\right)
      =\det(1-Tt)^{-1}.
\]
Applying this identity to each cohomology group gives the determinant formula, hence rationality.

\subsubsection*{First reduction}

The first reduction is from arbitrary constructible sheaves to sheaves which are locally constant on a dense open subset and zero or controlled on its complement. This is done by noetherian induction on the support and by the localization triangle
\[
   j_!j^*F\to F\to i_*i^*F\to.
\]

\subsubsection*{Second reduction}

The second reduction passes to a finite Galois cover which trivializes the locally constant sheaf. The equivariant form of the trace formula expresses the downstairs trace as an average of traces upstairs, weighted by the relevant group elements. This is where the Nielsen--Wecken and local term formulas enter.

\subsubsection*{Third reduction}

The third reduction handles ramification. By compactifying and using the curve case, local ramification terms are shown to contribute rational factors. The integrality and invariance of the local terms ensure that these factors are independent of auxiliary choices.

\subsubsection*{End of the proof}

After the reductions, the trace formula applies directly. The determinant expression obtained from compactly supported cohomology is a rational function, and the coefficient comparison of logarithmic derivatives identifies it with the Euler product. Thus the $L$-function is rational.

\subsection*{Bibliography of Expos\'e XV}

The bibliography includes Grothendieck's Bourbaki expos\'e on the Lefschetz formula and rationality of $L$-functions, the preceding SGA trace-formula expos\'es, and the finiteness and constructibility theorems for \etale cohomology.

\subsection*{Terminological index}

The terminological index includes the following principal entries in modern English usage: Artin conductor, constructible $\ell$-adic sheaf, Euler--Poincare characteristic, Frobenius morphism, geometric Frobenius, Grothendieck group, Lefschetz local term, Nielsen--Wecken class, perfect complex, pseudo-coherent complex, relative Frobenius, Swan representation, and trace formula.

\subsection*{Index of notation}

The main notation in this range includes $F_X$ for absolute Frobenius, $F_{X/S}$ for relative Frobenius, $H^i_c$ for compactly supported cohomology, $K(\cC)$ for a Grothendieck group, $\Sw_y$ and $\Art_y$ for local Swan and Artin representations, $Z(X,F;t)$ and $L(X,F;t)$ for zeta and $L$-functions, and $\D^{\mathrm{parf}}$ for perfect complexes.


\end{document}
